% Lecture notes on the classification of homogeneous plane continua,
% following the formal (Lean 4) proof in this repository.
% Volume 5: Plane topology II.
%
% Compile with:  pdflatex main && pdflatex main   (two runs for the contents and references)
\documentclass[11pt,oneside]{book}

\usepackage[T1]{fontenc}
\usepackage{lmodern}
\usepackage[margin=1.15in]{geometry}
\usepackage{amsmath,amssymb,amsthm}
\usepackage{enumitem}
\usepackage{booktabs}
\usepackage{longtable}
\usepackage{tikz}
\usepackage[hidelinks]{hyperref}

% ---------------------------------------------------------------- theorem styles
\theoremstyle{plain}
\newtheorem{theorem}{Theorem}[chapter]
\newtheorem{lemma}[theorem]{Lemma}
\newtheorem{proposition}[theorem]{Proposition}
\newtheorem{corollary}[theorem]{Corollary}
\newtheorem{claim}[theorem]{Claim}
\theoremstyle{definition}
\newtheorem{definition}[theorem]{Definition}
\newtheorem{example}[theorem]{Example}
\newtheorem{exercise}[theorem]{Exercise}
\theoremstyle{remark}
\newtheorem{remark}[theorem]{Remark}

% ---------------------------------------------------------------- notation
\newcommand{\N}{\mathbb{N}}
\newcommand{\Z}{\mathbb{Z}}
\newcommand{\R}{\mathbb{R}}
\newcommand{\diam}{\operatorname{diam}}
\newcommand{\dist}{\operatorname{dist}}
\newcommand{\interior}{\operatorname{int}}
\newcommand{\cl}[1]{\overline{#1}}
\newcommand{\abs}[1]{\left\lvert #1\right\rvert}
\newcommand{\set}[1]{\left\{ #1 \right\}}
\newcommand{\fl}[1]{\left\lfloor #1\right\rfloor}
\newcommand{\C}{\mathbb{C}}
% the principal logarithm and argument
\newcommand{\Log}{\operatorname{Log}}
\newcommand{\Arg}{\operatorname{Arg}}
\renewcommand{\Re}{\operatorname{Re}}
\renewcommand{\Im}{\operatorname{Im}}
% the sup norm of the plane and the euclidean modulus
\newcommand{\nrm}[1]{\lVert #1\rVert}
% the interior of a compact set of the plane (union of the bounded complementary components)
\newcommand{\inn}{\operatorname{ins}}
% crossing parity of a polygon or graph at a point
\newcommand{\parity}{\operatorname{par}}
% the component of x in the complement of A
\newcommand{\comp}[2]{\operatorname{C}_{#1}(#2)}
% the point set of a graph
\newcommand{\gs}[1]{\lvert #1\rvert}
% the degree of a vertex
\newcommand{\dg}{\operatorname{deg}}
% the odd set of a graph
\newcommand{\odd}{\operatorname{odd}}

% Names of declarations in the Lean development (typewriter, underscores allowed).
\newcommand{\lean}[1]{\texttt{\detokenize{#1}}}
% A Lean name ending in a subscript digit (such as IsEssentialSum with subscript 2).
\newcommand{\leansub}[2]{\texttt{\detokenize{#1}}\ensuremath{_{#2}}}
% A paragraph pointing to the formal counterpart of a result.
\newcommand{\formal}[1]{\par\smallskip\noindent{\small\textit{Formal counterpart:} #1}\par\smallskip}
% References to numbered results of Volumes 1 and 2.
\newcommand{\vone}[1]{Volume~1, #1}
\newcommand{\vtwo}[1]{Volume~2, #1}
\newcommand{\vthree}[1]{Volume~3, #1}
\newcommand{\vfour}[1]{Volume~4, #1}
% the point set of a set of elements, the element of a point
\newcommand{\pts}[1]{#1^{\ast}}
\newcommand{\el}[1]{[#1]}
% closure in Moore's sense, boundary, exterior
\newcommand{\bd}{\operatorname{bd}}
\newcommand{\ext}{\operatorname{ext}}
\newcommand{\fr}{\operatorname{fr}}
% open arc
\newcommand{\oarc}[1]{\mathring{#1}}

\setlist{itemsep=2pt,topsep=4pt}

\title{\textbf{Homogeneous Plane Continua}\\[6pt]
\Large Lecture notes following the formal proof\\[18pt]
\large Volume 5\\[4pt]
\Large Plane topology II: decompositions of the plane,\\ the Jordan curve theorem and\\ Moore's characterization of the plane}
\author{}
\date{}

\begin{document}

\frontmatter
\maketitle
\tableofcontents

\mainmatter
% The introduction is Chapter 0, as in Volumes 1--4.
\setcounter{chapter}{-1}
%======================================================================
\chapter{Introduction}\label{ch:intro}
%======================================================================

\section{What this volume is about}

Volume~4 developed the plane topology that can be done with logarithms and polygons. This volume
adds the three remaining pieces used by the classification of homogeneous plane continua: the
Jordan curve theorem for arbitrary simple closed curves, R.~L.~Moore's axiomatic characterization of
the plane, and Moore's theorem that the space of pieces of a suitable decomposition of the plane into
continua is again a plane. The last is the form in which the volume is used later: it turns a
monotone image of a plane continuum into a plane set.

The volume follows two papers of Moore, and the order of the formal proof:
\emph{Concerning upper semi-continuous collections of continua} (1925), and \emph{Concerning a set of
postulates for plane analysis situs} (1919), which rests on his \emph{On the foundations of plane
analysis situs} (1916, cited as F.A.). The 1925 paper proves its main theorem by verifying the
axioms of the 1919 paper; on the way, its Theorem~11 is a Jordan curve theorem for curves of
elements, which for the decomposition into points is the classical one.

The volume has nine chapters.

\begin{itemize}
\item \textbf{Chapter~\ref{ch:arcs}: Arcs and simple closed curves.} Irreducible continua between two
points with no other non-cut points are arcs; the cut-point characterizations of arcs and circles;
limits of arcs; replacing a piece of an arc.
\item \textbf{Chapter~\ref{ch:usc}: Decompositions of the plane into non-separating continua.} Moore
decompositions, the quotient space of elements, its metrizability, the correspondence of components,
and Moore's Theorems~1--4.
\item \textbf{Chapter~\ref{ch:elements}: Arcs of elements.} Moore's Theorems~5--10 and~12, and the
decomposition into points.
\item \textbf{Chapter~\ref{ch:jordan}: Mullikin's theorem and the Jordan curve theorem.} Mullikin's
theorem on two non-separating continua with disconnected intersection, Moore's Theorem~11, and the
Jordan curve theorem.
\item \textbf{Chapter~\ref{ch:axioms}: Moore's axioms and the bridge to topology.} The eight axioms,
the topology generated by the regions, metrizability, arcs in domains (F.A.\ Theorem~15).
\item \textbf{Chapter~\ref{ch:curves}: Closed curves in a Moore space.} F.A.\ Theorems~20--43:
interiors and exteriors, cross-cuts, closed curves in regions, chords.
\item \textbf{Chapter~\ref{ch:rulings}: Double rulings and Moore's Theorem~D.} Quadrilaterals, double
rulings, a corrected Theorem~C, and fine rulings of a quadrilateral.
\item \textbf{Chapter~\ref{ch:plane}: Grids, the Schoenflies theorem and Moore's characterization of the
plane.} A compact set with a grid is a square; the Schoenflies theorem and the annulus lemma in a
Moore space; exhaustion by closed curves; every Moore space is a plane.
\item \textbf{Chapter~\ref{ch:quotient}: The space of elements is a plane.} Moore's Theorems~13--22,
monotone decompositions of a plane continuum, and the cross-cut, Schoenflies and collar theorems in
the plane.
\end{itemize}

The main results are:

\begin{theorem}[Jordan curve theorem; Theorem~\ref{thm:jct}]
If $J\subseteq\R^2$ is homeomorphic to a circle, $\R^2\setminus J$ has exactly two components, one bounded
and one unbounded, and $J$ is the frontier of each.
\end{theorem}

\begin{theorem}[Moore's characterization of the plane; Theorem~\ref{thm:moore-plane}]
A set with a class of regions satisfying Moore's Axioms~1--8 is homeomorphic to $\R^2$, with the
topology generated by the regions.
\end{theorem}

\begin{theorem}[Schoenflies; Theorems~\ref{thm:schoenflies} and~\ref{thm:plane-schoenflies}]
The closure of the bounded complementary domain of a simple closed curve in the plane is a closed disc
(a square), by a homeomorphism that can be prescribed on one arc of the curve.
\end{theorem}

\begin{theorem}[Moore 1925; Theorem~\ref{thm:m25-22}]
If the plane is decomposed into compact connected sets, none of which separates the plane, in an upper
semicontinuous way, then the space of these sets with the quotient topology is homeomorphic to $\R^2$.
\end{theorem}

\begin{corollary}[Corollary~\ref{cor:monotone-plane}]
A monotone image of a plane continuum whose fibres do not separate the plane is homeomorphic to a compact
subset of the plane.
\end{corollary}

\section{Where these results are used}

\begin{itemize}
\item Moore's Theorem~22 and the monotone-image corollary are used in Hagopian's arguments on
indecomposable homogeneous plane continua (Volume~6), where a plane continuum is decomposed into layers or
into the fibres of a monotone map, and in Jones's work on decomposable homogeneous continua (Volume~7):
in each case the space of pieces is recognised as a plane, or as a plane continuum.
\item The Jordan curve theorem and the circle characterization (Chapter~\ref{ch:arcs}) are used throughout
Hagopian's constructions and in Jones's study of circle-like continua (Volumes~6 and~7).
\item The cross-cut, Schoenflies and collar theorems in the plane (Chapter~\ref{ch:quotient}) are used in
Hagopian's constructions (Volume~6) and in Jones's arguments (Volume~7).
\end{itemize}

\section{Where the formal proof differs from the literature}\label{sec:intro-diff}

\begin{enumerate}
\item \textbf{Topology instead of Moore's primitive notions.} Moore works with limit elements, domains and
regions defined by his axioms. The formal proof translates them into the quotient topology
(Proposition~\ref{prop:moore-notions}) and into the topology generated by the regions
(Proposition~\ref{prop:bridge}), and then uses standard topology: metrizability, compactness, the
cut-point characterizations of arcs and circles.
\item \textbf{Printed statements that are false.} Several statements are false as printed and are proved here in
a corrected form; in each case the formal development proves the negation of the printed form.
\begin{itemize}
\item Moore 1925, Theorem~2, with ``complementary domain'' in Moore's sense (Remark~\ref{rem:m25-2-printed}).
\item Moore 1925, Theorems~10 and~11, with Moore's printed definition of a simple closed curve of
elements, and Theorem~7, with ``$p$'' for ``$x$'' (Chapter~\ref{ch:elements}, after
Definition~\ref{def:arc-elt} and after Theorem~\ref{thm:m25-5-10}).
\item Mullikin's theorem as quoted by Moore (Remark~\ref{rem:mullikin-printed}).
\item F.A.\ Theorems~4, 14, 19 and~35, which fail for the empty region, and Theorem~43 for an empty family
(Remark~\ref{rem:empty-region}, Theorems~\ref{thm:fa35-36} and~\ref{thm:fa43}).
\item Moore 1919, Theorem~C (Theorem~\ref{thm:theoremC}): a curve may lie within the quadrilateral without
crossing it suitably. The formal proof of Theorem~D uses a corrected Theorem~C for curves that cross the
quadrilateral transversally.
\end{itemize}
\item \textbf{The Jordan curve theorem without the characterization of the plane.} The Jordan curve theorem is
obtained from Moore's Theorem~11 for the decomposition into points (Chapter~\ref{ch:jordan}), using
Janiszewski's and Borsuk's theorems from Volume~4 but none of Chapters~\ref{ch:axioms}--\ref{ch:plane}.
\item \textbf{Coordinates by a grid.} Moore 1919 reaches the plane through ``nets'' of curves (Theorems~E--G).
The formal proof proves a general recognition theorem for squares (a compact set with two families of grid
lines, Theorem~\ref{thm:grid-square}), derives the Schoenflies theorem in a Moore space from Theorem~D, and
pastes squares and collars along an exhaustion (Chapter~\ref{ch:plane}); Theorems~E and~G become corollaries.
\item \textbf{Moore's Axioms~6 and~7 for the space of elements.} Moore obtains them through F.A.\ and the 1919
paper; here they are proved directly from Theorems~12, 15, 18 and~19 (Theorem~\ref{thm:m25-21}).
\item \textbf{Proofs in outline.} Chapters~\ref{ch:curves}--\ref{ch:plane} contain some long case analyses
(several of F.A.'s theorems on closed curves, the corrected Theorem~C, Theorem~D). For these the notes give
the proofs in outline, marked as such, and indicate which earlier results they use; the complete arguments are
in the formal development, at the declarations named in the ``Formal counterpart'' lines.
\end{enumerate}

\section{Prerequisites and conventions}\label{sec:intro-conv}

The reader should know the basic metric topology of \vone{Chapter~1} and the results of Volume~4 on
logarithms, Brouwer's theorem, Janiszewski's theorem and Borsuk's theorem. From Volume~3 we use
irreducible continua, non-cut points and upper semicontinuous decompositions.

\begin{itemize}
\item \textbf{The plane.} As in Volume~4, $\R^2$ carries the sup metric, so that balls are squares.
\item \textbf{Elements.} For a decomposition of the plane, $\el x$ is the element containing $x$, $G$ is the space of
elements, and $\pts H$ is the point set of a set $H$ of elements (Definition~\ref{def:moore-decomp}).
\item \textbf{Arcs and curves.} An \emph{arc} from $a$ to $b$ is the image of an embedding of $[0,1]$ with $0\mapsto a$,
$1\mapsto b$; a \emph{simple closed curve} is the union of two arcs with only their ends in common, equivalently
a set homeomorphic to a circle (Chapter~\ref{ch:arcs}). $\oarc{\gamma}$ is the arc minus its ends.
\item \textbf{Moore spaces.} In Chapters~\ref{ch:axioms}--\ref{ch:plane}, $S$ is a set with regions satisfying some of
Axioms~1--8; $I(J)$ is the interior of a closed curve $J$, and ``within'', ``without'' mean in the interior,
in the exterior.
\item \textbf{Frontier.} $\fr A$ is the topological frontier; $\bd M$ is Moore's boundary.
\end{itemize}

\paragraph{Formal counterparts.} As in the earlier volumes, many results are followed by a line
``Formal counterpart:'' naming the corresponding declaration of the Lean development (all names
begin with \lean{HPC.}). These lines may be skipped. Appendix~\ref{app:formal} collects them.

\paragraph{Exercises.} Each chapter ends with exercises. None of them is used later.

%======================================================================
\chapter{Arcs, cut points and simple closed curves}\label{ch:arcs}
%======================================================================

Everything in this chapter is general topology: no plane and no axioms of Moore's enter. We
develop the facts about arcs and simple closed curves that the rest of the volume uses constantly:
sub-arcs, first points, gluing, the characterization of arcs as irreducible continua and as
continua cut by each of their interior points, the characterization of simple closed curves as
continua cut by every pair of points, limits of sequences of arcs, and the replacement of a null
family of sub-arcs of an arc.

Throughout, $X$ is a Hausdorff space and $I=[0,1]$.

\section{Arcs and simple closed curves}

\begin{definition}\label{def:arc}
Let $A\ne B$ be points of $X$. A set $C\subseteq X$ is an \emph{arc from $A$ to $B$} if there is an
embedding $f\colon I\to X$ (a homeomorphism onto its image) with $f(I)=C$, $f(0)=A$ and $f(1)=B$.
$A$ and $B$ are the \emph{ends} of $C$; $\oarc C=C\setminus\set{A,B}$ is the \emph{open arc}. A
\emph{simple closed curve} is a set $J=C_1\cup C_2$ where $C_1,C_2$ are arcs from $A$ to $B$ with
$C_1\cap C_2=\set{A,B}$.
\end{definition}

\formal{\lean{HPC.Moore.Gen.TArc}, \lean{HPC.Moore.Gen.TSCC}.}

Since $I$ is compact and $X$ Hausdorff, a continuous injective $f\colon I\to X$ is automatically an
embedding (a continuous bijection from a compact space onto a Hausdorff space is a homeomorphism);
so to produce an arc it suffices to produce a continuous injection (\lean{HPC.Moore.Gen.TArc.of_inj}).
An arc from $A$ to $B$ is an arc from $B$ to $A$ (use $t\mapsto f(1-t)$), and it is compact,
closed and connected.

\begin{lemma}[Sub-arcs]\label{lem:subarc}
Let $f\colon I\to X$ be an embedding.
\begin{enumerate}[label=(\alph*)]
\item For $s<t$ in $I$, $f([s,t])$ is an arc from $f(s)$ to $f(t)$.
\item If $C$ is an arc from $A$ to $B$ and $P\ne Q$ are points of $C$, there is an arc from $P$ to
$Q$ contained in $C$.
\item The open arc $\oarc C$ of an arc $C$ is nonempty, connected and dense in $C$; so are
$C\setminus\set A$ and $C\setminus\set B$.
\end{enumerate}
\end{lemma}

\begin{proof}
(a) Compose $f$ with the affine homeomorphism $u\mapsto s+u(t-s)$ of $I$ onto $[s,t]$. (b) Write
$P=f(s)$, $Q=f(t)$ and apply (a) to $[\min(s,t),\max(s,t)]$, reversing if $s>t$. (c)
$\oarc C=f((0,1))$ is the image of a nonempty connected set, dense in $I$; similarly for
$f((0,1])$ and $f([0,1))$.
\end{proof}

\formal{\lean{HPC.Moore.Gen.TArc.image_Icc}, \lean{HPC.Moore.Gen.TArc.exists_sub},
\lean{HPC.Moore.Gen.TArc.nonempty_diff}, \lean{HPC.Moore.Gen.TArc.isPreconnected_diff},
\lean{HPC.Moore.Gen.TArc.subset_closure_diff}.}

The next lemma is used dozens of times in Chapters~\ref{ch:curves}--\ref{ch:quotient}.

\begin{lemma}\label{lem:conn-frontier}
Let $C,U\subseteq X$ with $C$ connected and $C\cap\fr U=\emptyset$. If $C$ meets $U$ then
$C\subseteq U$; if $C$ meets $X\setminus\cl U$ then $C\subseteq X\setminus\cl U$.
\end{lemma}

\begin{proof}
$C$ lies in $X\setminus\fr U=\interior U\cup(X\setminus\cl U)$, a union of two disjoint open sets,
so it lies in one of them.
\end{proof}

\formal{\lean{HPC.Moore.Gen.conn_sub_of_frontier}, \lean{HPC.Moore.Gen.conn_sub_ext_of_frontier}.}

\begin{lemma}[First points, gluing, joining]\label{lem:arc-ops}
\leavevmode
\begin{enumerate}[label=(\alph*)]
\item \emph{(First point.)} Let $C$ be an arc from $A$ to $B$ and $Z$ a closed set meeting $C$ with
$A\notin Z$. There are $P\in Z\cap C$ and an arc $AP\subseteq C$ with $AP\cap Z=\set P$.
\item \emph{(Gluing.)} If $D$ is an arc from $P$ to $Q$, $E$ an arc from $Q$ to $R$ and
$D\cap E=\set Q$, then $D\cup E$ is an arc from $P$ to $R$.
\item \emph{(Joining.)} If $\alpha$ is an arc from $P$ to $Q$, $\beta$ an arc from $R$ to $T$,
$\alpha\cap\beta\ne\emptyset$ and $P\ne R$, then $\alpha\cup\beta$ contains an arc from $P$ to
$R$.
\end{enumerate}
\end{lemma}

\begin{proof}
(a) Let $f$ parametrize $C$. The set $f^{-1}(Z)$ is closed, nonempty and misses $0$; let $t$ be its
least element and $P=f(t)$. Then $AP=f([0,t])$ works. (b) With $f$ for $D$ and $g$ for $E$, the map
$h(u)=f(2u)$ for $u\le\frac12$, $h(u)=g(2u-1)$ for $u\ge\frac12$ is continuous and injective (the
two halves only share $Q=h(\frac12)$). (c) If $P\in\beta$, take the sub-arc of $\beta$ from $P$ to
$R$. Otherwise apply (a) to $\alpha$ and the closed set $\beta$: we get $X_0\in\beta$ and an arc
$PX_0\subseteq\alpha$ meeting $\beta$ only in $X_0$. If $X_0=R$ we are done; otherwise glue $PX_0$
to the sub-arc $X_0R$ of $\beta$ by (b).
\end{proof}

\formal{\lean{HPC.Moore.Gen.TArc.first_point}, \lean{HPC.Moore.Gen.TArc.union},
\lean{HPC.Moore.Gen.TArc.join}.}

\section{Irreducible continua are arcs}

A set $C$ containing $A\ne B$ is \emph{irreducible between $A$ and $B$} if $C$ is connected and no
connected proper subset of $C$ contains both $A$ and $B$ (\vthree{Definition~2.1}). Every arc is
irreducible between its ends: a connected subset of $f(I)$ containing $f(0)$ and $f(1)$ is $f$ of a
connected subset of $I$ containing $0$ and $1$, which is all of $I$. The converse, in metrizable
spaces, is the main theorem of this section. In Moore's 1916 paper an \emph{arc} is \emph{defined}
as a bounded closed irreducible connected set (following Lennes), and the theorem identifies
Moore's arcs with the topological ones.

Fix a compact connected separable Hausdorff space $Y$ irreducible between $a\ne b$, i.e.\ every
connected $H\subseteq Y$ containing $a$ and $b$ is $Y$.

\begin{lemma}\label{lem:insert-conn}
Let $Z$ be a connected space, $p\in Z$, and $Z\setminus\set p=U\cup V$ with $U,V$ disjoint and open.
Then $U\cup\set p$ is connected.
\end{lemma}

\begin{proof}
$U\cup\set p=Z\setminus V$ is closed. If $U\cup\set p=F_1\cup F_2$ with $F_1,F_2$ disjoint, nonempty
and closed, say $p\in F_1$, then $F_2\subseteq U$ is closed in $Z$, and $F_2=U\setminus F_1$ is open
in $Z$. So $F_2$ is a nonempty proper clopen subset of $Z$, which is impossible.
\end{proof}

\formal{\lean{HPC.Moore.ac_insert_isPreconnected}.}

\begin{lemma}[Cuts]\label{lem:cut}
\leavevmode
\begin{enumerate}[label=(\alph*)]
\item $Y\setminus\set a$ and $Y\setminus\set b$ are connected.
\item For every $p\ne a,b$ there are disjoint connected open sets $U\ni a$, $V\ni b$ with
$Y\setminus\set p=U\cup V$.
\end{enumerate}
\end{lemma}

\begin{proof}
The basic observation is: \emph{if $Y\setminus\set p=U_2\cup W$ with $U_2,W$ disjoint and open,
$U_2\ne\emptyset$, then $W\cup\set p$ does not contain both $a$ and $b$}; for $W\cup\set p$ is
connected (Lemma~\ref{lem:insert-conn}) and is a proper subset of $Y$.

(a) If $Y\setminus\set b$ were not connected, write $Y\setminus\set b=U_2\cup W$ (disjoint, open,
nonempty), with $a\in W$. Then $W\cup\set b$ contains $a$ and $b$, contradicting the observation.

(b) $Y\setminus\set p$ is a proper subset of $Y$ containing $a$ and $b$, so it is not connected:
$Y\setminus\set p=U\cup V$ with $U,V$ disjoint, open and nonempty. By the observation (applied with
$(U_2,W)=(U,V)$ and with $(U_2,W)=(V,U)$) neither $U\cup\set p$ nor $V\cup\set p$ contains both $a$
and $b$; so $a$ and $b$ lie in different pieces, say $a\in U$, $b\in V$. If $U$ were not connected,
$U=U_1\cup U_2$ with $U_1,U_2$ disjoint, open and nonempty, $a\in U_1$; then $Y\setminus\set
p=U_2\cup(U_1\cup V)$ and $U_1\cup V\cup\set p$ contains $a$ and $b$, against the observation. So $U$
is connected, and likewise $V$.
\end{proof}

\formal{\lean{HPC.Moore.ac_compl_isPreconnected}, \lean{HPC.Moore.ac_cut}.}

For $x,y\in Y$ let $L(x,y)$ be the component of $Y\setminus\set y$ containing $x$ (empty if $x=y$).
Define $x\prec y$ if $x\in L(a,y)$, i.e.\ $x\ne y$ and $x$ lies on the $a$-side of $y$.

\begin{proposition}\label{prop:arc-order}
$\prec$ is a strict linear order on $Y$ with least element $a$ and greatest element $b$. The rays
$\set{x:x\prec y}=L(a,y)$ and $\set{x:y\prec x}=L(b,y)$ are open, and the order is dense: if $x\prec
y$ there is $z$ with $x\prec z\prec y$.
\end{proposition}

\begin{proof}
We use two facts. (i) By Lemma~\ref{lem:cut}, for $y\ne a,b$ the set $Y\setminus\set y$ is the
union of the disjoint connected open sets $L(a,y)$ and $L(b,y)$ (the sets $U,V$ of the lemma are
connected, so they are these components); $Y\setminus\set b=L(a,b)$ and $Y\setminus\set a=L(b,a)$.
In particular $b\notin L(a,y)$ for every $y$, and $L(a,a)=\emptyset$. (ii) A connected set
containing $x$ and missing $z$ lies in $L(x,z)$. By (i) and Lemma~\ref{lem:insert-conn}, the sets
$A_y=L(a,y)\cup\set y$ ($y\ne a$) and $B_y=L(b,y)\cup\set y$ ($y\ne b$) are connected.

\emph{Asymmetry.} Suppose $x\prec y$ and $y\prec x$. By (i), $x,y\notin\set{a,b}$. The connected set
$B_y$ contains $b$ and $y$ and misses $x$ (as $x\in L(a,y)$), so by (ii) $B_y\subseteq L(y,x)$,
which is $L(a,x)$ because $y\in L(a,x)$. Then $b\in L(a,x)$, against (i).

\emph{Transitivity.} Let $x\prec y\prec z$. Then $y\ne a$ (as $L(a,a)=\emptyset$), so $A_y$ is
connected; it contains $a$ and $x$, and misses $z$ ($z\ne y$, and $z\in L(a,y)$ would mean $z\prec
y$, against asymmetry). By (ii), $A_y\subseteq L(a,z)$, so $x\prec z$.

\emph{Totality.} Let $y\ne z$ with $y\not\prec z$. If $z=a$, then $a\in L(a,y)$, i.e.\ $z\prec y$.
The case $z=b$ is impossible, since $L(a,b)=Y\setminus\set b\ni y$. Otherwise $y\in L(b,z)$ by (i),
so $A_z$ is a connected set containing $a$ and $z$ and missing $y$, and $A_z\subseteq L(a,y)$ by
(ii): $z\prec y$.

\emph{Rays.} $x\in L(b,y)$ if and only if $y\prec x$: if $x\in L(b,y)$ then $x\ne y$ and $x\notin
L(a,y)$, so $y\prec x$ by totality; if $y\prec x$ then $y\ne b$, $x\ne y$ and $x\not\prec y$, so
$x\in L(b,y)$ by (i). Hence both rays are open. For $x\ne a$, $a\in L(a,x)$, i.e.\ $a\prec x$; for
$x\ne b$, $b\in L(b,x)$, i.e.\ $x\prec b$.

\emph{Density.} If $x\prec y$ and no $z$ lies between, then $Y$ is the union of
$\set{w:w\preceq x}=Y\setminus L(b,x)$ and $\set{w:y\preceq w}=Y\setminus L(a,y)$, two disjoint
nonempty closed sets, against connectedness.
\end{proof}

\formal{\lean{HPC.Moore.acL}, \lean{HPC.Moore.acL_asymm}, \lean{HPC.Moore.acL_trans},
\lean{HPC.Moore.acL_total}, \lean{HPC.Moore.acL_ray}, \lean{HPC.Moore.ac_exists_between}.}

\begin{theorem}[Irreducible continua are arcs]\label{thm:irr-arc}
A compact connected separable Hausdorff space $Y$ irreducible between $a\ne b$ is homeomorphic to
$I$ by a homeomorphism sending $0\mapsto a$, $1\mapsto b$.
\end{theorem}

\begin{proof}
Let $D\subseteq Y$ be countable and dense, and $D'=D\setminus\set{a,b}$. With the order of
Proposition~\ref{prop:arc-order}, $D'$ is a countable dense linear order without end points (between
two points of $Y$ there is a nonempty open interval, which meets $D$). It is nonempty, since
$Y\ne\set{a,b}$ (a two-point Hausdorff space is not connected). By Cantor's back-and-forth theorem
there is an order isomorphism $\varphi\colon D'\to\mathbb Q\cap(0,1)$. Define
$F(y)=\sup\set{\varphi(d):d\in D',\ d\prec y}$ (with $\sup\emptyset=0$). Then $F$ is
order preserving, strictly increasing (between $x\prec y$ there are two points of $D'$), $F(a)=0$
and $F(b)=1$. The preimage of an open ray $\set{t<s}$ or $\set{t>s}$ is a union of open rays of
$Y$, so $F$ is continuous. $F(Y)$ is a connected subset of $I$ containing $0$ and $1$, so
$F$ is onto. A continuous bijection from a compact space onto a Hausdorff space is a homeomorphism;
take its inverse.
\end{proof}

\formal{\lean{HPC.Moore.ac_exists_strictMono}, \lean{HPC.Moore.ac_homeomorph}.}

\begin{corollary}[Arc characterization]\label{cor:arc-char}
Let $X$ be metrizable, $A\ne B$ in $X$ and $C\subseteq X$. Then $C$ is an arc from $A$ to $B$ if and
only if $C$ is compact, connected, contains $A$ and $B$, and no connected proper subset of $C$
contains both $A$ and $B$.
\end{corollary}

\begin{proof}
One direction was noted above. Conversely, $C$ with the subspace topology is compact metrizable,
hence separable, and Theorem~\ref{thm:irr-arc} applies.
\end{proof}

\formal{\lean{HPC.Moore.isIrreducibleConnected_iff_arc}.}

\section{Cut-point arcs}

\begin{definition}\label{def:cutarc}
A set $H$ is a \emph{cut-point arc from $A$ to $B$} if $A\ne B$, $H$ is compact and connected,
$A,B\in H$, and $H\setminus\set x$ is not connected for every $x\in H\setminus\set{A,B}$.
\end{definition}

\formal{\lean{HPC.Moore.IsCutArc}.}

This is Moore's definition of an arc in his 1920 and 1925 papers (Chapter~\ref{ch:elements}).

\begin{proposition}\label{prop:cutarc}
Let $X$ be a $T_1$ space.
\begin{enumerate}[label=(\alph*)]
\item A cut-point arc from $A$ to $B$ is irreducible between $A$ and $B$.
\item If $X$ is metrizable, a cut-point arc from $A$ to $B$ is an arc from $A$ to $B$. Conversely,
the image of an embedding $f\colon I\to X$ is a cut-point arc from $f(0)$ to $f(1)$, and so is
$f([s,t])$ for $s<t$.
\item If $f\colon I\to X$ is an embedding and $L\subseteq f(I)$ is closed, connected and has two
points, then $L=f([s,t])$ for some $s<t$. If $f([s,t])$ is a cut-point arc from $a$ to $b$, then
$\set{a,b}=\set{f(s),f(t)}$.
\item \emph{(Nadler 6.7.)} If $K$ is a continuum and $H\subsetneq K$ is connected, some $x\in
K\setminus H$ has $K\setminus\set x$ connected.
\end{enumerate}
\end{proposition}

\begin{proof}
(d) is \vthree{Theorem~3.7}.
(a) Let $L\subsetneq H$ be connected with $A,B\in L$, and take $x\in H\setminus L$. By (d) applied to
the continuum $H$ and the connected set $L$ we may take $x$ with $H\setminus\set x$ connected; then
$x\ne A,B$ (as $A,B\in L$), against the definition.
(b) By (a) and Corollary~\ref{cor:arc-char}. For the converse: if $x=f(u)$ with $s<u<t$, then
$f([s,t])\setminus\set x=f([s,u))\cup f((u,t])$, two disjoint nonempty sets, each closed in the
union (they are compact images); so it is not connected.
(c) $f^{-1}(L)$ is a compact connected subset of $I$ with two points, i.e.\ an interval $[s,t]$,
$s<t$. In $f([s,t])$ removing $f(s)$ or $f(t)$ leaves a connected set, and removing any other point
does not; so the two points that can be removed without disconnecting are the ends.
\end{proof}

\formal{(a) \lean{HPC.Moore.IsCutArc.irreducible}; (b) \lean{HPC.Moore.IsCutArc.exists_embedding},
\lean{HPC.Moore.isCutArc_range}, \lean{HPC.Moore.isCutArc_image_Icc}; (c)
\lean{HPC.Moore.eq_image_Icc_of_subset_range}, \lean{HPC.Moore.ends_of_image_Icc}; (d)
\lean{HPC.Moore.exists_nonCut_diff}.}

\section{Simple closed curves}

\begin{theorem}[Moore 1920, Theorem~4]\label{thm:circle-char}
Let $J$ be a compact connected subset of a Hausdorff space with at least two points, such that
$J\setminus\set{x,y}$ is not connected for any two distinct $x,y\in J$. Then for any distinct
$p,q\in J$ there are cut-point arcs $a,b$ from $p$ to $q$ with $a\cup b=J$ and $a\cap b=\set{p,q}$.
\end{theorem}

\begin{proof}
Work in the space $Y=J$, which is compact, connected and Hausdorff.

\emph{Step 1: no point of $Y$ is a cut point.} Suppose $Y\setminus\set x=U\cup V$ with $U,V$
disjoint, open and nonempty. By Lemma~\ref{lem:insert-conn}, $U\cup\set x=Y\setminus V$ is a
continuum with at least two points, so it has two non-cut points (\vthree{Theorem~3.6}), one of
them, $u$, different from $x$; so $u\in U$ and $(U\cup\set x)\setminus\set u$ is connected.
Likewise there is $v\in V$ with $(V\cup\set x)\setminus\set v$ connected. Then
$Y\setminus\set{u,v}$ is the union of these two connected sets, which share $x$; so it is
connected, against the hypothesis.

\emph{Step 2: the halves.} $Y\setminus\set{p,q}$ is nonempty (otherwise $\set q=Y\setminus\set p$
would be a proper nonempty clopen set) and not connected, so $Y\setminus\set{p,q}=U\cup V$ with
$U,V$ disjoint, open and nonempty. Put $a=Y\setminus V=U\cup\set{p,q}$ and $b=Y\setminus
U=V\cup\set{p,q}$. Then $a,b$ are closed, $a\cup b=Y$ and $a\cap b=\set{p,q}$. We show that $a$ is
a cut-point arc from $p$ to $q$; the argument for $b$ is the same.

\emph{$a$ is connected.} Suppose $a=N\cup K$ with $N,K$ disjoint, closed and nonempty. If $p,q\in
N$, then $K\subseteq U$, so $K=U\setminus N$ is open in $Y$, and $K$ is a nonempty proper clopen
subset of $Y$: impossible. If $p\in K$ and $q\in N$, then one of the two pieces has a point other
than its marked point: if $K=\set p$ then $N\supseteq U\ne\emptyset$ has points other than $q$.
Suppose $K$ has a point $k\ne p$ (the other case is symmetric, exchanging $p$ and $q$). Then
\[
Y\setminus\set p=(K\setminus\set p)\cup\big((N\cup b)\setminus\set p\big),
\]
both sets are closed in $Y\setminus\set p$, they are disjoint (since $K\cap b\subseteq a\cap
b=\set{p,q}$ and $q\in N$), and nonempty ($k$, resp.\ $q$). This contradicts Step~1.

\emph{Each $x\in U$ cuts $a$.} Suppose $a\setminus\set x$ were connected. Take $y\in V$; we show
$S=Y\setminus\set{x,y}$ is connected, contradicting the hypothesis. Let $S\subseteq u\cup v$ with
$u,v$ open and $S\cap u\cap v=\emptyset$. The connected set $a\setminus\set x\subseteq S$ lies in
one of them, say $u$. Put $Q=S\cap v$, an open set. Every point of $S$ outside $V$ lies in
$a\setminus\set x\subseteq u$, so $Q\subseteq V$. If $z\ne y$ is in $\cl Q$, then $z\notin u$
(as $Q$ misses the open set $u$), $z\in\cl V\subseteq b$, so $z\ne x$ (as $x\in U$ misses $b$); thus
$z\in S\setminus u\subseteq v$ and $z\in Q$. So $Q$ is closed and open in $Y\setminus\set y$, which
is connected by Step~1. As $p\in Y\setminus\set y$ and $p\notin Q$, $Q=\emptyset$. Hence
$S\subseteq u$, and $S$ is connected.

The points of $a\setminus\set{p,q}$ are the points of $U$; so $a$ is a cut-point arc.
\end{proof}

\formal{\lean{HPC.Moore.circle_noCut}, \lean{HPC.Moore.circle_half}, \lean{HPC.Moore.circle_char}.}

\begin{proposition}\label{prop:scc-cut}
Let $J=X_1\cup X_2$ be a simple closed curve ($X_1,X_2$ arcs from $p$ to $q$, $X_1\cap
X_2=\set{p,q}$). Then $J\setminus\set{x,y}$ is not connected for any distinct $x,y\in J$, and $J$
is homeomorphic to the circle.
\end{proposition}

\begin{proof}
Let $f$, $g$ parametrize $X_1$, $X_2$ from $p$ to $q$. If $x=f(s)$, $y=f(t)$ with $s<t$ lie on one
arc, the closed sets $f([s,t])$ and $f([0,s]\cup[t,1])\cup X_2$ cover $J$ and meet only in $\set{x,y}$,
and each contains a point of $J\setminus\set{x,y}$. If $x=f(s)\in X_1\setminus X_2$ and $y=g(t)\in
X_2\setminus X_1$, use $f([0,s])\cup g([0,t])$ and $f([s,1])\cup g([t,1])$. In both cases
$J\setminus\set{x,y}$ is the union of two disjoint nonempty relatively closed sets.

For the second claim, the map $h\colon[0,1]\to J$ equal to $f(2u)$ for $u\le\frac12$ and to
$g(2-2u)$ for $u\ge\frac12$ is continuous, onto, and injective except that $h(0)=h(1)=p$; it
induces a continuous bijection from the circle $[0,1]/\set{0\sim1}$ onto $J$, which is a
homeomorphism by compactness.
\end{proof}

\formal{\lean{HPC.Moore.Gen.tscc_cut}; \lean{HPC.Jones.glue_circle},
\lean{HPC.Jones.tscc_homeomorph_circle}.}

So in a metrizable space the following are equivalent for a set $J$ with at least two points:
$J$ is a simple closed curve; $J$ is homeomorphic to the circle; $J$ is a continuum cut by every
pair of its points (Theorem~\ref{thm:circle-char}, Propositions~\ref{prop:cutarc}(b)
and~\ref{prop:scc-cut}).

\section{Limits of arcs and replacement of sub-arcs}

Two constructions with infinitely many arcs are needed in Chapter~\ref{ch:curves}.

\begin{theorem}[Limit arcs]\label{thm:limit-arc}
Let $X$ be metrizable, $B\in X$, and let $\varepsilon_n$ ($n\ge0$) be arcs from $y_n$ to $y_{n+1}$
such that $\varepsilon_n\cap\varepsilon_{n+1}=\set{y_{n+1}}$, $\varepsilon_n\cap\varepsilon_m=\emptyset$
for $m\ge n+2$, $B\notin\varepsilon_n$, and the $\varepsilon_n$ converge to $B$ (every
neighbourhood of $B$ contains $\varepsilon_n$ for all large $n$). Then
$\tau=\set B\cup\bigcup_n\varepsilon_n$ is an arc from $y_0$ to $B$.
\end{theorem}

\begin{proof}
$\tau$ is compact: an open cover has a member containing $B$, which contains all but finitely many
$\varepsilon_n$, and the remaining ones are compact. The same argument shows that
$\set B\cup\bigcup_{m\ge k}\varepsilon_m$ is compact for every $k$. The sets
$\Gamma_n=\varepsilon_0\cup\dots\cup\varepsilon_n$ are arcs from $y_0$ to $y_{n+1}$ (gluing,
Lemma~\ref{lem:arc-ops}(b)); they increase, so their union is connected, and $\tau$, which lies in its
closure, is connected.

$\tau$ is irreducible between $y_0$ and $B$. Let $z\in\tau\setminus\set{y_0,B}$ and let $n$ be least
with $z\in\varepsilon_n$. Let $g$ parametrize the arc $\Gamma_{n+1}$ from $y_0$ to $y_{n+2}$; then
$z=g(c)$ with $0<c<1$ ($z\ne y_{n+2}$ since $\varepsilon_{n+2}\cap\varepsilon_n=\emptyset$). The sets
$F_1=g([0,c])$ and $F_2=g([c,1])\cup\set B\cup\bigcup_{m\ge n+2}\varepsilon_m$ are compact, cover
$\tau$, and $F_1\cap F_2=\set z$ ($\varepsilon_m$, $m\ge n+2$, meets $\Gamma_{n+1}$ at most in
$y_{n+2}=g(1)$). So $\tau\setminus\set z$ is the union of the disjoint relatively closed sets
$F_1\setminus\set z\ni y_0$ and $F_2\setminus\set z\ni B$, and every connected subset of $\tau$
containing $y_0$ and $B$ contains $z$. Hence such a set is all of $\tau$, and
Corollary~\ref{cor:arc-char} applies.
\end{proof}

\formal{\lean{HPC.Moore.Gen.tarc_of_seq}; with trimming of arcs that meet more than consecutively,
\lean{HPC.Moore.Gen.exists_tarc_of_seq}.}

In applications the arcs $\beta_n$ come with $\beta_n\cap\beta_m=\emptyset$ for $m\ge n+2$ only;
cutting each $\beta_n$ at its first point on $\beta_{n+1}$ (Lemma~\ref{lem:arc-ops}(a)) produces arcs
$\varepsilon_n\subseteq\beta_n$ satisfying the hypotheses, and so $\set B\cup\bigcup\beta_n$ contains
an arc from $x_0$ to $B$.

\begin{theorem}[Replacement of a null family of sub-arcs]\label{thm:arc-replace}
Let $X$ be a metric space and $f\colon I\to X$ an embedding. Let $(u_i,v_i)$, $i\in\iota$, be
pairwise disjoint open subintervals of $I$, and $Z_i$ arcs from $f(u_i)$ to $f(v_i)$ whose open
arcs are pairwise disjoint and disjoint from $f(I)$. Suppose the family is \emph{null}: for every
$\varepsilon>0$ only finitely many $Z_i$ contain a point at distance $\ge\varepsilon$ from
$f(u_i)$. Then
\[
D=f\Big(I\setminus\bigcup_i(u_i,v_i)\Big)\cup\bigcup_i Z_i
\]
is an arc from $f(0)$ to $f(1)$.
\end{theorem}

\begin{proof}
Let $z_i\colon I\to Z_i$ parametrize $Z_i$ from $f(u_i)$ to $f(v_i)$, and define $h\colon I\to X$
by $h(t)=z_i\big((t-u_i)/(v_i-u_i)\big)$ for $t\in(u_i,v_i)$ and $h(t)=f(t)$ for $t$ in no interval.
Then $h(I)=D$, and $h$ is injective by the disjointness hypotheses. It remains to show that $h$ is
continuous; then $D$ is an arc. On each closed interval $[u_i,v_i]$ the map $h$ is continuous (the
two definitions agree at $u_i,v_i$). Let $s_k\to t$; we show $h(s_k)\to h(t)$, and may assume that
$t$ is in no open interval and that every $s_k$ lies in some interval $(u_{i_k},v_{i_k})$ (the
$s_k$ in no interval satisfy $h(s_k)=f(s_k)\to f(t)$). If some index $i$ occurs infinitely often,
then $t\in\set{u_i,v_i}$ and along those $k$, $h(s_k)\to h(t)$ by continuity on $[u_i,v_i]$. So we
may assume that every index occurs only finitely often. Since the intervals are disjoint, only
finitely many have length $\ge\delta$, for every $\delta>0$; so $v_{i_k}-u_{i_k}\to0$, hence
$u_{i_k}\to t$ and $f(u_{i_k})\to f(t)$. By nullity, $\sup_{x\in Z_{i_k}}d(x,f(u_{i_k}))\to0$. As
$h(s_k)\in Z_{i_k}$, $h(s_k)\to f(t)=h(t)$.
\end{proof}

\formal{\lean{HPC.Moore.Gen.arc_replace}.}

\section*{Exercises}
\addcontentsline{toc}{section}{Exercises}

\begin{exercise}
Show that in Lemma~\ref{lem:arc-ops}(c) the hypothesis $P\ne R$ is needed, and that the arc from $P$
to $R$ need not be unique.
\end{exercise}

\begin{exercise}
Show that the conclusion of Corollary~\ref{cor:arc-char} fails for the ``long line'' compactified
by two end points (a compact connected Hausdorff space irreducible between its end points which is
not metrizable). Which step of the proof of Theorem~\ref{thm:irr-arc} fails?
\end{exercise}

\begin{exercise}
Let $J$ be a simple closed curve and $p,q\in J$ distinct. Show that there are exactly two arcs in
$J$ from $p$ to $q$, and that their union is $J$.
\end{exercise}

\begin{exercise}
Show that the null hypothesis in Theorem~\ref{thm:arc-replace} cannot be omitted: replace the
intervals $(1/(n+1),1/n)$ of the segment $[0,1]\times\set0$ by semicircles of radius $1$ (suitably
tilted so that they are disjoint) and show that the result is not compact.
\end{exercise}

%======================================================================
\chapter{Decompositions of the plane into non-separating continua}\label{ch:usc}
%======================================================================

This chapter begins the study of Moore's 1925 paper \emph{Concerning upper semi-continuous
collections of continua}. Its main theorem (Theorem~\ref{thm:m25-22} in Chapter~\ref{ch:quotient})
says that if the plane is cut up into compact connected pieces, none of which separates the plane,
in an upper semicontinuous way, then the space of pieces is again a plane. Here we set up the space
of pieces as a quotient space, show that it is metrizable and locally connected, translate Moore's
notions (limit element, domain, boundary) into topology, and prove Moore's Theorems~1--4. The
plane tools are those of Volume~4: Brouwer's theorem, Janiszewski's theorem and enclosing polygons.

\section{The standing hypothesis}

As in Volume~4, $\R^2$ carries the sup metric $d$, so that balls are squares; for a set $A$ and a
point $z$, $\dist(z,A)=\inf_{a\in A}d(z,a)$.

\begin{definition}\label{def:moore-decomp}
A \emph{Moore decomposition} of the plane is an equivalence relation on $\R^2$ whose classes, the
\emph{elements}, satisfy:
\begin{enumerate}[label=(\arabic*)]
\item every element is compact and connected;
\item no element separates the plane;
\item \emph{(upper semicontinuity, Moore's form)} for every element $g$ and every $e>0$ there is
$\delta>0$ such that every element containing a point at distance $<\delta$ from $g$ lies within
distance $<e$ of $g$ (every point of it is at distance $<e$ from $g$).
\end{enumerate}
We write $\el x$ for the element containing $x$, $G$ for the set of elements with the quotient
topology, and $\pi\colon\R^2\to G$, $x\mapsto\el x$, for the quotient map. For $H\subseteq G$,
$\pts H=\pi^{-1}(H)$ is the \emph{point set} of $H$ (Moore writes $\bar H$). $H$ is \emph{bounded} if
$\pts H$ is bounded.
\end{definition}

\formal{\lean{HPC.Moore.MooreDecomp}, \lean{HPC.Moore.elt}, \lean{HPC.Moore.ptsOf}.}

By definition of the quotient topology, $H$ is open (closed) in $G$ if and only if $\pts H$ is open
(closed) in $\R^2$, and $\pi(\pts H)=H$.

The two extreme examples: the decomposition into points (every element a single point), for which
$G$ is the plane itself (Section~\ref{sec:points}); and, for a plane continuum $M$ with a monotone map
$q\colon M\to Y$ whose fibres do not separate the plane, the decomposition into the fibres of $q$ and
the single points off $M$ (Section~\ref{sec:monotone-quotient}).

\begin{definition}\label{def:tube}
For $U\subseteq\R^2$ the \emph{tube} of $U$ is $\mathrm{tube}(U)=\set{g\in G:g\subseteq U}$.
\end{definition}

\formal{\lean{HPC.Moore.tube}.}

\begin{proposition}\label{prop:closed-map}
\leavevmode
\begin{enumerate}[label=(\alph*)]
\item $\pi$ is a closed map. Consequently the tube of an open set is open.
\item The tubes of the open $\varepsilon$-neighbourhoods $N_\varepsilon(g)=\set{z:\dist(z,g)<\varepsilon}$,
$\varepsilon>0$, form a neighbourhood base at the element $g$.
\item Conversely, for a decomposition into compact sets, if $\pi$ is closed then Moore's condition
(3) holds.
\end{enumerate}
\end{proposition}

\begin{proof}
(a) Let $F\subseteq\R^2$ be closed; we show that $\pi^{-1}(\pi(F))$, the union of the elements meeting
$F$, is closed. Let $x$ lie in no element meeting $F$. The compact set $\el x$ misses the closed set
$F$, so $\dist(z,\el x)\ge e$ for all $z\in F$, for some $e>0$. Take $\delta$ as in (3). If an element
$g$ contains a point at distance $<\delta$ from $\el x$, then $g\subseteq N_e(\el x)$, so $g$ misses
$F$. So $N_\delta(\el x)$ is a neighbourhood of $x$ missing $\pi^{-1}(\pi(F))$. For the tube:
$\mathrm{tube}(U)=G\setminus\pi(\R^2\setminus U)$.

(b) If $V\ni g$ is open in $G$, $\pts V$ is an open set containing the compact set $g$, so it contains
some $N_\varepsilon(g)$, and $g\in\mathrm{tube}(N_\varepsilon(g))\subseteq V$.

(c) Given $g$ and $e>0$, $\mathrm{tube}(N_e(g))$ is open, so its point set is an open set containing
the compact $g$, and contains some $N_\delta(g)$. An element meeting $N_\delta(g)$ lies in the tube,
i.e.\ in $N_e(g)$.
\end{proof}

\formal{\lean{HPC.Moore.isClosedMap_mk_plane}, \lean{HPC.Moore.isOpen_tube},
\lean{HPC.Moore.mem_nhds_iff_tube}, \lean{HPC.Moore.usc_of_isClosedMap}.}

So Moore's distance form of upper semicontinuity agrees, for decompositions into compact sets,
with the topological one of \vthree{Definition~6.1}.

\begin{theorem}\label{thm:quot-metrizable}
$G$ is a $T_1$, regular, second countable space; hence it is metrizable.
\end{theorem}

\begin{proof}
\emph{$T_1$.} Points of $G$ are closed because elements are closed.

\emph{Regular.} Let $g\in G$ and $F\subseteq G$ closed with $g\notin F$. The closed sets $g$ and $\pts F$
of the plane are disjoint, so they have disjoint open neighbourhoods $U\supseteq g$, $V\supseteq\pts
F$. Then $\mathrm{tube}(U)\ni g$ and $\mathrm{tube}(V)\supseteq F$ are open and disjoint (an element in
both would lie in $U\cap V=\emptyset$).

\emph{Second countable.} Let $\mathcal B$ be the countable family of finite unions of open squares
with rational centres and radii. Given $g\in V$ open, the compact set $g\subseteq\pts V$ is covered by
finitely many rational squares inside $\pts V$; their union $B\in\mathcal B$ satisfies $g\in
\mathrm{tube}(B)\subseteq V$. So the countably many open sets $\mathrm{tube}(B)$, $B\in\mathcal B$, form a
basis.

A regular second countable $T_1$ space is metrizable (Urysohn's metrization theorem).
\end{proof}

\formal{\lean{HPC.Moore.t1Space_quot}, \lean{HPC.Moore.regularSpace_quot},
\lean{HPC.Moore.secondCountable_quot}, \lean{HPC.Moore.quotient_metrizable_plane}.}

\section{Connectedness and components}

\begin{proposition}\label{prop:pts-conn}
\leavevmode
\begin{enumerate}[label=(\alph*)]
\item If $C\subseteq G$ is connected, so is $\pts C$.
\item \emph{(Components correspond.)} For every $W\subseteq G$ and $x\in\R^2$,
$\pts{(\text{component of }\el x\text{ in }W)}=\text{component of }x\text{ in }\pts W$.
\item Components of open subsets of $G$ are open; $G$ is locally connected.
\item If $C\subseteq G$ is open, $\fr C=\pi(\fr\pts C)$.
\end{enumerate}
\end{proposition}

\begin{proof}
(a) Let $\pts C\subseteq U\cup V$ with $U,V$ open and $\pts C\cap U\cap V=\emptyset$. Every element
$g\in C$ is connected and lies in $U\cup V$, so (being covered by the two disjoint relatively open
sets $g\cap U$, $g\cap V$) it lies in $U$ or in $V$. So $C\subseteq\mathrm{tube}(U)\cup\mathrm{tube}(V)$,
two open sets (Proposition~\ref{prop:closed-map}) whose intersection misses $C$. As $C$ is connected it
lies in one of them, and then $\pts C$ lies in $U$ or in $V$.

(b) Let $K$ be the component of $\el x$ in $W$ and $L$ the component of $x$ in $\pts W$. By (a),
$\pts K$ is a connected subset of $\pts W$ containing $x$, so $\pts K\subseteq L$. Conversely $\pi(L)$
is connected, lies in $W$ and contains $\el x$, so $\pi(L)\subseteq K$ and $L\subseteq\pts K$.

(c) If $W$ is open, $\pts W$ is open, its components are open (balls are connected), and by (b) the
point set of each component of $W$ is one of them.

(d) $\pts C$ is open and is a union of elements. Since $\pi$ is closed and continuous,
$\cl C=\pi(\cl{\pts C})$. Hence $\fr C=\cl C\setminus C=\pi(\cl{\pts C})\setminus C=\pi(\cl{\pts
C}\setminus\pts C)=\pi(\fr\pts C)$, using that $\pts C$ is a union of elements.
\end{proof}

\formal{\lean{HPC.Moore.isPreconnected_ptsOf}, \lean{HPC.Moore.ptsOf_connectedComponentIn},
\lean{HPC.Moore.isOpen_connectedComponentIn_quot}, \lean{HPC.Moore.locallyConnectedSpace_quot},
\lean{HPC.Moore.frontier_eq_image}.}

\begin{corollary}\label{cor:quot-tools}
\leavevmode
\begin{enumerate}[label=(\alph*)]
\item \emph{(Brouwer for elements.)} If $F\subseteq G$ is closed and connected and $q\notin F$, the
frontier of the component of $q$ in $G\setminus F$ is connected.
\item \emph{(Janiszewski for elements.)} Let $A,B\subseteq G$ be closed and bounded with $A\cap B$
connected or empty, and $x,y\notin A\cup B$. If neither $A$ nor $B$ separates $x$ from $y$ in $G$,
then $A\cup B$ does not.
\item \emph{(The unbounded component.)} If $H\subseteq G$ is bounded, exactly one component of
$G\setminus H$ is unbounded.
\item For every element $g$, $G\setminus\set g$ is connected.
\end{enumerate}
\end{corollary}

\begin{proof}
(a) $\pts F$ is closed and connected (Proposition~\ref{prop:pts-conn}(a)). If $C$ is the component of
$q$ in $G\setminus F$, then $\pts C$ is a component of $\R^2\setminus\pts F$ (Proposition~\ref{prop:pts-conn}(b)),
whose frontier is connected by Brouwer's theorem (\vfour{Theorem~1.12}). By
Proposition~\ref{prop:pts-conn}(d), $\fr C$ is its image under $\pi$, hence connected.

(b) By Proposition~\ref{prop:pts-conn}(b), $A$ separates $x$ from $y$ in $G$ exactly when $\pts A$
separates any point of $x$ from any point of $y$ in the plane. $\pts A$, $\pts B$ are compact, and
$\pts A\cap\pts B=\pts{(A\cap B)}$ is connected or empty by Proposition~\ref{prop:pts-conn}(a).
Apply Janiszewski's theorem (\vfour{Theorem~2.7}).

(c) Choose $R$ with $\pts H\subseteq\bar B(0,R)$. The set $W=\pi(\R^2\setminus\bar B(0,R))$ is
connected (the complement of a square is connected, \vfour{Lemma~2.2}), it misses $H$ (an element of
$H$ lies in $\pts H\subseteq\bar B(0,R)$), and its point set is unbounded. So $W$ lies in one
component $C_\infty$ of $G\setminus H$, which is unbounded. An unbounded component has point set
meeting $\R^2\setminus\bar B(0,R)$, so it meets $W\subseteq C_\infty$ and equals $C_\infty$.

(d) $\pts{(G\setminus\set g)}=\R^2\setminus g$ is connected because $g$ does not separate the plane;
its image is $G\setminus\set g$.
\end{proof}

\formal{\lean{HPC.Moore.quot_brouwer}, \lean{HPC.Moore.quot_janiszewski},
\lean{HPC.Moore.m25_unbounded_unique}, \lean{HPC.Moore.isConnected_compl_singleton_quot}.}

\section{Moore's notions are the topological ones}

Moore works directly with elements and distances. An element $p$ is a \emph{limit element} of a set
$K\subseteq G$ if for every $e>0$ some element of $K$ other than $p$ lies within distance $e$ of $p$
(every point of it is at distance $<e$ from $p$). A \emph{domain} of elements is a connected set $D$
such that each element of $D$ is not a limit element of $G\setminus D$. The \emph{boundary} of
$H\subseteq G$ consists of the elements of $H$ that are limit elements of $G\setminus H$ together with
the elements not in $H$ that are limit elements of $H$.

\begin{proposition}\label{prop:moore-notions}
$p$ is a limit element of $K$ if and only if $p$ is an accumulation point of $K$ in $G$. Hence the
domains are the connected open subsets of $G$, and the boundary of $H$ is its frontier $\fr H$.
\end{proposition}

\begin{proof}
An element $k$ lies within distance $e$ of $p$ exactly when $k\in\mathrm{tube}(N_e(p))$, and these tubes
form a neighbourhood base at $p$ (Proposition~\ref{prop:closed-map}(b)).
\end{proof}

\formal{\lean{HPC.Moore.isLimitElt_iff}, \lean{HPC.Moore.m25_domain_iff},
\lean{HPC.Moore.m25_boundary_iff}.}

\begin{remark}
Read literally, Moore's definition of boundary (p.~417) has ``a limit element of $H$'' in both
clauses, which would make the first clause empty; the reading above (``of $G-H$'' in the first
clause) is the intended one, and it gives the frontier.
\end{remark}

\section{Moore's Theorems 1--4}

\begin{theorem}[Moore, Theorem~1]\label{thm:m25-1}
Let $K\subseteq\R^2$ and let $H=\pi(K)$ be the set of elements meeting $K$. If $K$ is closed, so is
$H$; if $K$ is connected, so is $H$.
\end{theorem}

\begin{proof}
$\pi$ is closed (Proposition~\ref{prop:closed-map}) and continuous.
\end{proof}

\formal{\lean{HPC.Moore.m25_theorem1}.}

\begin{definition}\label{def:moore-region}
A \emph{Moore region} is an open connected bounded set $R\subseteq G$ whose frontier is connected.
\end{definition}

\formal{\lean{HPC.Moore.mooreRegion}, \lean{HPC.Moore.spaceOfElts}.}

\begin{theorem}[Moore, Theorem~3]\label{thm:m25-3}
For every element $p$ and every $e>0$ there is a Moore region $R\ni p$ all of whose elements lie in
$N_e(p)$.
\end{theorem}

\begin{proof}
$p$ is a continuum not separating the plane, inside the open set $N_e(p)$. By the enclosing polygon
theorem (\vfour{Theorem~4.14}) there is a simple closed polygon $\Gamma$ missing $p$ with
$p\subseteq\inn\Gamma\subseteq N_e(p)$. Let $H=\pi(\Gamma)$, the set of elements meeting $\Gamma$; it is
closed and connected (Theorem~\ref{thm:m25-1}), and $p\notin H$. Let $R$ be the component of $p$ in
$G\setminus H$. It is open and connected (Proposition~\ref{prop:pts-conn}(c)), and its frontier is
connected by Corollary~\ref{cor:quot-tools}(a). By Proposition~\ref{prop:pts-conn}(b), $\pts R$ is the
component of the plane set $p$ in $\R^2\setminus\pts H$; it misses $\Gamma$, is connected and meets
$\inn\Gamma$, so it lies in $\inn\Gamma\subseteq N_e(p)$. So $R$ is bounded and every element of $R$ lies in
$N_e(p)$.
\end{proof}

\formal{\lean{HPC.Moore.exists_region_in_tube}, \lean{HPC.Moore.m25_theorem3}.}

\begin{theorem}[Moore, Theorem~4]\label{thm:m25-4}
The Moore regions form a basis of the topology of $G$. In particular $p$ is a limit element of $K$
if and only if every Moore region containing $p$ contains an element of $K$ other than $p$; and the
closure of a set in Moore's sense (the set together with its limit elements) is its topological
closure.
\end{theorem}

\begin{proof}
Moore regions are open, and by Theorem~\ref{thm:m25-3} and Proposition~\ref{prop:closed-map}(b) every
neighbourhood of $p$ contains one containing $p$. The rest follows from
Proposition~\ref{prop:moore-notions}.
\end{proof}

\formal{\lean{HPC.Moore.isTopologicalBasis_mooreRegion}, \lean{HPC.Moore.m25_theorem4},
\lean{HPC.Moore.spaceOfElts_topology}, \lean{HPC.Moore.cl_eq_closure_quot}.}

\begin{theorem}[Moore, Theorem~2]\label{thm:m25-2}
Let $M\subseteq G$ be closed, connected and bounded, $D$ a bounded component of $G\setminus M$, $E$
the unbounded component of $G\setminus\fr D$ and $K=\fr E$ (the \emph{outer boundary} of $D$). Then
$K$ is closed and connected, and $K\setminus\set p$ is connected for every $p\in K$.
\end{theorem}

\begin{proof}
$D$ and $E$ are open and connected; $\fr D\subseteq M$ and $K=\fr E\subseteq\fr D$ (the frontier of a
component of the complement of a closed set lies in that set). $K$ is closed, and it is connected by
Corollary~\ref{cor:quot-tools}(a) applied twice: $\fr D$ is connected, and then so is the frontier of
the component $E$ of its complement. $K$ is bounded, as $K\subseteq M$. Fix $x\in D$ and $y\in E$.

\emph{$D\cap E=\emptyset$.} $E$ is connected and misses $\fr D$; if it met $D$ it would lie in $D$
(Lemma~\ref{lem:conn-frontier}), and would be bounded.

\emph{$K\setminus\set p\neq\emptyset$.} Otherwise $\fr E=\set p$. Then $E$ is open and closed in the
connected set $G\setminus\set p$ (Corollary~\ref{cor:quot-tools}(d)), so $E=G\setminus\set p\supseteq D$
($p\notin D$ since $p\in\fr D$ and $D$ is open), against $D\cap E=\emptyset$.

\emph{$K\setminus\set p$ is connected.} Suppose $K\setminus\set p=H\cup N$ with $H,N$ nonempty, disjoint
and relatively closed. Then $H\cup\set p$ and $N\cup\set p$ are closed (a limit point of $H$ lies in
the closed set $K$ and not in $N$), bounded, and meet in $\set p$.

\emph{Claim: $F=H\cup\set p$ does not separate $x$ from $y$.} Suppose it does, and let $D_x$ be the
component of $x$ in $G\setminus F$; then $y\notin D_x$, and $\fr D_x\subseteq F$. Since $D$ is connected
and misses $K\supseteq F$, $D\subseteq D_x$. Take $q\in N$. Then $q\in K\subseteq\fr D\subseteq\cl D\subseteq
\cl{D_x}$ and $q\notin F\supseteq\fr D_x$, so $q\in D_x$. As $q\in\fr E$, the open set $D_x$ meets $E$.
$E$ is connected and misses $K\supseteq\fr D_x$, so $E\subseteq D_x$ (Lemma~\ref{lem:conn-frontier}), and
$y\in D_x$: a contradiction.

Likewise $N\cup\set p$ does not separate $x$ from $y$. Their intersection $\set p$ is connected, so by
Janiszewski's theorem for elements (Corollary~\ref{cor:quot-tools}(b)) $K$ does not separate $x$ from
$y$: the component $C$ of $x$ in $G\setminus K$ contains $y$. But $C$ is connected, misses $K=\fr E$
and meets $E$ (at $y$), so $C\subseteq E$ and $x\in E$, against $D\cap E=\emptyset$.
\end{proof}

\formal{\lean{HPC.Moore.m25_theorem2}.}

\begin{remark}\label{rem:m25-2-printed}
Moore defines a \emph{complementary domain} of a closed set $H$ as a domain whose boundary lies
in $H$, not necessarily a component of $G\setminus H$. With that definition Theorem~2 is false: in the
decomposition into points, let $M$ be the closed unit disc, $D$ the open disc minus the segment
$T=[0,\frac12]\times\set0$, and $E$ the plane minus $T$; then $E$ is an unbounded ``complementary
domain'' of $\fr D$, its boundary is $T$, and removing $(\frac14,0)$ disconnects $T$
(\lean{HPC.Moore.not_m25_theorem2_printed}). With complementary domains read as components, as above,
the theorem holds; this is the form used later (in the proof of Theorem~9, Chapter~\ref{ch:elements}).
\end{remark}

\section{The claim of p.~420}\label{sec:p420}

Moore states without proof that the space of elements, with the Moore regions as regions, satisfies
Axioms~1, 2, 4 and~5 of his 1916 paper and its Theorem~4, and that consequently Theorems~1--15 of
that paper hold in it. Moore's axioms are the subject of Chapter~\ref{ch:axioms}; we record here the
part that only needs the results above. For the definitions see Section~\ref{sec:axioms}.

\begin{proposition}\label{prop:p420}
The space of elements with the Moore regions satisfies Axioms~1, 2, 4 and~5 of
Definition~\ref{def:moore-axioms}, and every nonempty Moore region is infinite.
\end{proposition}

\begin{proof}[Proof sketch]
\emph{Axiom~1.} $G$ is metrizable (Theorem~\ref{thm:quot-metrizable}); fix a metric. For each $k$ the
set of elements meeting the square $\bar B(0,k)$ is compact (its point set is closed and bounded, by
upper semicontinuity), so it is covered by finitely many Moore regions lying in balls of radius
$1/(k+1)$ (Theorems~\ref{thm:m25-3} and~\ref{thm:m25-4}). Listing these finite families one after the
other gives a sequence $K_1,K_2,\dots$ with the properties of Axiom~1; Moore's closure is the
topological closure (Theorem~\ref{thm:m25-4}).
\emph{Axiom~2.} Moore regions are connected, and Moore's connectedness is the topological one once
Axiom~1 holds (Proposition~\ref{prop:bridge}).
\emph{Axiom~4.} The closure of a Moore region $R$ is the image of the compact set $\cl{\pts R}$.
\emph{Axiom~5.} The elements containing the points $(n,0)$, $n\in\N$, are infinitely many (only finitely
many integer points lie in a bounded element), and have no limit element: a tube about any element
meets only finitely many of them.
\emph{Theorem~4 of 1916.} A finite nonempty region would be a compact open nonempty proper subset of the
connected space $G$.
\end{proof}

\formal{\lean{HPC.Moore.m25_axioms1245}, \lean{HPC.Moore.m25_FA_1_15}.}

The remaining axioms need more work and are proved, for a different system of regions, in
Chapter~\ref{ch:quotient}.

\section*{Exercises}
\addcontentsline{toc}{section}{Exercises}

\begin{exercise}
Show that the decomposition of the plane into the vertical lines is upper semicontinuous in the
topological sense, but violates condition (1) of Definition~\ref{def:moore-decomp}, and that its quotient
is the real line.
\end{exercise}

\begin{exercise}
Let $M$ be the closed unit square and decompose the plane into $M$ and single points. Verify the
standing hypothesis, and show directly (without Theorem~\ref{thm:m25-22}) that the quotient is
homeomorphic to the plane.
\end{exercise}

\begin{exercise}
Show that Proposition~\ref{prop:pts-conn}(a) fails if the elements are not assumed connected.
\end{exercise}

\begin{exercise}
In Theorem~\ref{thm:m25-2}, show by an example in the decomposition into points that $K$ may be a
proper subset of $\fr D$.
\end{exercise}

%======================================================================
\chapter{Arcs and simple closed curves of elements}\label{ch:elements}
%======================================================================

We keep the standing hypothesis of Chapter~\ref{ch:usc}: $G$ is the space of elements of a Moore
decomposition of the plane, a connected, locally connected metrizable space
(Theorem~\ref{thm:quot-metrizable}, Proposition~\ref{prop:pts-conn}). This chapter covers Moore's
Theorems~5--10 and~12. The central one is Theorem~9: \emph{an arc of elements does not separate
$G$}. With the general topology of Chapter~\ref{ch:arcs} most of the others become short.

\section{Moore's definitions}

\begin{definition}\label{def:arc-elt}
Let $A\ne B$ be elements.
\begin{itemize}
\item A \emph{simple continuous arc of elements from $A$ to $B$} is a closed, connected, bounded set
$H\subseteq G$ containing $A$ and $B$ such that $H\setminus\set x$ is not connected for every $x\in
H\setminus\set{A,B}$.
\item A \emph{simple closed curve of elements} is a closed, connected, bounded set $J\subseteq G$ with
at least two elements such that $J\setminus\set{x,y}$ is not connected for any two distinct $x,y\in
J$.
\end{itemize}
\end{definition}

\formal{\lean{HPC.Moore.IsArcElt}, \lean{HPC.Moore.IsSCCElt}.}

Moore's printed definition of a simple closed curve of elements (p.~417) asks only that $J$ be
bounded and disconnected by the omission of any two of its elements. Then two disjoint segments
(in the decomposition into points) form a ``simple closed curve'', and Theorems~10 and~11 fail for
it (\lean{HPC.Moore.not_m25_theorem10_printed}, \lean{HPC.Moore.not_m25_theorem11_printed}); a single
element must also be excluded (\lean{HPC.Moore.not_m25_theorem11_singleton}). The corrected
definition above adds ``closed, connected, with at least two elements''.

\begin{proposition}\label{prop:arc-elt-top}
\leavevmode
\begin{enumerate}[label=(\alph*)]
\item A closed bounded set of elements is compact.
\item $H$ is a simple continuous arc of elements from $A$ to $B$ if and only if $H$ is an arc from $A$
to $B$ in the space $G$ (Definition~\ref{def:arc}).
\item $J$ is a simple closed curve of elements if and only if $J$ is a simple closed curve in the
space $G$; equivalently, $J$ is homeomorphic to the circle.
\end{enumerate}
\end{proposition}

\begin{proof}
(a) If $H$ is closed and bounded, $\pts H$ is closed and bounded, hence compact, and $H=\pi(\pts H)$.

(b) A simple continuous arc of elements is, by (a), a cut-point arc in the metrizable space $G$, hence
an arc (Proposition~\ref{prop:cutarc}(b)). Conversely an arc $H$ in $G$ is a cut-point arc, compact and
so closed; its point set is compact, because a compact set of elements has compact point set (cover
$\pts H$ by open sets; each element of $H$ is covered by finitely many, whose union contains a tube
about that element, Proposition~\ref{prop:closed-map}; finitely many tubes cover $H$), hence bounded.

(c) If $J$ is a simple closed curve of elements and $p\ne q$ are in $J$, Theorem~\ref{thm:circle-char}
writes $J$ as the union of two cut-point arcs from $p$ to $q$ meeting only in $p,q$; they are arcs by
(b). Conversely a simple closed curve in $G$ is compact (hence closed and, as in (b), bounded),
connected, has two points, and is cut by every pair of points (Proposition~\ref{prop:scc-cut}). The
last equivalence is Proposition~\ref{prop:scc-cut}.
\end{proof}

\formal{\lean{HPC.Moore.isCompact_of_isClosed_bounded}, \lean{HPC.Moore.isArcElt_iff_tarc},
\lean{HPC.Moore.isSCCElt_iff_tscc}.}

From now on we use ``arc'' and ``simple closed curve'' for sets of elements as well.

\section{Theorems 5--8 and 10}

\begin{theorem}[Moore, Theorems 5--8 and 10]\label{thm:m25-5-10}
\leavevmode
\begin{enumerate}[label=(\alph*)]
\item \emph{(Theorem~5.)} An arc of elements minus one of its ends is connected.
\item \emph{(Theorem~6.)} A closed connected subset of an arc of elements with at least two elements
is an arc of elements.
\item \emph{(Theorem~7.)} Let $pq$ be an arc of elements, $x\in pq$ with $x\ne p,q$, and $H\subseteq
pq$. Then $x$ is a limit element of $H$ if and only if every arc of elements $L\subseteq pq$ containing
$x$, $x$ not an end of $L$, contains an element of $H$ other than $x$.
\item \emph{(Theorem~8.)} If $K$ is a closed, connected, bounded set of elements and $H\subsetneq K$ is
connected, some element $x\in K\setminus H$ has $K\setminus\set x$ connected.
\item \emph{(Theorem~10.)} If $J$ is a simple closed curve of elements and $p\ne q$ are in $J$, then
$J=a\cup b$ with $a,b$ arcs of elements from $p$ to $q$ and $a\cap b=\set{p,q}$.
\end{enumerate}
\end{theorem}

\begin{proof}
By Proposition~\ref{prop:arc-elt-top} we may work with topological arcs. (a) Lemma~\ref{lem:subarc}(c).
(b) Proposition~\ref{prop:cutarc}(c): such a set is $f([s,t])$, $s<t$, for a parametrization $f$.
(c) The sub-arcs of $pq$ having $x=f(u)$ as an interior point are the sets $f([s,t])$ with $s<u<t$
(Proposition~\ref{prop:cutarc}(c)), and these form a neighbourhood base of $x$ in $pq$; and $x$ is a
limit element of $H$ iff it is an accumulation point of $H$ (Proposition~\ref{prop:moore-notions}).
(d) Proposition~\ref{prop:cutarc}(d), since $K$ is compact. (e) Proposition~\ref{prop:arc-elt-top}(c).
\end{proof}

\formal{\lean{HPC.Moore.m25_theorem5}, \lean{HPC.Moore.m25_theorem6}, \lean{HPC.Moore.m25_theorem7},
\lean{HPC.Moore.m25_theorem8}, \lean{HPC.Moore.m25_theorem10}.}

\begin{remark}
In Theorem~7 Moore prints ``then $p$ is a limit element of the set $H$''; with $p$ an end of the arc
the statement is false (\lean{HPC.Moore.not_m25_theorem7_printed}), and $x$ is meant. Moore's proof of
Theorem~8 appeals to Theorem~5 for the claim that a certain set ``would be a simple continuous arc'';
the formal proof replaces this by the non-cut point theorem of Volume~3.
\end{remark}

\section{Theorem 9: arcs of elements do not separate}

\begin{theorem}[Moore, Theorem~9]\label{thm:m25-9}
If $K$ is an arc of elements, then $G\setminus K$ is connected.
\end{theorem}

\begin{proof}
$K$ is bounded, so by Corollary~\ref{cor:quot-tools}(c) the complement $G\setminus K$ has exactly one
unbounded component. It is enough to show that \emph{no component of $G\setminus K$ is bounded}: then
all components coincide with the unbounded one, and $G\setminus K$ (nonempty, since $\pts K$ is
bounded) is connected.

Suppose the component $D$ of some $x$ in $G\setminus K$ is bounded. Then $\fr D\subseteq K$ is bounded;
let $E$ be the unbounded component of $G\setminus\fr D$ and $K'=\fr E$. By Theorem~\ref{thm:m25-2}
(applied with $M=K$), $K'$ is closed and connected and $K'\setminus\set r$ is connected for every $r\in
K'$. As in the proof of Theorem~\ref{thm:m25-2}, $D\cap E=\emptyset$.

$K'$ has at least two elements. It is nonempty: otherwise $E$ would be open and closed in the
connected space $G$, so $E=G\supseteq D$. It is not a single element $p$: otherwise $E$ would be open
and closed in the connected set $G\setminus\set p$ (Corollary~\ref{cor:quot-tools}(d)), so
$E=G\setminus\set p\supseteq D$ (note $p\in\fr D$, so $p\notin D$).

Now $K'\subseteq\fr D\subseteq K$ is a closed connected subset of the arc $K$ with two elements, so it is
an arc by Theorem~\ref{thm:m25-5-10}(b). An arc has an interior point (Lemma~\ref{lem:subarc}(c)), and
removing it disconnects the arc; but $K'\setminus\set r$ is connected for every $r\in K'$. This
contradiction proves the theorem.
\end{proof}

\formal{\lean{HPC.Moore.m25_theorem9}.}

Moore gives no proof of Theorem~9, only the remark that it ``may be established by an argument
largely similar to that used in the proof of Lemma~8'' of an earlier paper, ``with the help of
Theorems~2 and~6''. The argument above uses exactly Theorems~2 and~6.

\begin{corollary}\label{cor:pts-nonsep}
If $H\subseteq G$ and $G\setminus H$ is connected, then $\pts H$ does not separate the plane. In
particular the point set of an arc of elements does not separate the plane.
\end{corollary}

\begin{proof}
$\pts{(G\setminus H)}=\R^2\setminus\pts H$ is connected by Proposition~\ref{prop:pts-conn}(a).
\end{proof}

\formal{\lean{HPC.Moore.not_separatesPlane_ptsOf}.}

\section{Theorem 12}

\begin{theorem}[Moore, Theorem~12]\label{thm:m25-12}
Let $D_1,D_2$ be bounded connected open sets of elements such that $\fr D_1$ is connected and $\fr
D_2\subseteq D_1$. Then $D_2\subseteq D_1$.
\end{theorem}

\begin{proof}
$D_2$ is nonempty (it is connected) and not all of $G$ (it is bounded), so $\fr D_2\ne\emptyset$ ($G$ is
connected). Since $D_1$ is open, $\fr D_1\cap D_1=\emptyset$, hence $\fr D_1\cap\fr D_2=\emptyset$. The
connected set $\fr D_1$ misses $\fr D_2$, so by Lemma~\ref{lem:conn-frontier} it lies either in $D_2$ or
in $G\setminus\cl{D_2}$.

If $\fr D_1\subseteq D_2$: then $\cl{D_1\cup D_2}\subseteq D_1\cup\fr D_1\cup D_2\cup\fr D_2\subseteq D_1\cup
D_2$, so $D_1\cup D_2$ is a nonempty open and closed subset of $G$, hence equal to $G$; but it is
bounded, and $G$ is not.

So $\fr D_1\cap D_2=\emptyset$. The connected set $D_2$ misses $\fr D_1$ and meets $D_1$ (near a point
of $\fr D_2\subseteq D_1$, which is a limit of points of $D_2$, and $D_1$ is open). By
Lemma~\ref{lem:conn-frontier}, $D_2\subseteq D_1$.
\end{proof}

\formal{\lean{HPC.Moore.m25_theorem12}.}

\section{The decomposition into points}\label{sec:points}

The decomposition of the plane into its single points satisfies the standing hypothesis, and the
quotient map is a homeomorphism $\R^2\to G$. So every theorem of this and the next chapters, applied to
this decomposition, is a theorem about the plane; and every counterexample in the plane to a printed
statement is a counterexample to that statement in Moore's setting. This is how the formal development
refutes the printed forms of Theorems~2, 7, 10 and~11, and how it obtains the Jordan curve theorem in
Chapter~\ref{ch:jordan}.

\formal{\lean{HPC.Moore.mooreDecomp_ptSetoid}, \lean{HPC.Moore.ptHomeo}.}

\section*{Exercises}
\addcontentsline{toc}{section}{Exercises}

\begin{exercise}
Deduce from Theorem~\ref{thm:m25-9} that no arc in the plane separates the plane. (This is a
classical theorem; compare the proof with the one by logarithms: an arc is contractible, so every
nonvanishing function on it has a logarithm, and \vfour{Theorem~2.5} applies.)
\end{exercise}

\begin{exercise}
Show that in Theorem~\ref{thm:m25-12} the hypothesis that $\fr D_1$ is connected cannot be omitted:
in the decomposition into points take $D_1=\set{z:1<\nrm z<3}$ and $D_2=\set{z:\nrm z<2}$.
\end{exercise}

\begin{exercise}
Prove Theorem~\ref{thm:m25-5-10}(c) directly from the definitions, without using the identification
of arcs of elements with topological arcs.
\end{exercise}

%======================================================================
\chapter{Mullikin's theorem and the Jordan curve theorem}\label{ch:jordan}
%======================================================================

This chapter proves Moore's Theorem~11: \emph{a simple closed curve of elements has exactly two
complementary domains, one bounded and one unbounded, and is the boundary of each}. Applied to the
decomposition of the plane into points, it is the Jordan curve theorem. The plane input is a theorem
of A.~M.~Mullikin (1922) on the union of two non-separating continua, which we prove with
logarithms, as in Volume~4, instead of Mullikin's polygons.

Points of the plane are identified with complex numbers (\vfour{Section~0.4}). Recall from Volume~4:
two logarithms of the same function on a connected set differ by a constant in $2\pi i\Z$
(\vfour{Lemma~1.3}); if $y$ is in the component of $x$ in the complement of a closed set $K$, then
$(z-x)/(z-y)$ has a logarithm on $K$ (\vfour{Theorem~2.4}); if that component is bounded, then for
every integer $a\ne0$ and every $u$ continuous and nowhere zero on its closure, $(z-x)^au(z)$ has no
logarithm on $K$ (\vfour{Theorem~2.5}); and if $K$ is compact with connected complement, every
continuous nowhere zero function on $K$ has a logarithm on $K$ (Borsuk, \vfour{Theorem~2.10}).

\section{Mullikin's theorem}

\begin{theorem}[Mullikin 1922, Theorem~5, corrected]\label{thm:mullikin}
Let $M_1,M_2\subseteq\R^2$ be continua neither of which separates the plane, and suppose
$M_1\cap M_2=K_1\cup K_2$ with $K_1,K_2$ disjoint, closed and connected. Then $\R^2\setminus(M_1\cup
M_2)$ has exactly two components: it is the union of two disjoint connected open sets.
\end{theorem}

The theorem has two halves, proved separately.

\begin{proposition}[At least two components]\label{prop:mullikin-two}
Let $M_1,M_2$ be continua with $M_1\cap M_2=K_1\cup K_2$, where $K_1,K_2$ are disjoint, nonempty and
closed. Then $M_1\cup M_2$ separates the plane.
\end{proposition}

\begin{proof}
Suppose $\R^2\setminus(M_1\cup M_2)$ is connected. By Urysohn's lemma there is a continuous
$\varphi\colon\R^2\to[0,1]$ with $\varphi=0$ on $K_1$ and $\varphi=1$ on $K_2$. Define $f$ on $M_1\cup
M_2$ by $f=1$ on $M_1$ and $f=e^{2\pi i\varphi}$ on $M_2$; on $M_1\cap M_2$ both formulas give $1$, and
$M_1,M_2$ are closed, so $f$ is continuous and nowhere zero. By Borsuk's theorem $f$ has a logarithm
$g$ on the compact set $M_1\cup M_2$. Take $k_1\in K_1$, $k_2\in K_2$. On the connected set $M_1$, $g$ and
the constant $0$ are logarithms of $1$, so $g(k_1)=g(k_2)$. On the connected set $M_2$, $g$ and
$2\pi i\varphi$ are logarithms of $e^{2\pi i\varphi}$, so $g(k_2)-g(k_1)=2\pi i(\varphi(k_2)-\varphi(k_1))=2\pi
i$. Contradiction.
\end{proof}

\formal{\lean{HPC.Moore.not_isConnected_compl_union}.}

\begin{proposition}[At most two components; Eilenberg]\label{prop:mullikin-three}
Let $M_1,M_2$ be compact sets with connected complements, and $M_1\cap M_2=K_1\cup K_2$ with $K_1,K_2$
connected. Then $\R^2\setminus(M_1\cup M_2)$ has at most two components.
\end{proposition}

\begin{proof}
Put $M=M_1\cup M_2$, a compact set. Suppose $M$ had three components $C_p,C_q,C_y$ of its complement,
containing $p,q,y$. At most one component is unbounded (\vfour{Lemma~2.3}); renaming, we may assume
$C_p$ and $C_q$ are bounded.

For $x\in\set{p,q}$ and $i\in\set{1,2}$, the point $x$ and $y$ lie off $M_i$, whose complement is
connected; so $h_x(z)=(z-x)/(z-y)$ has a logarithm $g_{ix}$ on $M_i$ (\vfour{Theorem~2.4}). On each
connected set $K_j$ the two logarithms $g_{1x},g_{2x}$ of $h_x$ differ by a constant: $g_{1x}-g_{2x}=2\pi
i\,n_{jx}$ on $K_j$ with $n_{jx}\in\Z$ (if $K_j$ is empty there is nothing to say; take $n_{jx}=0$).

Choose integers $a,b$, not both zero, with $a(n_{1p}-n_{2p})+b(n_{1q}-n_{2q})=0$ (if both differences
vanish take $a=1,b=0$; otherwise $a=n_{1q}-n_{2q}$, $b=-(n_{1p}-n_{2p})$). Define $G$ on $M$ by
\[
G=a\,g_{1p}+b\,g_{1q}\ \text{on }M_1,\qquad G=a\,g_{2p}+b\,g_{2q}+2\pi i(a\,n_{1p}+b\,n_{1q})\ \text{on }M_2\setminus M_1.
\]
On $K_1$ the two expressions agree by definition of the $n_{1x}$, and on $K_2$ they agree by the choice
of $a,b$; so the second expression, extended to all of $M_2$, agrees with the first on $M_1\cap M_2$,
and $G$ is continuous on $M$ ($M_1,M_2$ closed). Moreover $e^{G}=h_p^a\,h_q^b$ on $M$.

Say $a\neq0$ (the case $b\ne0$ is symmetric, exchanging $p$ and $q$). Then
$h_p^ah_q^b=(z-p)^a\,u(z)$ with $u(z)=(z-y)^{-a}h_q(z)^b$. The closure of $C_p$ lies in $C_p\cup M$,
which contains neither $y$ nor $q$, so $u$ is continuous and nowhere zero on $\cl{C_p}$. By
\vfour{Theorem~2.5}, $(z-p)^au(z)$ has no logarithm on $M$; but $G$ is one. Contradiction.
\end{proof}

\formal{\lean{HPC.Moore.exists_jump}, \lean{HPC.Moore.no_log_pow_of_bounded_component},
\lean{HPC.Moore.not_two_bounded_components}, \lean{HPC.Moore.compl_union_two_components}.}

\begin{proof}[Proof of Theorem~\ref{thm:mullikin}]
$K_1,K_2$ are nonempty (connected sets are nonempty), so the complement of $M_1\cup M_2$ has at least
two components (Proposition~\ref{prop:mullikin-two}) and at most two
(Proposition~\ref{prop:mullikin-three}). The components are open, since $M_1\cup M_2$ is closed.
\end{proof}

\formal{\lean{HPC.Moore.mullikin_theorem5}.}

\begin{remark}\label{rem:mullikin-printed}
\leavevmode
\begin{enumerate}
\item Mullikin states the theorem with $K_1,K_2$ merely mutually exclusive connected sets. This is
false: for $M_1=M_2$ a segment, $M_1\cap M_2$ is the union of two disjoint connected (not closed)
pieces, and the complement is connected (\lean{HPC.Moore.not_mullikin_theorem5_printed}). Moore uses
the theorem with $K_1,K_2$ two elements, which are closed.
\item Mullikin's Theorem~4 says that the union of two plane continua, neither separating the plane,
separates the plane if and only if their intersection is not connected. With her convention that the
empty set is connected, the case $M_1\cap M_2=\emptyset$ says that the union of two disjoint
non-separating continua does not separate the plane; this is true (Janiszewski's theorem with empty
intersection, \lean{HPC.Moore.mullikin_theorem4_disjoint}). With the modern convention (the empty set
is not connected) the same sentence would be false for two points
(\lean{HPC.Moore.not_mullikin_theorem4_disjoint_mathlib}).
\end{enumerate}
\end{remark}

\section{Theorem 11}

We return to the standing hypothesis of Chapter~\ref{ch:usc}.

\begin{lemma}\label{lem:frontier-scc}
Let $J$ be a simple closed curve of elements and $I\subseteq G$ a nonempty open set with $I\cap
J=\emptyset$ and $\fr I\subseteq J$. Then $\fr I=J$.
\end{lemma}

\begin{proof}
Put $B=\fr I$, a closed subset of $J$, and suppose $x\in J\setminus B$. Since $J$ is homeomorphic to
a circle (Proposition~\ref{prop:arc-elt-top}(c)) and $B$ is closed, there are elements $p_1\ne q_1$ of
$J$ such that one of the two arcs of $J$ from $p_1$ to $q_1$ (Theorem~\ref{thm:m25-5-10}(e)), call it
$C$, contains $x$ as an interior point and misses $B$. Let $A$ be the other arc; then $B\subseteq A$ and
$x\notin A$. By Theorem~\ref{thm:m25-9}, $G\setminus A$ is connected. Now $I\subseteq G\setminus A$, $I$
is open, and its closure $I\cup B$ meets $G\setminus A$ only in $I$; so $I$ is a nonempty open and
closed subset of $G\setminus A$, hence $I=G\setminus A$. But $x\in G\setminus A$ and $x\in J$, while
$I\cap J=\emptyset$.
\end{proof}

\formal{\lean{HPC.Moore.frontier_eq_of_scc}.}

\begin{theorem}[Moore, Theorem~11]\label{thm:m25-11}
Let $J$ be a simple closed curve of elements. Then $G\setminus J=I\cup E$ with $I,E$ disjoint,
open and connected, $I$ bounded, $E$ unbounded, and $\fr I=\fr E=J$.
\end{theorem}

\begin{proof}
Choose distinct $p,q\in J$ and, by Theorem~\ref{thm:m25-5-10}(e), arcs $a,b$ of elements from $p$
to $q$ with $a\cup b=J$ and $a\cap b=\set{p,q}$. By Theorem~\ref{thm:m25-9} and
Corollary~\ref{cor:pts-nonsep}, neither $\pts a$ nor $\pts b$ separates the plane. They are continua
(compact by Proposition~\ref{prop:arc-elt-top}, connected by Proposition~\ref{prop:pts-conn}(a)), and
$\pts a\cap\pts b=\pts{\set{p,q}}=p\cup q$ is the union of the two disjoint continua $p$ and $q$
(elements are compact and connected). By Mullikin's theorem, $\R^2\setminus\pts J=U\cup V$ with
$U,V$ disjoint, connected and open.

Each element $g\notin J$ is connected and lies in $U\cup V$, so it lies in $U$ or in $V$; hence $U$
and $V$ are unions of elements. Put $I'=\pi(U)$, $E'=\pi(V)$: they are open (their point sets are $U,V$),
connected, disjoint, and $I'\cup E'=G\setminus J$. As $J$ is bounded, exactly one component of
$G\setminus J$ is unbounded (Corollary~\ref{cor:quot-tools}(c)), and not both are bounded (else $G=I'\cup
E'\cup J$ would be bounded). Let $I$ be the bounded one and $E$ the unbounded one.

$I$ is open, and it is closed in $G\setminus J$ (its complement there is the open set $E$); so $\cl
I\subseteq I\cup J$ and $\fr I\subseteq J$. Lemma~\ref{lem:frontier-scc} gives $\fr I=J$; likewise $\fr
E=J$.
\end{proof}

\formal{\lean{HPC.Moore.m25_theorem11}, \lean{HPC.Moore.sccInterior_spec}.}

\begin{remark}
Moore's proof says that $J$ is the boundary of each domain because ``it easily follows, with the help
of Theorems~10 and~6, that $B$ is a simple continuous arc''. Theorem~6 needs $B=\fr I$ to be
connected, which is not known at that point. Lemma~\ref{lem:frontier-scc} replaces this step; it uses
Theorems~9 and~10 only.
\end{remark}

The bounded domain $I$ is the \emph{interior} of $J$ and $E$ its \emph{exterior}.

\section{The Jordan curve theorem}

\begin{theorem}[Jordan curve theorem]\label{thm:jct}
Let $J\subseteq\R^2$ be homeomorphic to a circle. Then $\R^2\setminus J=I\cup E$ with $I,E$ disjoint,
open and connected, $I$ bounded and $E$ unbounded, and $\fr I=\fr E=J$.
\end{theorem}

\begin{proof}
Apply Theorem~\ref{thm:m25-11} to the decomposition of the plane into points
(Section~\ref{sec:points}): there $G$ is the plane itself, and by
Proposition~\ref{prop:arc-elt-top}(c) the simple closed curves of elements are exactly the subsets
of the plane homeomorphic to a circle.
\end{proof}

\formal{\lean{HPC.Jones.jordan_curve_theorem}.}

The proof uses: logarithms and Borsuk's theorem (Volume~4), Brouwer's theorem and Janiszewski's
theorem (Volume~4), the characterizations of arcs and circles (Chapter~\ref{ch:arcs}), Theorems~2,
6, 9 and~10 (Chapters~\ref{ch:usc} and~\ref{ch:elements}) and Mullikin's theorem. It does not use
Moore's characterization of the plane (Chapters~\ref{ch:axioms}--\ref{ch:plane}). In the decomposition
into points, Theorem~2 and Theorem~9 are statements about plane continua and plane arcs, and the
whole argument is a proof of the Jordan curve theorem by Janiszewski's and Borsuk's theorems, in the
spirit of Eilenberg's.

\section*{Exercises}
\addcontentsline{toc}{section}{Exercises}

\begin{exercise}
Deduce from Proposition~\ref{prop:mullikin-two} that a simple closed curve in the plane separates the
plane, using only that an arc does not separate the plane.
\end{exercise}

\begin{exercise}
Generalize Proposition~\ref{prop:mullikin-three}: if $M_1\cap M_2$ has $n$ components, the complement of
$M_1\cup M_2$ has at most $n$ components.
\end{exercise}

\begin{exercise}
Show that the theta-curve (a circle together with a diameter) has exactly three complementary
domains, using Theorem~\ref{thm:jct} and Proposition~\ref{prop:mullikin-three}.
\end{exercise}

%======================================================================
\chapter{Moore's axioms and the bridge to topology}\label{ch:axioms}
%======================================================================

To prove that the space of elements of a Moore decomposition is a plane, Moore uses his own
axiomatic characterization of the plane (1916, 1919): a set of ``points'' with a class of subsets
called ``regions'' satisfying eight axioms is a plane. Chapters~\ref{ch:axioms}--\ref{ch:plane}
prove this characterization; Chapter~\ref{ch:quotient} verifies the axioms for the space of elements.

This chapter states the axioms (Moore 1916, \S2) and shows that Moore's notions agree with those of
the topology generated by the regions (the ``bridge''). Then the topology does the work: the space is
metrizable, closures of regions are compact, two points of a connected open set are joined by an arc
(Theorem~15, by Moore's nested chains), and Moore's arcs are the topological ones. We cite Moore's
paper \emph{On the foundations of plane analysis situs} (1916) as F.A., as Moore does.

\section{Points, regions and the axioms}\label{sec:axioms}

A \emph{Moore space} is a set $S$ together with a class of subsets of $S$, the \emph{regions}.

\begin{definition}[Moore's derived notions]\label{def:moore-notions}
Let $M\subseteq S$.
\begin{itemize}
\item $P$ is a \emph{limit point} of $M$ if every region containing $P$ contains a point of $M$
other than $P$.
\item The \emph{boundary} $\bd M$ is the set of points $X$ such that every region containing $X$
contains a point of $M$ and a point not in $M$. $M'=M\cup\bd M$, and the \emph{exterior} of $M$ is
$S\setminus M'$.
\item $M$ is \emph{connected} if, however $M$ is divided into two disjoint nonempty sets, one of them
contains a limit point of the other. $M$ is \emph{closed} if it contains all its limit points.
\item $M$ is \emph{bounded} if $M\subseteq(R_1\cup\dots\cup R_n)'$ for finitely many regions; $M$ has
the \emph{Heine--Borel property} if every cover of $M$ by regions has a finite subcover.
\item A \emph{domain} is a connected set each point of which lies in a region contained in the set.
\item An \emph{arc from $A$ to $B$} ($A\ne B$) is a bounded, closed, connected set containing $A$ and
$B$ with no connected proper subset containing both. A \emph{simple closed curve} is the union of two
arcs from $A$ to $B$ with only $A$ and $B$ in common.
\item A \emph{simple chain from $A$ to $B$} is a finite sequence of regions $R_0,\dots,R_{n-1}$ such
that $A$ lies in $R_0$ only, $B$ in $R_{n-1}$ only, and for $i<j$, $R_i\cap R_j\ne\emptyset$ if and only
if $j=i+1$.
\end{itemize}
\end{definition}

\formal{\lean{HPC.Moore.Space}, \lean{HPC.Moore.Space.IsLimitPt}, \lean{HPC.Moore.Space.bdry},
\lean{HPC.Moore.Space.cl}, \lean{HPC.Moore.Space.ext}, \lean{HPC.Moore.Space.IsConn},
\lean{HPC.Moore.Space.IsBoundedM}, \lean{HPC.Moore.Space.IsDomain}, \lean{HPC.Moore.Space.IsArc},
\lean{HPC.Moore.Space.IsSCC}, \lean{HPC.Moore.Space.IsSimpleChain}.}

In Moore's definition of ``connected'' the two parts must be nonempty (otherwise no set would be
connected); so read, the empty set is connected, i.e.\ Moore's ``connected'' is ``preconnected''.

\begin{definition}[Moore's axioms, F.A.\ pp.~132--133]\label{def:moore-axioms}
\leavevmode
\begin{description}
\item[Axiom 1.] There is a sequence of regions $K_1,K_2,\dots$ such that (1) every point lies in $K_n$
for arbitrarily large $n$; (2) if $P\ne\bar P$ are points of a region $R$, then for all large $n$,
$P\in K_n$ implies $K_n'\subseteq R\setminus\set{\bar P}$.
\item[Axiom 2.] Every region is connected.
\item[Axiom 3.] The exterior of every region is connected.
\item[Axiom 4.] For every region $R$, $R'$ has the Heine--Borel property.
\item[Axiom 5.] There is an infinite set of points with no limit point.
\item[Axiom 6.] If $R,\bar R$ are regions and $P\in\bar R$ is a boundary point of $R$, there are regions
$K,\bar K\subseteq\bar R$ with $P\in\bar K$, $K\subseteq R$, and every boundary point of $R$ in $\bar K$ a
boundary point of $K$.
\item[Axiom 7.] The same as Axiom~6 with ``$L$ in the exterior of $R$'' in place of ``$K\subseteq R$''.
\item[Axiom 8.] Every simple closed curve is the boundary of some region.
\end{description}
Axiom~$5_1$ (a footnote of F.A.) says that there are at least three points; it follows from Axiom~5.
\end{definition}

\formal{\lean{HPC.Moore.Space.Axiom1} \dots\ \lean{HPC.Moore.Space.Axiom8},
\leansub{HPC.Moore.Space.Axiom5}{1}, \leansub{HPC.Moore.axiom5}{1}\lean{_of_axiom5}.}

The model to keep in mind: $S=\R^2$, the regions the interiors of simple closed curves (or the open
squares). The \emph{fundamental regions} $K_n$ of Axiom~1 play the role of a countable basis of
shrinking sets.

\section{The bridge}

Let $\tau$ be the topology on $S$ generated by the regions (the smallest topology in which every region
is open).

\begin{proposition}[The bridge]\label{prop:bridge}
Assume Axiom~1.
\begin{enumerate}[label=(\alph*)]
\item The regions form a basis of $\tau$: if $P$ lies in two regions, some region about $P$ lies in
both (F.A.\ Theorem~6).
\item $P$ is a limit point of $M$ in Moore's sense if and only if it is an accumulation point of $M$ in
$\tau$. Consequently $M'$ is the closure of $M$, Moore's closed sets are the closed sets, $\bd M$ is the
frontier of $M$ minus the interior points, the exterior of a region is the complement of its closure,
Moore's connected sets are the preconnected sets, and the domains are the open connected sets.
\item The Heine--Borel property is compactness.
\item $\tau$ is $T_1$ and regular.
\item With Axioms~3 and $5_1$: every nonempty region is infinite (F.A.\ Theorem~4), the fundamental regions
form a countable basis, and $\tau$ is metrizable.
\item With Axiom~4: closures of regions are compact; a set is bounded if and only if it lies in a
compact set.
\end{enumerate}
\end{proposition}

\begin{proof}
(a) Let $P\in K\cap L$. If $K$ has a point $\bar P\ne P$, Axiom~1 gives $n$ with $P\in K_n$ and
$K_n'\subseteq K\setminus\set{\bar P}$; applying Axiom~1 again inside $L$ (if $L\ne\set P$) and taking $n$
large enough for both, $K_n\subseteq K\cap L$. If $K=\set P$ or $L=\set P$, that region itself will do.

(b) With (a), ``every region containing $P$'' means ``every basic neighbourhood of $P$''. The other
statements are rewritings; for example, $M$ is connected in Moore's sense iff in every partition into
two nonempty pieces one meets the closure of the other, i.e.\ iff $M$ is preconnected.

(c) Covers by regions refine arbitrary open covers, since regions form a basis.

(d) Points are closed: if $\bar P\ne P$, Axiom~1 (2) applied to a region containing both (if there is
none, any region about $P$ misses $\bar P$) gives a region about $P$ missing $\bar P$. Regularity:
given $P$ in a region $R$, if $R$ has a second point, Axiom~1 (2) gives $K_n\ni P$ with $K_n'\subseteq R$.

(e) Let $O\in R$. By Axiom~$5_1$ there are points $P_1\ne P_2$ different from $O$. By (d) there is a
region $R_1\ni P_1$ missing $O$ and $P_2$, and a region $K\ni P_1$ with $K'\subseteq R_1$. Then $O$ and $P_2$
lie in the exterior of $K$, which is connected (Axiom~3). Divide it into $\set O$ and the rest (which
contains $P_2$): one part contains a limit point of the other. A point is never a limit point of the
single point $O$ (by (d)), so $O$ is a limit point of the exterior of $K$. Hence every region about $O$, in
particular $R$, contains infinitely many points (a region about a limit point of $M$ contains infinitely
many points of $M$, by (d) and Axiom~1). So there are no isolated points; every region about $P$ has a
second point, Axiom~1 (2) applies, and the $K_n$ containing $P$ with large $n$ lie in any given region
about $P$: the $K_n$ form a countable basis. A regular second countable $T_1$ space is metrizable.

(f) Axiom~4 and (c). A bounded set lies in the compact set $(R_1\cup\dots\cup R_n)'=\cl{R_1}\cup\dots\cup
\cl{R_n}$; conversely a compact set is covered by finitely many regions, so it is bounded.
\end{proof}

\formal{\lean{HPC.Moore.Space.topology}, \lean{HPC.Moore.fa_theorem6},
\lean{HPC.Moore.isTopologicalBasis_region}, \lean{HPC.Moore.isLimitPt_iff},
\lean{HPC.Moore.cl_eq_closure}, \lean{HPC.Moore.isClosedM_iff}, \lean{HPC.Moore.isConn_iff},
\lean{HPC.Moore.isDomain_iff}, \lean{HPC.Moore.heineBorel_iff}, \lean{HPC.Moore.t1Space},
\lean{HPC.Moore.regularSpace}, \lean{HPC.Moore.fa_theorem4}, \lean{HPC.Moore.secondCountable},
\lean{HPC.Moore.metrizable}, \lean{HPC.Moore.isCompact_closure_region},
\lean{HPC.Moore.isBoundedM_of_isCompact}.}

What matters for the rest of the volume is the conclusion: \emph{under Axioms~1, 3, 4
and $5_1$, $S$ is a metrizable space in which the regions form a basis of open sets with compact
closures.}

\begin{remark}\label{rem:empty-region}
Moore's Theorem~4 (``every region contains infinitely many points'') and Theorem~19 (``every region has
a boundary point'') are printed for all regions, and his proofs begin ``there exists in $R$ at least one
point''. The axioms do not exclude $\emptyset$ as a region: adding $\emptyset$ to the regions changes none of
the derived notions and preserves Axioms~1--8. So the printed statements fail in any model with
$\emptyset$ added, and ``nonempty'' must be added (\lean{HPC.Moore.not_fa_theorem4_printed},
\lean{HPC.Moore.not_fa_theorem19_printed}). Similarly Theorem~14 needs every $M_n$ nonempty
(\lean{HPC.Moore.not_fa_theorem14_printed}).
\end{remark}

\section{Elementary consequences}

The following theorems of F.A.\ \S\S3--5 are, after the bridge, standard facts about a metrizable space
with a basis of connected open sets with compact closures. We list them with the topological reason.

\begin{theorem}[F.A.\ Theorems~1--3, 5, 7--9, 11--14, 17--19, 22]\label{thm:fa-elementary}
Assume Axioms~1--5.
\begin{enumerate}[label=(\alph*)]
\item (1) No point of a region is a boundary point of it. (2) Every region about a limit point of $M$
contains infinitely many points of $M$. (3) A point not in a finite set $M$ lies in a region missing
$M$. (7) A limit point of $M\cup N$ is a limit point of $M$ or of $N$. (8), (9) Limit points are limits
of sequences, and a convergent sequence has no other limit point.
\item (11), (12) Closed bounded sets are compact. (13) Every infinite bounded set has a limit point.
(14) A decreasing sequence of nonempty bounded sets $M_n\supseteq M_{n+1}'$ has nonempty closed
intersection.
\item (17), (18) The exterior of a region contains an infinite set without limit point and is unbounded.
(19) Every nonempty region has a boundary point. (22) The whole space is connected.
\end{enumerate}
\end{theorem}

\begin{proof}
(a) and (b) follow from the bridge (Proposition~\ref{prop:bridge}) and metrizability, and (14) is the
nested intersection property of compact sets.
(17) The infinite set $M$ of Axiom~5 meets the bounded set $R'$ in a finite set (by (13)), so $M\setminus
R'$ will do; (18) follows by (13).
(22) If $S=U\cup V$ with $U,V$ disjoint, closed and nonempty, take a region $R\subseteq U$ (a region about a
point of $U$, inside the open set $U$). Its exterior is connected (Axiom~3) and contains $V$, and so lies in
$V$; hence $U\subseteq R'$, which is compact (Axiom~4). Likewise $V$ is compact, so $S$ is compact, and
then every infinite set has a limit point, against Axiom~5.
(19) Let $P\in R$; by (22) $R$ is not closed and open unless $R=S$, which is excluded by (18); so $\fr R$ is
nonempty, and as $R$ is open, $\fr R=\bd R$.
\end{proof}

\formal{\lean{HPC.Moore.fa_theorem1}, \lean{HPC.Moore.fa_theorem2}, \lean{HPC.Moore.fa_theorem3},
\lean{HPC.Moore.fa_theorem5}, \lean{HPC.Moore.fa_theorem7}, \lean{HPC.Moore.fa_theorem8},
\lean{HPC.Moore.fa_theorem9}, \lean{HPC.Moore.fa_theorem11}, \lean{HPC.Moore.fa_theorem12},
\lean{HPC.Moore.fa_theorem13}, \lean{HPC.Moore.fa_theorem14}, \lean{HPC.Moore.fa_theorem17},
\lean{HPC.Moore.fa_theorem18}, \lean{HPC.Moore.fa_theorem19}, \lean{HPC.Moore.isConn_univ},
\lean{HPC.Moore.fa_theorem22}.}

Each formal statement carries only the axioms its proof uses; for example Theorems~11--14 use only
Axioms~1 and~4.

\section{Simple chains and Theorem 15}

\begin{theorem}[F.A.\ Theorem~10]\label{thm:fa10}
Let $M$ be connected, $A\ne B$ points of $M$, and $\mathcal G$ a set of regions covering $M$. There is a
simple chain from $A$ to $B$ all of whose links belong to $\mathcal G$.
\end{theorem}

\begin{proof}
Call a \emph{weak chain} from $A$ to $B$ a finite sequence of sets with $A$ in the first, $B$ in the last,
and consecutive members meeting. Let $S_A$ be the set of points $X$ of $M$ such that there is a weak chain
of members of $\mathcal G$ from $A$ to $X$. Then $S_A$ and $M\setminus S_A$ are both relatively open in $M$
(if $X\in M$ lies in $R\in\mathcal G$, every point of $M\cap R$ is in $S_A$ iff $X$ is: append $R$). So
neither contains a limit point of the other, and since $M$ is connected and $A\in S_A$, $S_A=M\ni B$.

Now take a weak chain from $A$ to $B$ of members of $\mathcal G$ of minimal length; it is a simple chain.
Indeed, if $A$ lay in a link $R_i$ with $i>0$, the chain $R_i,\dots,R_{n-1}$ would be shorter; similarly
for $B$; and if $R_i\cap R_j\ne\emptyset$ with $j>i+1$, then $R_0,\dots,R_i,R_j,\dots,R_{n-1}$ would be a
shorter weak chain. Note that the minimal chain is a \emph{sublist} of the original, with the order of the
links preserved.
\end{proof}

\formal{\lean{HPC.Moore.fa_theorem10}, \lean{HPC.Moore.exists_simpleChain_sublist}.}

\begin{theorem}[F.A.\ Theorem~15]\label{thm:fa15}
Assume Axioms~1--4 and~$5_1$. Two distinct points of a domain $M$ are joined by an arc lying in $M$.
\end{theorem}

\begin{proof}
By the bridge, $S$ is metrizable (fix a metric $d$), $M$ is open and connected, and the regions are
connected open sets with compact closures, with small regions about every point inside every open set.

\emph{The chains.} We construct simple chains $C_1,C_2,\dots$ of regions from $A$ to $B$ such that
\begin{enumerate}[label=(\arabic*)]
\item each link of $C_1$ has its closure in $M$; each link of $C_{n+1}$ has its closure in a single link
of $C_n$, its \emph{owner};
\item the owner map is monotone: if a link precedes another in $C_{n+1}$, the owner of the first is not
after the owner of the second;
\item every link of $C_n$ lies in a ball of radius $1/n$.
\end{enumerate}
$C_1$: apply Theorem~\ref{thm:fa10} to $M$ and the regions of diameter $<1$ with closure in $M$. Given
$C_n=(R_0,\dots,R_{k-1})$, choose points $p_i\in R_{i-1}\cap R_i$ ($0<i<k$), $p_0=A$, $p_k=B$. Inside the
connected open set $R_i$, Theorem~\ref{thm:fa10} gives a simple chain from $p_i$ to $p_{i+1}$ of regions of
radius $<1/(n+1)$ with closures in $R_i$ (if $p_i=p_{i+1}$, a single such region). Laying these chains end
to end gives a weak chain from $A$ to $B$, in which the links coming from $R_i$ come before those coming
from $R_{i+1}$; its shortest sub-chain (proof of Theorem~\ref{thm:fa10}) is a sublist, so (1)--(3) hold
with the owner of a link the region $R_i$ it came from.

\emph{The continuum.} Let $F_n$ be the union of the closures of the links of $C_n$. Each $F_n$ is compact
and connected (closures of links are connected and consecutive ones meet), $F_{n+1}\subseteq F_n$ by (1),
and $A,B\in F_n$. So $C=\bigcap_nF_n$ is a continuum containing $A$ and $B$, and $C\subseteq F_1\subseteq M$.

\emph{The order.} A point lies in at most two links of a simple chain, and if in two, they are
consecutive. For $X,Y\in C$ write $X\prec Y$ if at some stage $n$ every link of $C_n$ containing $X$
precedes every link containing $Y$. By (2) this persists to all later stages. If $X\ne Y$, then at a stage
with $2/n<d(X,Y)/2$ no link containing $X$ meets a link containing $Y$, so the (one or two, consecutive)
links containing $X$ all precede, or all follow, those containing $Y$: $X\prec Y$ or $Y\prec X$. Also $A\prec
P$ and $P\prec B$ for $P\ne A,B$ ($A$ lies in the first link only, $B$ in the last only).

\emph{Irreducibility.} Fix $P\in C\setminus\set{A,B}$ and let $S_A=\set{X\in C:X\prec P}$,
$S_B=\set{X\in C:P\prec X}$. These are disjoint and cover $C\setminus\set P$, and $A\in S_A$, $B\in S_B$. We
show $S_A$ is open in $C$ (likewise $S_B$). Let $X\in S_A$; choose $n$ so large that $X\prec P$ at stage $n$
and no two meeting links of $C_n$ contain $X$ and $P$ respectively. Let $L$ be a link of $C_n$ containing $X$
and $Y\in C\cap L$. The links containing $Y$ are $L$ or adjacent to $L$; the links containing $P$ follow $L$
and are not adjacent to it. So every link containing $Y$ precedes every link containing $P$, $Y\prec P$, and
$C\cap L\subseteq S_A$ is a neighbourhood of $X$ in $C$. Hence $C\setminus\set P$ is the union of two disjoint
relatively open sets, one containing $A$ and the other $B$, and every connected subset of $C$ containing $A$
and $B$ contains $P$. So $C$ is irreducible between $A$ and $B$, and it is an arc by
Corollary~\ref{cor:arc-char}.
\end{proof}

\formal{\lean{HPC.Moore.exists_chains}, \lean{HPC.Moore.exists_irreducible},
\lean{HPC.Moore.fa_theorem15}.}

\begin{corollary}[F.A.\ Theorem~16]\label{cor:fa16}
Two points of a region are joined by an arc in the region; two points of the exterior of a region are
joined by an arc in the exterior.
\end{corollary}

\begin{proof}
A region is a domain (Axiom~2); its exterior is open (the closure is closed) and connected (Axiom~3).
\end{proof}

\formal{\lean{HPC.Moore.fa_theorem16}.}

\section{Moore's arcs and simple closed curves}

\begin{proposition}\label{prop:moore-arcs}
Assume Axioms~1, 3, 4 and $5_1$. A set is an arc from $A$ to $B$ in Moore's sense if and only if it is an
arc from $A$ to $B$ in the topology $\tau$ (Definition~\ref{def:arc}); and it is a simple closed curve in
Moore's sense if and only if it is one in $\tau$.
\end{proposition}

\begin{proof}
Bounded and closed means compact (Proposition~\ref{prop:bridge}), connected means connected, and
Corollary~\ref{cor:arc-char} applies in the metrizable space $(S,\tau)$.
\end{proof}

\formal{\lean{HPC.Moore.isArc_iff}, \lean{HPC.Moore.isSCC_iff_tscc}.}

From now on all notions are topological: ``region'' means a member of the given basis, ``arc'' and
``simple closed curve'' have their topological meaning, and $\cl M$, $\fr M$ are closure and frontier.
Moore's unnumbered facts about arcs (F.A.\ pp.~139--140: splitting an arc at a point, first and last points
on a closed set, the two arcs of a simple closed curve) are the results of Chapter~\ref{ch:arcs}
(\lean{HPC.Moore.fa_sec5_arc_split}, \lean{HPC.Moore.fa_sec5_first_point}, \lean{HPC.Moore.fa_scc_two_arcs}).

\section*{Exercises}
\addcontentsline{toc}{section}{Exercises}

\begin{exercise}
Check that $\R^2$ with the open squares as regions satisfies Axioms~1--5. (Axioms~6--8 for this choice of
regions are not obvious: Axiom~8 says that every simple closed curve bounds a square!) Show that with the
interiors of simple closed curves as regions all eight axioms hold, using the Jordan curve theorem and the
Schoenflies theorem.
\end{exercise}

\begin{exercise}
Show that the sphere $S^2$ with the interiors of simple closed curves (either complementary domain) as
regions satisfies Axioms~1--4 and 6--8 but not Axiom~5. Which step of the proof of Theorem~\ref{thm:fa-elementary}(22)
uses Axiom~5?
\end{exercise}

\begin{exercise}
Show that the construction in the proof of Theorem~\ref{thm:fa15} works in any connected, locally connected,
locally compact metric space, and deduce that such a space is arcwise connected.
\end{exercise}

%======================================================================
\chapter{Simple closed curves and their interiors}\label{ch:curves}
%======================================================================

This chapter follows F.A.\ \S\S6--7 (Theorems~20--43). It develops, from Moore's axioms, the plane
topology of simple closed curves: every simple closed curve has exactly one interior; a cross-cut
divides an interior into two interiors (the $\theta$-curve theorems); every point of a curve is
accessible from both sides; there are closed curves about every point inside every region; and two
interiors that meet can be merged into one (Theorems~41 and~42). These are the tools for the double
rulings of Chapter~\ref{ch:rulings}.

Throughout, $S$ is a space satisfying Axioms~1--8, with the topology of the bridge
(Proposition~\ref{prop:bridge}): $S$ is metrizable, the regions form a basis of connected open sets
with compact closures, and Moore's arcs and simple closed curves are the topological ones. ``Closed
curve'' means simple closed curve. For an arc we write $AXB$ (an arc from $A$ to $B$ through $X$),
$\oarc{AXB}$ for the open arc, and $AXBCA$ for the closed curve formed by two arcs $AXB$ and $ACB$ with
only $A,B$ in common.

Many arguments of F.A.\ are long case analyses with figures. We give complete proofs of the results
that are used later and whose proofs are short, and proofs in outline of the others, indicating the
construction and the key step; the formal proofs are complete.

\section{The interior of a closed curve}

\begin{theorem}[F.A.\ Theorems~20 and~21]\label{thm:fa20-21}
\leavevmode
\begin{enumerate}[label=(\alph*)]
\item \emph{(Theorem~20.)} Every boundary point of a region $R$ is a limit point of its exterior
$S\setminus\cl R$. Hence $\fr(S\setminus\cl R)=\fr R$.
\item \emph{(Theorem~21.)} If $K,R$ are regions and $\fr R\subseteq\cl K$, then $R\subseteq K$.
\end{enumerate}
\end{theorem}

\begin{proof}
(a) Let $P\in\fr R$ and let $\bar R$ be a region about $P$. Axiom~7 gives regions $L,\bar L\subseteq\bar
R$ with $L$ in the exterior of $R$, $P\in\bar L$, and every boundary point of $R$ in $\bar L$ a boundary
point of $L$. So $P\in\fr L$, and the region $\bar L\ni P$ contains a point of $L$, which lies in the
exterior of $R$ and in $\bar R$. As $\bar R$ was arbitrary, $P$ is a limit point of the exterior.

(b) The exterior $E_K=S\setminus\cl K$ is connected (Axiom~3) and does not lie in the compact set $\cl R$
(it is unbounded, Theorem~\ref{thm:fa-elementary}(18)). It misses $\fr R\subseteq\cl K$, so by
Lemma~\ref{lem:conn-frontier} it lies in $S\setminus\cl R$. Hence $R\subseteq\cl K$. If a point of $R$ were
on $\fr K$, it would be a limit point of $E_K$ by (a); but $R$ is open and misses $E_K$. So
$R\subseteq K$.
\end{proof}

\formal{\lean{HPC.Moore.fa_theorem20}, \lean{HPC.Moore.fa_theorem21}.}

\begin{corollary}[The interior of a closed curve]\label{cor:interior}
Every closed curve $J$ is the frontier of exactly one region, its \emph{interior} $I(J)$. The
\emph{exterior} $E(J)=S\setminus\cl{I(J)}$ is open and connected, $S\setminus J=I(J)\cup E(J)$ (disjoint), and
$\fr I(J)=\fr E(J)=J$. If $J'$ is a closed curve with $J'\subseteq\cl{I(J)}$, then $I(J')\subseteq I(J)$.
\end{corollary}

\begin{proof}
Existence is Axiom~8; if $R,R'$ are regions with frontier $J$, Theorem~\ref{thm:fa20-21}(b) gives
$R\subseteq R'$ and $R'\subseteq R$. $I(J)$ is connected (Axiom~2), $E(J)$ is connected (Axiom~3), and the
rest follows from Theorem~\ref{thm:fa20-21}.
\end{proof}

\formal{\lean{HPC.Moore.existsUnique_interior}, \lean{HPC.Moore.jordan}, \lean{HPC.Moore.Space.IsInteriorOf}.}

``Within $J$'' means in $I(J)$, ``without $J$'' means in $E(J)$. In F.A.\ the interior of a closed curve
is used from Theorem~23 on but never defined; Corollary~\ref{cor:interior} supplies the definition.

\paragraph{The topological setting.} From here on the axioms enter only through the following facts, and
the formal proofs of Theorems~24--34 are carried out in an arbitrary Hausdorff, locally connected space
with a family of open sets (``regions'') satisfying them:
\begin{enumerate}[label=(T\arabic*)]
\item two distinct points of a connected open set are joined by an arc in it (Theorem~\ref{thm:fa15});
\item if $K,R$ are regions and $\fr R\subseteq\cl K$ then $R\subseteq K$ (Theorem~\ref{thm:fa20-21});
\item every closed curve is the frontier of a region (Axiom~8);
\item regions and their exteriors are connected, closures of regions are compact, $S$ is not compact.
\end{enumerate}

\formal{\lean{HPC.Moore.Gen.ArcConn}, \lean{HPC.Moore.Gen.Th21Prop}, \lean{HPC.Moore.Gen.Ax8Prop},
\lean{HPC.Moore.Gen.TInt}.}

Two constructions recur. \emph{Accessibility:} if $W$ is a connected open set, $F\in W$, $Q\in\fr W$ and
$N$ a neighbourhood of $Q$, there is an arc $FH$ with $H\in\fr W\cap N$ and $FH\setminus\set H\subseteq W$ (join
$F$ by an arc in $W$ to a point of $W$ in a small connected neighbourhood of $Q$, continue in that
neighbourhood towards $Q$, and stop at the first point off $W$). \emph{Crossing:} two such arcs from a
common point $F$, glued, give an arc $HK$ through $W$ with ends on $\fr W$ near two prescribed boundary
points.

\formal{\lean{HPC.Moore.Gen.exists_access}, \lean{HPC.Moore.Gen.exists_cross}.}

\begin{theorem}[F.A.\ Theorem~23]\label{thm:fa23}
No nonempty region lies in an arc.
\end{theorem}

\begin{proof}[Proof (outline)]
Suppose the region $R$ lies in the arc $AB$. Take $O\in R$, $O\ne A,B$ (regions are infinite), a region $K$
about $O$ whose closure misses $A$ and $B$, and an arc $AXB$ in the exterior of $K$
(Corollary~\ref{cor:fa16}) from $A$ to $B$. Cutting at a last and a first point, one finds points $A_1$, $B_1$
on $AB$ on either side of $O$ and a sub-arc of $AXB$ from $A_1$ to $B_1$ meeting $AB$ only in $A_1,B_1$; with
the interval $A_1OB_1$ of $AB$ it forms a closed curve $J$. $O$ is a boundary point of $I(J)$, so a small region
about $O$, inside $R$ and missing the rest of $J$, meets $I(J)$; but it lies in $AB\subseteq J$, which misses
$I(J)$.
\end{proof}

\formal{\lean{HPC.Moore.fa_theorem23}.}

\section{Theta-curves: Theorems 24--27}

In this section $J=ACB\cup ADB$ is a closed curve with interior $R$ ($A,B$ separate $C,D$ on $J$), and
$AXB$ is an arc from $A$ to $B$ with $\oarc{AXB}$ disjoint from $J$. Put $R_1=I(AXBCA)$, $R_2=I(AXBDA)$.

\begin{theorem}[F.A.\ Theorem~24]\label{thm:fa24}
If $\oarc{AXB}\subseteq R$, then (1) $R_1\subseteq R$; (2) $\oarc{ADB}$ lies without $AXBCA$; (3)
$R_1\cap R_2=\emptyset$.
\end{theorem}

\begin{proof}
(1) $\fr R_1=AXBCA\subseteq\cl R$; apply (T2). (2) $\oarc{ADB}$ misses $\cl{R_1}=R_1\cup AXBCA$: it lies on $J=\fr R$, so it misses the open set $R\supseteq R_1$;
and it misses $AXBCA$, since $\oarc{AXB}\subseteq R$ misses $J$ and $ACB\cap ADB=\set{A,B}$. So it lies in $E(AXBCA)$.
(3) If $R_1\cap R_2\ne\emptyset$: $R_2\subseteq R$ by (1) for $R_2$, so $R_2$ misses $J\supseteq ACB$,
and it misses $\oarc{AXB}\subseteq\fr R_2$; so the connected set $R_2$ misses $\fr R_1=AXBCA$ and meets $R_1$,
whence $R_2\subseteq R_1$ (Lemma~\ref{lem:conn-frontier}). Then $D\in\fr R_2\subseteq\cl{R_1}$, against (2).
\end{proof}

\formal{\lean{HPC.Moore.fa_theorem24}, \lean{HPC.Moore.Gen.th24_sub}, \lean{HPC.Moore.Gen.th24_ext},
\lean{HPC.Moore.Gen.th24_disj}.}

\begin{theorem}[F.A.\ Theorem~25]\label{thm:fa25}
If $\oarc{AXB}\subseteq R$, then $R=\oarc{AXB}\cup R_1\cup R_2$.
\end{theorem}

\begin{proof}
``$\supseteq$'' is Theorem~\ref{thm:fa24}(1). The points of $R$ on $\fr R_1=AXBCA$ are those of $\oarc{AXB}$, so
$R\cap\cl{R_1}=R_1\cup\oarc{AXB}$, and likewise for $R_2$. Hence, with $Y=R\setminus(\oarc{AXB}\cup\cl{R_1}\cup\cl
{R_2})$, the set $R\setminus\oarc{AXB}$ is the union of the three disjoint open sets $R_1,R_2,Y$. Suppose
$Y\ne\emptyset$, and take $E\in Y$.

\emph{Step 1.} An arc in $R$ from $E$ to a point of $R_1$ (T1) must meet $\oarc{AXB}$; let $O$ be its first point on
the closed set $AXB$ ($A,B\notin R$). The sub-arc $EO$ minus $O$ is connected, lies in $R\setminus\oarc{AXB}$ and
contains $E$, so it lies in $Y$.

\emph{Step 2.} $E$ and the points of $\oarc{ADB}$ lie in the connected open set $E(AXBCA)$ (by
Theorem~\ref{thm:fa24}(2), and $E\notin\cl{R_1}$). An arc in it from $E$ to a point of $\oarc{ADB}$ has a first
point $L$ on $J$. The part before $L$ is connected, misses $J=\fr R$ and starts in $R$, so lies in $R$; it misses
$\cl{R_1}\supseteq\oarc{AXB}$; so it lies in $R_2\cup Y$, hence in $Y$. And $L\in J\setminus\cl{R_1}=\oarc{ADB}$.

\emph{Step 3.} Choose points $\bar H,H,O,K,\bar K$ in this order on $\oarc{AXB}$ and, by the crossing construction, arcs
$HP_1K$ with open part in $R_1$ and $\bar HP_2\bar K$ with open part in $R_2$ (the ends can be prescribed near
given points of $\oarc{AXB}$ because these points are on $\fr R_1$ and $\fr R_2$; one aims at points of $\oarc{AXB}$
on either side of $O$, which is why the order can be arranged). Then
$\bar J=\bar HP_2\bar K\cup\bar KK\cup KP_1H\cup H\bar H$ (with $\bar KK$, $H\bar H$ sub-arcs of $AXB$) is a closed curve
lying in $R$, so its interior $\bar R$ lies in $R$ by (T2).

\emph{Step 4.} $O\notin\bar J$, and $O$ is joined to $L$ by $OE\cup EL$, which misses $\bar J$ (its points other than $O,L$
lie in $Y$, and $L\notin R$). As $L\notin\cl{\bar R}\subseteq R$, $O$ lies in $E(\bar J)$ (Lemma~\ref{lem:conn-frontier}).
The open sub-arc of $AXB$ between $H$ and $K$ misses $\bar J$ and contains $O$, so it lies in $E(\bar J)$; the open
sub-arcs from $A$ to $\bar H$ and from $\bar K$ to $B$ miss $\bar J$ and have the ends $A,B\notin\cl{\bar R}$, so they lie
in $E(\bar J)$ too. Hence $\bar R$ misses $\oarc{AXB}$ and $\bar R\subseteq R_1\cup R_2\cup Y$. But $\bar R$ is connected and
meets $R_1$ (near $P_1\in\fr\bar R\cap R_1$) and $R_2$ (near $P_2$): a contradiction.
\end{proof}

\formal{\lean{HPC.Moore.fa_theorem25}, \lean{HPC.Moore.Gen.th25}.}

\begin{theorem}[F.A.\ Theorem~26]\label{thm:fa26}
If $ACB$, $ADB$, $AFB$ are three arcs, no two with a common point other than $A,B$, then two of the interiors of
$ACBFA$, $ACBDA$, $ADBFA$ meet.
\end{theorem}

\begin{proof}[Proof (outline)]
Suppose the three interiors $R=I(ACBFA)$, $R_1=I(ACBDA)$, $R_2=I(ADBFA)$ are pairwise disjoint. Then each open arc
lies without the curve formed by the other two (for $\oarc{ADB}$: were it within $ACBFA$, Theorem~\ref{thm:fa24} would put
$R_1$ inside $R$). The set $Y$ of points without all three curves is nonempty (the closures of the three interiors are
compact, $S$ is not). From a point $P\in Y$, an arc without $ACBFA$ reaches $\oarc{ADB}$ at a first point $O$, through $Y$.
Around $O$ one builds, by crossing arcs through $R_1$ and $R_2$, a closed curve $\bar J$ on $R_1\cup R_2\cup\oarc{ADB}$ as in
the proof of Theorem~\ref{thm:fa25}. If $I(\bar J)$ met $ACB\cup AFB$ it would contain it and then all of $Y$, which is
unbounded; if it met $Y$ it would meet $ACB$ through arcs in $Y$; so $I(\bar J)\subseteq R\cup R_1\cup R_2$, while it meets
$R_1$ and $R_2$ and is connected. Contradiction. (This simplifies Moore's three cases: his Case~I is reduced to Case~II and
his auxiliary curves are not needed.)
\end{proof}

\formal{\lean{HPC.Moore.fa_theorem26}, \lean{HPC.Moore.Gen.th26}.}

\begin{theorem}[F.A.\ Theorem~27]\label{thm:fa27}
If $\oarc{AXB}$ lies without $J$, then (1) $D$ is without $AXBCA$ or $C$ is without $AXBDA$; (2) if $D$ is without
$AXBCA$, then $C$ is within $AXBDA$ and $I(AXBDA)=\oarc{ACB}\cup R\cup R_1$.
\end{theorem}

\begin{proof}
(1) If $D$ is not without $AXBCA$, then $\oarc{ADB}$ (connected, missing $AXBCA$) is within $AXBCA$, and
Theorem~\ref{thm:fa24}(2), applied to the curve $AXBCA$ and the cross-cut $\oarc{ADB}$, puts $C$ without $AXBDA$.
(2) $R$ and $R_1$ are disjoint ($R_1$ is connected, misses $J$, and is not inside $R$, as its boundary point $X$ is
without $J$). If $C$ were also without $AXBDA$, the three interiors $R,R_1,R_2$ would be pairwise disjoint
(the cases are symmetric), against Theorem~\ref{thm:fa26}. So $C$ is within $AXBDA$; then $\oarc{ACB}$ is a
cross-cut of $AXBDA$ and Theorem~\ref{thm:fa25} gives $I(AXBDA)=\oarc{ACB}\cup R\cup R_1$.
\end{proof}

\formal{\lean{HPC.Moore.fa_theorem27}, \lean{HPC.Moore.Gen.th27}.}

\begin{corollary}[Splitting an interior]\label{cor:split-interior}
If $J$ is a closed curve, $AXB$ an arc with $A,B\in J$ and $\oarc{AXB}\subseteq I(J)$, and $C_1,C_2$ the two arcs of $J$
from $A$ to $B$, then $I(J)\setminus AXB$ is the disjoint union of $I(AXB\cup C_1)$ and $I(AXB\cup C_2)$.
\end{corollary}

\formal{\lean{HPC.Moore.split_interior}, \lean{HPC.Moore.split_interior_of}.}

\section{Arcs near a curve: Theorems 28--30}

\begin{theorem}[F.A.\ Theorem~28]\label{thm:fa28}
Let $O$ be a point of the closed curve $J$, $\bar R$ a region about $O$, and $M$ either $I(J)$ or $E(J)$. There is an
arc $AXB$ with $A,B\in J$, $\oarc{AXB}\subseteq M\cap\bar R$, such that of the two arcs of $J$ from $A$ to $B$ the one
containing $O$ (in its interior) lies in $\bar R$.
\end{theorem}

\begin{proof}[Proof (outline)]
Axiom~6 (for $M=I(J)$) or Axiom~7 (for $M=E(J)$) gives regions $L\subseteq M\cap\bar R$ and $\bar L\subseteq\bar R$ about $O$ such
that every point of $J$ in $\bar L$ is a boundary point of $L$. Take an arc $A_1OB_1$ of $J$ in $\bar L$, points $A_2,B_2$
with $A_1A_2OB_2B_1$ in order, small regions $R_A\ni A_2$, $R_B\ni B_2$ meeting $J$ only in $A_1O$ and $OB_1$, points
$C\in L\cap R_A$, $D\in L\cap R_B$, an arc from $C$ to $D$ in $L$ and arcs from $C$ to $A_2$ in $R_A$, from $D$ to $B_2$ in
$R_B$. The arc $AXB$ is extracted from their union; Moore's last step needs care, since the arc $CA_2$ may leave $M$
before reaching $J$: one takes for $A$ the \emph{first point of $CA_2$ outside $M$}, which lies on $\fr M=J$, hence on the
open arc $A_1O$; likewise $B$.
\end{proof}

\formal{\lean{HPC.Moore.fa_theorem28}.}

Theorem~29 (if $A,B$ separate $C,D$ on $J$, and $\oarc{AXB}$, $\oarc{CYD}$ are both within or both without $J$, they meet)
follows from Theorems~\ref{thm:fa24}, \ref{thm:fa25} and~\ref{thm:fa27}. In the case ``within'', if the arcs did not
meet, the connected set $\oarc{CYD}$ would lie in $R_1\cup R_2$ (Theorem~\ref{thm:fa25}), hence in one of them, say $R_1$; but its
limit point $D$ lies without $AXBCA$ (Theorem~\ref{thm:fa24}(2)), and symmetrically $C$ lies without $AXBDA$. The case ``without'' uses
Theorem~\ref{thm:fa27} in the same way. Theorem~30 (the interiors of $AXBDA$ and $CYDBC$ meet) is not proved by Moore; the formal proof
examines the four combinations of the alternatives of Theorem~\ref{thm:fa27}.

\formal{\lean{HPC.Moore.fa_theorem29}, \lean{HPC.Moore.fa_theorem30}.}

\section{The two sides of an arc and accessibility: Theorems 31--34}

\begin{theorem}[F.A.\ Theorem~31; the pocket construction]\label{thm:fa31}
Let $f\colon I\to S$ parametrize an arc, $0<b<1$, and $R$ a region about $f(b)$. There are closed curves $J_1=F_1\cup P_1$
and $J_2=F_2\cup P_2$, where $P_1\supseteq P_2$ are sub-arcs of the arc around $f(b)$ (in its interior) and $F_1,F_2$ are
arcs meeting the given arc only in their ends, such that $I(J_1)$ and $I(J_2)$ are disjoint, miss the arc, and together
with the open arc $\oarc{P_2}$ form a region about $f(b)$ inside $R$.
\end{theorem}

\begin{proof}[Proof (outline)]
A closed curve $J_0$ through $f(b)$ containing a sub-arc of the given arc is built as in Theorem~\ref{thm:fa23}. Theorem~28
inside $J_0$ at $f(b)$ gives an arc $F_1$ within $J_0$, closing a curve $J_1$; Theorem~28 again, for $J_1$ from outside,
gives $F_2$ without $J_1$, closing a curve $J_2$. By Theorem~\ref{thm:fa27}, the interior of the curve formed by $F_1$, $F_2$
and the two short pieces of the arc is $I(J_1)\cup I(J_2)\cup\oarc{P_2}$, and Theorem~\ref{thm:fa24}(3) makes $I(J_1)$ and
$I(J_2)$ disjoint.
\end{proof}

\formal{\lean{HPC.Moore.fa_theorem31}, \lean{HPC.Moore.th31_pocket}.}

\begin{corollary}[The two sides]\label{cor:sides}
If $f(b)$ is an interior point of an arc $\alpha$ and $W$ an open set about $f(b)$, there are an open set $N$ with $f(b)\in
N\subseteq W$ and two connected open sets $I_1,I_2\subseteq W\setminus\alpha$ with $N\setminus\alpha\subseteq I_1\cup I_2$, each of
which has every point of some open sub-arc of $\alpha$ about $f(b)$ as a limit point.
\end{corollary}

\formal{\lean{HPC.Moore.sides}.}

\begin{theorem}[Propagation and accessibility; F.A.\ Theorem~32]\label{thm:fa32}
Let $f$ parametrize an arc $\alpha=AB$ ($A=f(0)$, $B=f(1)$), $O$ an open set, and $Q$ the component of $O\setminus\alpha$
containing a point $x_0$.
\begin{enumerate}[label=(\alph*)]
\item \emph{(Propagation.)} If $T\subseteq(0,1)$ is an interval with $f(T)\subseteq O$ and some point of $f(T)$ is a limit point
of $Q$, then every point of $f(T)$ is a limit point of $Q$.
\item \emph{(Accessibility.)} If $f(u)\in O$ for all $u\in(t_0,1)$ and some such $f(u)$ is a limit point of $Q$, then $x_0$ is
joined to $B$ by an arc lying, except for $B$, in $Q$.
\end{enumerate}
\end{theorem}

\begin{proof}
(a) The set of $u\in T$ with $f(u)\in\cl Q$ is closed in $T$ and nonempty. It is open: if $f(u)\in\cl Q$, take $N,I_1,I_2$
as in Corollary~\ref{cor:sides} with $W\subseteq O$. $Q$ meets $N\setminus\alpha$, hence $I_1$ or $I_2$, say $I_1$; as $I_1$ is
connected and lies in $O\setminus\alpha$, $I_1\subseteq Q$, and every point of an open sub-arc about $f(u)$ is in $\cl{I_1}\subseteq\cl
Q$. As $T$ is connected, the set is all of $T$.

(b) (Outline.) By (a), $Q$ has every point of $f((t_0,1))$ as a limit point. Choose open sets $O_k\subseteq O$, $O_k$ within
distance $1/k$ of $B$ and containing $f((t_k,1))$ for suitable $t_k\to1$; by (a) the component $Q_k$ of $O_k\setminus\alpha$ adjacent to
$f((t_k,1))$ lies in $Q$ and, by connectedness, in $Q_{k-1}$. Pick $x_k\in Q_k$ and arcs $\beta_k$ in $Q_k$ from $x_k$ to
$x_{k+1}$ (T1), arranged so that non-consecutive ones are disjoint. The $\beta_k$ converge to $B$, and
Theorem~\ref{thm:limit-arc} (with trimming) gives an arc from $x_0$ to $B$ in $Q\cup\set B$.
\end{proof}

\formal{\lean{HPC.Moore.side_propagate}, \lean{HPC.Moore.accessible}, \lean{HPC.Moore.fa_theorem32}.}

Moore's own proof of Theorem~32 covers the arc by chains of the closed curves of Theorem~31 and uses a least-upper-bound
argument with two classes of points; the connectedness argument (a) replaces it.

\begin{theorem}[F.A.\ Theorems~33 and~34]\label{thm:fa33-34}
\leavevmode
\begin{enumerate}[label=(\alph*)]
\item If $A\in\fr R$ for a region $R$, $B\in R$, and $AB$ is an arc with $AB\setminus\set A\subseteq R$, then
$R\setminus(AB\setminus\set A)$ is connected.
\item If $AB$ is an arc lying in a region $R$, then $R\setminus AB$ is connected. In particular this holds for $R$ the
interior of a closed curve.
\end{enumerate}
\end{theorem}

\begin{proof}
We use the following consequence of Theorem~\ref{thm:fa32} (\emph{reaching an end}): if $O$ is open, $\alpha=AB$ an arc
with $\alpha\setminus\set A\subseteq O$, and $P\in O\setminus\alpha$ is joined in $O$ by an arc to a point of $\alpha$ other
than $A$, then $P$ is joined to $B$ by an arc in $O$ meeting $\alpha$ only in $B$. Indeed, the first point $X$ of that arc on
$\alpha$ is a limit point of the component $Q$ of $P$ in $O\setminus\alpha$; if $X=B$ we are done, otherwise propagate along
$\alpha$ from $X$ towards $B$ and apply accessibility.

(a) (Kline's proof.) Let $C,D\in R\setminus(AB\setminus\set A)$. An arc from $C$ to $B$ in $R$ meets $AB$, not at
$A\notin R$; so $C$ is joined to $B$ by an arc $\gamma$ meeting $AB$ only in $B$. Then $\Gamma=\gamma\cup AB$ is an arc from
$C$ to $A$ with $\Gamma\setminus\set A\subseteq R$. An arc from $D$ to $C$ in $R$ meets $\Gamma$ not at $A$; reaching the end $C$
of $\Gamma$, $D$ is joined to $C$ by an arc in $R$ meeting $\Gamma$ only in $C$. This arc misses $AB$, so $C$ and $D$ are in the
same component of $R\setminus(AB\setminus\set A)$.

(b) The same argument, with both ends of $AB$ in $R$: reaching an end gives an arc from $C$ to an end of $AB$, say $B$,
meeting $AB$ only there, and $\Gamma=\gamma\cup AB$; then $D$ reaches an end of $\Gamma$. If it reaches $C$ we are done. If it
reaches $A$, then $D$ and $C$ are the ends of an arc $D\dots A\dots B\dots C$ in $R$ containing $AB$, and the two sides of this
arc near an interior point (Corollary~\ref{cor:sides}) together with accessibility join $D$ to $C$ in $R\setminus AB$.
\end{proof}

\formal{\lean{HPC.Moore.reach_end}, \lean{HPC.Moore.joinZ}, \lean{HPC.Moore.fa_theorem33},
\lean{HPC.Moore.fa_theorem34_region}, \lean{HPC.Moore.fa_theorem34}.}

Moore gives two proofs of Theorem~34. The first (for interiors of closed curves) is not needed; the second, ``a
modification of a proof due to Mr.~J.~R.~Kline'', is the one formalized. Kline's argument does not cover the case in
which the arc from $D$ first meets $\Gamma$ at its end $A$; that case is the last sentence of the proof.

\section{Curves about points: Theorems 35, 36, 39 and 40}

\begin{theorem}[F.A.\ Theorems~39 and~40; accessibility of a closed curve]\label{thm:fa39-40}
Let $J$ be a closed curve and $M$ either $I(J)$ or $E(J)$.
\begin{enumerate}[label=(\alph*)]
\item Every point $O\in J$ is joined to every point $P\in M$ by an arc lying in $M\cup\set O$. In particular every point not on
$J$ is joined to $O$ by an arc meeting $J$ only in $O$.
\item Any two points $A\ne B$ of $J$ are joined by an arc whose open part lies in $I(J)$ and by one whose open part lies in
$E(J)$.
\end{enumerate}
\end{theorem}

\begin{proof}
(a) Let $f$ parametrize an arc $\alpha$ of $J$ with $f(1)=O$. The set $W=(S\setminus J)\cup f((0,1))$ is open (its
complement is the compact set $J\setminus f((0,1))$), it contains $f(u)$ for $0<u<1$, and $W\setminus\alpha=S\setminus J$. So the
component of $P$ in $W\setminus\alpha$ is the component of $P$ in $S\setminus J$, which is $M$ (Corollary~\ref{cor:interior}).
Since $\fr M=J$, the points $f(u)$ are limit points of $M$. Theorem~\ref{thm:fa32}(b) joins $P$ to $O$ by an arc lying, except for
$O$, in $M$.
(b) Take $P\in M$ and arcs from $P$ to $A$ and to $B$ as in (a); their union contains an arc from $A$ to $B$
(Lemma~\ref{lem:arc-ops}(c)), whose open part lies in $M$.
\end{proof}

\formal{\lean{HPC.Moore.access}, \lean{HPC.Moore.fa_theorem39}, \lean{HPC.Moore.fa_theorem40}.}

Moore proves Theorem~39 with a sequence of closed curves converging to $O$ and a limit arc; accessibility
(Theorem~\ref{thm:fa32}) does the same work.

\begin{theorem}[F.A.\ Theorems~35 and~36]\label{thm:fa35-36}
\leavevmode
\begin{enumerate}[label=(\alph*)]
\item Every nonempty region contains a closed curve.
\item If $O$ is a point of a region $R$, there is a closed curve $J\subseteq R$ with $O\in I(J)\subseteq R$.
\end{enumerate}
\end{theorem}

\begin{proof}[Proof (outline)]
(a) Take an arc in $R$ (Corollary~\ref{cor:fa16}) and, about an interior point of it, the curve $J_1$ of the pocket
construction (Theorem~\ref{thm:fa31}), which lies in $R$.

(b) Moore's proof, with one correction. Take a closed curve $J$ in a small region about a point $X\ne O$ of $R$, so that $O$ is
without $J$ and $I(J)\subseteq R$ (by (a) and (T2)); an arc $OA$ meeting $J$ only in $A$; a boundary point $C$ of $R$
(Theorem~\ref{thm:fa-elementary}(19)) and an arc $BC$ from $J$ to $C$ with $BC\setminus\set C\subseteq R$, missing $OA$ and meeting $J\cup I(J)$
only in $B$ (by Theorem~\ref{thm:fa33-34}(a) in $R\setminus OA$ and first points). Moore takes $BC$ inside a region about $C$ that need
not lie in $R$; the arc is cut at its first point outside $R$, which is again a boundary point of $R$, and that point is used as $C$.
Let $AF_1B$, $AH_1B$ be the two arcs of $J$ from $A$ to $B$ and $FZH$ an arc from $AF_1B$ to $AH_1B$, missing $OA\cup BC$, with open part
without $J$ and inside $R$ (Theorems~\ref{thm:fa20-21}, \ref{thm:fa33-34}). Then $B$ is without $ZHAFZ$ (the interior of that curve lies
in $R$, and the arc $BC$ reaches $C\notin R$ without meeting it), so by Theorem~\ref{thm:fa27} $A$ is within $ZHBFZ$, and so is $O$, since
$OA$ does not meet that curve. Finally $I(ZHBFZ)\subseteq R$ by (T2), as the curve lies in $R$.
\end{proof}

\formal{\lean{HPC.Moore.fa_theorem35}, \lean{HPC.Moore.fa_theorem36},
\lean{HPC.Moore.fa_theorem36_interior_subset}.}

\begin{remark}
As printed, Theorem~35 (``every region contains a closed curve'') fails for the empty region
(\lean{HPC.Moore.not_fa_theorem35_printed}); see Remark~\ref{rem:empty-region}. The statement that the interior lies in $R$ is
not part of the printed Theorem~36, but it is what Moore uses in 1925 (Theorem~19), and it follows from (T2).
\end{remark}

\section{Merging interiors: Theorems 37--43}

\begin{definition}\label{def:chord}
Let $J_1,J_2$ be closed curves. A \emph{chord} of $J_1$ over a closed set $K$ is an arc of $J_1$ with its ends on $K$ and its open
arc off $K$.
\end{definition}

Every point of $J_1\setminus K$ lies on exactly one chord over $K$ (cut $J_1$ at the nearest points of $K$ on either side), and two
chords whose open arcs meet are equal.

\begin{lemma}[Chords form a null family; Moore 1919, Theorem~A]\label{lem:chords-null}
Arcs of a closed curve $L$ with pairwise disjoint open arcs form a null family: for every $\varepsilon>0$ only finitely many of
them have a point at distance $\ge\varepsilon$ from their first end.
\end{lemma}

\begin{proof}
Parametrize $L$ by the circle $\R/\Z$. Each arc of $L$ is the image of a closed interval of the circle, and the open arcs being
disjoint, so are the open intervals; so for every $\delta>0$ only finitely many intervals have length $\ge\delta$. By uniform
continuity of the parametrization, an arc whose interval has length $<\delta$ has diameter $<\varepsilon$ for suitable $\delta$.
\end{proof}

\formal{\lean{HPC.Moore.chords_null}, \lean{HPC.Moore.m19_theoremA}.}

Moore's Theorem~A (1919) is the statement that a closed curve $L$ carries no infinite family of arcs, no two with a common
segment, each meeting two given disjoint closed sets; the lemma above is the form used here.

Theorems~37 and~38 of F.A.\ produce, for a point $O$ of $J$ not on another closed curve $C$, the innermost (resp.\ outermost) arc of
$C$ around $O$ with open part within (resp.\ without) $J$. The formal proofs replace Moore's limit-point arguments by the
observation that the chords of $C$ in question have only finitely many pairs of ends: distinct chords have disjoint open
intervals in a parametrization of $C$, and all have length bounded below, since their two ends lie in disjoint compact sets.

\formal{\lean{HPC.Moore.exists_maximal_chord}, \lean{HPC.Moore.fa_theorem37}, \lean{HPC.Moore.fa_theorem38}.}

\begin{theorem}[F.A.\ Theorem~41]\label{thm:fa41}
If the interiors $I_1,I_2$ of the closed curves $J_1,J_2$ meet, there is a closed curve $J\subseteq J_1\cup J_2$ with $I_1\cup
I_2\subseteq I(J)$.
\end{theorem}

\begin{proof}
\emph{Case I: $J_1\cap J_2$ has at most one point.} Then one interior contains the other, and $J$ is the bigger curve. Indeed,
if $J_2$ misses $I_1$, then $I_1$ (connected, meeting $I_2$, missing $J_2=\fr I_2$) lies in $I_2$. Otherwise $J_2$ has a point $x$
within $J_1$, and no point $y$ without $J_1$: the two arcs of $J_2$ from $x$ to $y$ would each meet $J_1$ at a point other than
$x,y$, giving two common points. So $J_2\subseteq\cl{I_1}$ and $I_2\subseteq I_1$ by (T2).

\emph{Case II, first sub-case: every point of $J_1$ is within or on $J_2$.} Then $I_1\subseteq I_2$ by (T2), and $J=J_2$;
symmetrically if $J_2\subseteq\cl{I_1}$.

\emph{Case II proper} (outline). Let $V$ be a component of $S\setminus(\cl{I_1}\cup\cl{I_2})$ with noncompact closure. Its frontier lies
on $J_1\cup J_2$, and $V$ misses the interior of every closed curve that misses $V$. There is a point $O$ of $J_2$ without $J_1$ which is
a limit point of $V$, and a second one $O'$. Split $J_2$ at $O,O'$ into two arcs; on each, replace every sub-arc cut off by a chord of $J_1$
over $J_2$ lying in $\cl V$ by that chord. The replaced sub-arcs have disjoint parameter intervals (no point inside such a sub-arc is a limit
point of $V$, by Theorem~\ref{thm:fa27}), the chords have disjoint open arcs and form a null family (Lemma~\ref{lem:chords-null}), so the
result is an arc from $O$ to $O'$ (Theorem~\ref{thm:arc-replace}). The two arcs so obtained meet only in $O,O'$, and their union $J$ is a
closed curve on $J_1\cup J_2$ containing the frontier of $V$. Then $V$ is the exterior of $J$, and $I_1\cup I_2$, which misses $J$ and $V$, lies
in $I(J)$.
\end{proof}

\formal{\lean{HPC.Moore.fa_theorem41_caseI_nested}, \lean{HPC.Moore.fa_theorem41_caseI},
\lean{HPC.Moore.fa_theorem41_of_subset}, \lean{HPC.Moore.fa_theorem41_caseII}, \lean{HPC.Moore.Gen.arc_replace},
\lean{HPC.Moore.fa_theorem41}.}

Moore ends the proof of Case~II with ``It may be proved that $J$ is a closed curve satisfying the requirements of Theorem~41''; the
replacement of a null family of arcs (Theorem~\ref{thm:arc-replace}) is what this sentence needs.

\begin{theorem}[F.A.\ Theorem~42]\label{thm:fa42}
Let $\mathcal G$ be a finite nonempty set of closed curves the union of whose interiors is connected. There is a closed curve lying on
the curves of $\mathcal G$ whose interior contains all their interiors.
\end{theorem}

\begin{proof}
Induction on the number of curves not yet absorbed. Suppose $\mathcal H\subsetneq\mathcal G$ is nonempty and $J_{\mathcal H}$ is a closed
curve on the curves of $\mathcal H$ whose interior contains theirs. If no curve of $\mathcal G\setminus\mathcal H$ had an interior meeting
$U=\bigcup_{\mathcal H}I(J)$, then $U$ and $V=\bigcup_{\mathcal G\setminus\mathcal H}I(J)$ would be disjoint nonempty open sets whose union
is the union of all the interiors, against connectedness. So some curve $q\in\mathcal G\setminus\mathcal H$ has an interior meeting
$U\subseteq I(J_{\mathcal H})$. Theorem~\ref{thm:fa41}
applied to $J_{\mathcal H}$ and $q$ absorbs $q$.
\end{proof}

\formal{\lean{HPC.Moore.fa_theorem42_of_41}, \lean{HPC.Moore.fa_theorem42}.}

The connectedness of the union of the interiors is essential; it is what guarantees an order in which each new interior meets the union
of the previous ones.

\begin{theorem}[F.A.\ Theorem~43, corrected]\label{thm:fa43}
Let $J$ be a closed curve, $O\in J$, and $\mathcal G$ a finite \emph{nonempty} set of closed curves whose interiors contain $O$ and none
of whose interiors contains $I(J)$. There are closed curves $J_1,J_2$ through $O$, lying on $J$ and the curves of $\mathcal G$, with $I(J_1)$
without $J$, $I(J_2)$ within $J$, and both inside every interior of a curve of $\mathcal G$.
\end{theorem}

\begin{proof}[Proof (outline)]
Moore gives no proof. The key step is the intersection analogue of Theorem~41: if $O\in J$ lies within $J'$ and $U$ is $I(J)$ or $E(J)$
(in the second case with $I(J)\not\subseteq I(J')$), and $J\cap J'$ has two points, then the component of $U\cap I(J')$ adjacent to $O$ is
the interior of a closed curve on $J\cup J'$ through $O$. Its proof is that of Case~II of Theorem~41 with this component in place of the
unbounded one. Iterating over $\mathcal G$ gives $J_2$ (starting from $J$) and $J_1$ (starting from the exterior side).
\end{proof}

\formal{\lean{HPC.Moore.dual_component}, \lean{HPC.Moore.fa_theorem43}.}

For $\mathcal G=\emptyset$ the printed Theorem~43 is false: its hypotheses hold vacuously and its conclusion would make the interior of $J$
lie in its own exterior (\lean{HPC.Moore.not_fa_theorem43_printed}).

\section*{Exercises}
\addcontentsline{toc}{section}{Exercises}

\begin{exercise}
In $\R^2$, verify Theorem~\ref{thm:fa25} for $J$ the unit circle and $AXB$ a diameter, and Theorem~\ref{thm:fa27} for $AXB$ a semicircle
outside the unit circle.
\end{exercise}

\begin{exercise}
Show by an example in $\R^2$ that Theorem~\ref{thm:fa42} fails if the union of the interiors is not connected.
\end{exercise}

\begin{exercise}
Prove Corollary~\ref{cor:split-interior} from Theorems~\ref{thm:fa24} and~\ref{thm:fa25}.
\end{exercise}

%======================================================================
\chapter{Double rulings and Moore's Theorem D}\label{ch:rulings}
%======================================================================

Moore's 1919 paper \emph{Concerning a set of postulates for plane analysis situs} proves that a space
satisfying Axioms~1--8 is a plane. Its device is the \emph{double ruling}: a grid of arcs across a
``quadrilateral'' (a closed curve with four marked points), dividing its interior into cells. Theorem~D says
that the interior of a quadrilateral can be ruled so finely that every cell lies inside some member of a given
cover. This chapter proves Theorem~D; Chapter~\ref{ch:plane} uses it to build coordinates.

The formal proof of Theorem~D takes a route different from Moore's: his Theorem~C, on which his proof rests, is
false as printed, and the formal proof goes through a corrected Theorem~C for curves that cross the quadrilateral
transversally.

Throughout, $S$ satisfies Axioms~1--8, and we use the results of Chapter~\ref{ch:curves} freely.

\section{Quadrilaterals and double rulings}

\begin{definition}\label{def:quad}
A \emph{quadrilateral} $ABCDA$ is a closed curve $J$ with interior $I$ and four points $A,B,C,D$ on it,
dividing $J$ into four arcs, the \emph{sides} $AB$, $BC$, $CD$, $DA$; consecutive sides meet only in their
common vertex, and opposite sides are disjoint.
\end{definition}

\formal{\lean{HPC.Moore.Space.Quad}.}

\begin{definition}\label{def:double-ruling}
A \emph{double ruling} of the quadrilateral $ABCDA$ with $n$ and $m$ arcs consists of
\begin{itemize}
\item points $X_1,\dots,X_n$ on $AB$ and $X_1',\dots,X_n'$ on $DC$, in the orders $A,X_1,\dots,X_n,B$ and
$D,X_1',\dots,X_n',C$, and pairwise disjoint arcs $X_iX_i'$ with open parts in $I$;
\item points $Y_1,\dots,Y_m$ on $BC$ and $Y_1',\dots,Y_m'$ on $AD$, in the orders $B,Y_1,\dots,Y_m,C$ and
$A,Y_1',\dots,Y_m',D$, and pairwise disjoint arcs $Y_jY_j'$ with open parts in $I$;
\item each $X_iX_i'$ meets each $Y_jY_j'$ in exactly one point.
\end{itemize}
Its \emph{subdivisions} (cells) are the components of $I$ minus the union of the arcs.
\end{definition}

\formal{\lean{HPC.Moore.Space.DoubleRuling}, \lean{HPC.Moore.Space.DoubleRuling.carrier},
\lean{HPC.Moore.Space.DoubleRuling.cells}.}

Moore's definition (footnote of p.~172) uses the same number $n$ of arcs in both directions and does not
require the $Y$-arcs to be disjoint. Allowing $n\ne m$ is a convenience only (a ruling can be padded to
$\max(n,m)$ arcs in each direction, \lean{HPC.Moore.doubleRuling_pad}); disjointness of the $Y$-arcs is needed.

It is often more convenient to forget the indexing: an \emph{arc ruling} is a pair $(\mathcal X,\mathcal Y)$ of finite
sets of arcs, the arcs of $\mathcal X$ pairwise disjoint from $AB$ to $CD$, those of $\mathcal Y$ pairwise disjoint from $BC$
to $DA$, all with open parts in $I$ and ends on the open sides, each arc of $\mathcal X$ meeting each arc of $\mathcal Y$ in
exactly one point. Indexing the arcs of $\mathcal X$ by the order of their ends on $AB$, their ends on $DC$ come in the same
order (an arc starting after $x$ on $AB$ lies in the half cut off by $x$ that contains $B$, so it ends after $x$ on $DC$);
so arc rulings are double rulings up to indexing.

\formal{\lean{HPC.Moore.Space.ArcRuling}, \lean{HPC.Moore.Space.ArcRuling.toDoubleRuling},
\lean{HPC.Moore.Space.DoubleRuling.toArcRuling}, \lean{HPC.Moore.Space.ArcRuling.param_mono}.}

\section{Counting the cells}

\begin{definition}\label{def:tiling}
Let $U$ be open and $F$ closed. A \emph{tiling} of $U$ relative to $F$ is a finite family of pairwise disjoint
interiors of closed curves lying on $F$ whose union is $U\setminus F$.
\end{definition}

\begin{lemma}\label{lem:tiling}
\leavevmode
\begin{enumerate}[label=(\alph*)]
\item If $\mathcal T$ is a tiling of $U$ relative to $F$, and $\alpha$ is an arc whose ends lie on $F$ and whose open
part lies in $U\setminus F$, then splitting the tile containing $\oarc\alpha$ (Corollary~\ref{cor:split-interior}) gives a tiling
relative to $F\cup\alpha$ with one more tile.
\item More generally, if $\alpha$ meets $F$ in exactly $k$ points, its ends among them, and the rest of $\alpha$ lies in
$U$, adding $\alpha$ adds $k-1$ tiles.
\item The tiles of a tiling are exactly the components of $U\setminus F$.
\end{enumerate}
\end{lemma}

\begin{proof}
(a) The tile $T=I(K)$ containing the connected set $\oarc\alpha$ has $\alpha$ as a cross-cut (its ends are on $F$, and lie on
$K=\fr T$ since $\oarc\alpha\subseteq T$ and the ends are not in any tile); Corollary~\ref{cor:split-interior} splits $T\setminus\alpha$
into two interiors of closed curves on $K\cup\alpha$. (b) Cut $\alpha$ at its $k$ points on $F$ into $k-1$ sub-arcs and apply (a)
to each in turn. (c) Each tile is connected and open, and the tiles are disjoint with union $U\setminus F$.
\end{proof}

\formal{\lean{HPC.Moore.Space.Tiling}, \lean{HPC.Moore.Space.Tiling.add_arc}, \lean{HPC.Moore.Space.Tiling.add_arc_finite}, \lean{HPC.Moore.Space.Tiling.components}.}

\begin{proposition}[Subdivision lemma]\label{prop:cells-count}
A double ruling with $n$ and $m$ arcs has exactly $(n+1)(m+1)$ subdivisions, each the interior of a closed curve lying on
$J$ and the arcs.
\end{proposition}

\begin{proof}
Start from the tiling $\set I$ of $I$ relative to $J$. Each arc $X_iX_i'$ meets $J$ and the previous $X$-arcs only in its two
ends, so it adds one tile; after the $X$-arcs there are $n+1$ tiles. Each arc $Y_jY_j'$ meets $J$, the $X$-arcs and the previous
$Y$-arcs in exactly $n+2$ points (its ends and its crossings with the $n$ arcs $X_i$), so it adds $n+1$ tiles. The total is
$1+n+m(n+1)=(n+1)(m+1)$, and by Lemma~\ref{lem:tiling}(c) the tiles are the subdivisions.
\end{proof}

\formal{\lean{HPC.Moore.Space.DoubleRuling.tiling}, \lean{HPC.Moore.doubleRuling_cells_ncard}.}

\section{Cutting, gluing, extending and refining}

An arc $x=XX'$ from a point $X$ of the open side $AB$ to a point $X'$ of the open side $CD$, with open part in $I$, cuts $ABCDA$
into the two quadrilaterals $AXX'DA$ and $XBCX'X$: by Corollary~\ref{cor:split-interior}, $I\setminus x$ is the disjoint union of
their interiors. Rotating the labels ($ABCDA\mapsto BCDAB$) gives the same for arcs from $BC$ to $DA$.

\formal{\lean{HPC.Moore.Space.Quad.splitX}, \lean{HPC.Moore.Space.Quad.rot}, \lean{HPC.Moore.Space.Quad.Cut}.}

\begin{lemma}[Cut and glue]\label{lem:cut-glue}
\leavevmode
\begin{enumerate}[label=(\alph*)]
\item A ruling containing the arc $x$ restricts to rulings of the two halves (the $Y$-arcs are cut at their crossing with $x$), and
off $x$ its arcs are those of the restricted rulings.
\item Conversely, if each half carries a ruling and the $Y$-arcs of the two halves end on $x$ at the same points, the $X$-arcs of both
halves together with $x$, and the unions of the $Y$-arcs meeting on $x$, form a ruling of $ABCDA$.
\item \emph{(Extension.)} Through every point $P$ of the open side $BC$ not on a $Y$-arc of a ruling $(\mathcal X,\mathcal Y)$ there is an arc
$y$ from $BC$ to $DA$ such that $(\mathcal X,\mathcal Y\cup\set y)$ is a ruling. Hence any finite set of points of the open side $BC$ can be
made ends of $Y$-arcs.
\end{enumerate}
\end{lemma}

\begin{proof}
(a) and (b) are bookkeeping with Corollary~\ref{cor:split-interior}. (c) Induction on the number of arcs. With no arcs,
Theorem~\ref{thm:fa39-40}(b) (in the form: a point of $J$ is joined to a point of the opposite side by an arc with open part in $I$) gives
$y$. Otherwise cut along an arc $x\in\mathcal X$ (or, if $\mathcal X=\emptyset$, along an arc of $\mathcal Y$, after rotating), extend in the half
containing $P$, then in the other half from the point where the new arc reached $x$, and glue by (b).
\end{proof}

\formal{\lean{HPC.Moore.Space.Quad.Cut.restrictA}, \lean{HPC.Moore.Space.Quad.Cut.restrictB},
\lean{HPC.Moore.Space.Quad.Cut.glue}, \lean{HPC.Moore.Space.ArcRuling.extendY},
\lean{HPC.Moore.Space.ArcRuling.extendY_finite}.}

\begin{definition}
Let $\mathrm{Good}$ be a property of subsets of $S$ inherited by subsets. A quadrilateral is \emph{rulable for Good} if it has a
double ruling all of whose subdivisions have property Good.
\end{definition}

\formal{\lean{HPC.Moore.Space.Rulable}, \lean{HPC.Moore.Space.DoubleRulable}.}

\begin{proposition}[Refinement]\label{prop:refine}
Let $(\mathcal X,\mathcal Y)$ be a ruling of $ABCDA$ such that every subdivision, made into a quadrilateral in any way, is rulable for
Good. Then $ABCDA$ has a ruling $(\mathcal X',\mathcal Y')$ with $\mathcal X\subseteq\mathcal X'$, $\mathcal Y\subseteq\mathcal Y'$, all of whose
subdivisions have property Good and lie in subdivisions of $(\mathcal X,\mathcal Y)$. In particular, if both halves of a cut are rulable for
Good, so is $ABCDA$.
\end{proposition}

\begin{proof}
Induction on the number of arcs. With no arcs the only subdivision is $I$ itself, rulable by hypothesis. Otherwise cut along an arc
$x\in\mathcal X$ (or of $\mathcal Y$, after rotating), refine the two halves by induction, use the extension lemma
(Lemma~\ref{lem:cut-glue}(c)) in each half to make the $Y$-arcs of the two halves end on $x$ at the same points (adding arcs keeps the
subdivisions Good, since Good is inherited by subsets), and glue (Lemma~\ref{lem:cut-glue}(b)).
\end{proof}

\formal{\lean{HPC.Moore.Space.ArcRuling.refine}, \lean{HPC.Moore.doubleRuling_refine},
\lean{HPC.Moore.Space.Quad.Cut.rulable}.}

Moore uses this refinement twice without proof (``the arcs of these rulings may be extended in such a way that there will result a new
double ruling $\tau_2$ of $ABCDA$ such that each arc of $\tau_1$ is also an arc of $\tau_2$'', pp.~175--176).

\section{Theorem C}

\begin{definition}\label{def:chord-quad}
A \emph{chord} of $ABCDA$ is an arc with its ends on $J$, different from $A,B,C,D$, and its open part in $I$. A closed curve $h$
\emph{touches $ABCDA$ from within} at a common point $P$ if $P$ is a common end of two arcs of $h$ meeting only in $P$ whose other
points lie in $I$. It \emph{crosses} $ABCDA$ at $P$ if near $P$ it consists of an arc with one half in $I$ and the other in $E(J)$.
\end{definition}

\formal{\lean{HPC.Moore.Space.IsChord}, \lean{HPC.Moore.Space.Crosses}, \lean{HPC.Moore.Space.IsCrossCut}.}

Moore's Theorem~C says: if $g$ is a closed curve meeting $ABCDA$ in finitely many points, none of them $A,B,C,D$, there is a double
ruling of $ABCDA$ each of whose subdivisions lies within or without $g$. \emph{As printed it is false}: let $g$ be a closed curve lying,
except for one point $P$ of the open side $AB$, within $ABCDA$ (in the plane: a circle inside a square, tangent to a side). Every point
of $g$ within $ABCDA$ must lie on an arc of the ruling (a subdivision containing it would meet $\fr I(g)=g$); near $P$ the only arc of
the ruling is the $X$-arc starting at $P$, if any; but $g$ leaves $P$ along two arcs within $ABCDA$ meeting only at $P$, which would both
have to start along this single arc. The configuration exists in every space satisfying Axioms~1--8.

\formal{\lean{HPC.Moore.no_ruling_of_two_branches}, \lean{HPC.Moore.tangent_configuration},
\lean{HPC.Moore.not_m19_theoremC_printed}.}

\begin{theorem}[Theorem~C, corrected]\label{thm:theoremC}
Let $h$ be a closed curve meeting $J$ in finitely many points, none of them $A,B,C,D$, and touching $ABCDA$ from within at none of
them. Then $ABCDA$ has a double ruling none of whose subdivisions meets $h$; so each subdivision lies within or without $h$.
\end{theorem}

The proof rests on two propositions.

\begin{proposition}[Chords]\label{prop:chords-rulable}
For every finite set $T$ of pairwise disjoint chords of $ABCDA$ there is a double ruling none of whose subdivisions meets a chord of $T$.
\end{proposition}

\begin{proof}[Proof (outline)]
A chord $t$ with ends $a,b$ divides $I$ into the interiors of $t\cup C$ and $t\cup C'$, where $C,C'$ are the two arcs of $J$ from $a$ to
$b$ (Corollary~\ref{cor:split-interior}); call the one whose arc of $J$ contains at most one vertex the \emph{pocket} of $t$ (for a chord
not joining opposite sides there is one) and that arc its \emph{base}. Give a chord weight $1$ if it joins opposite sides, $3$ if both
ends are on one side, and $2$ otherwise, and induct on the total weight. A chord joining opposite sides is itself an arc of a cut, and
the chords of the two halves have smaller total weight. Otherwise take a chord $t$ with maximal pocket, with an end $a$ on $AB$ (after
rotating); the bases of the maximal pockets are disjoint, so some point $d$ of the open side $CD$ lies on none of them; draw an arc $u$
from a point $c$ of $t$ to $d$ through the region outside all pockets (which is the interior of a closed curve, by removing the pockets
one at a time with Corollary~\ref{cor:split-interior}). Cut along $x=ac\cup u$: the remaining piece $cb$ of $t$ is a chord of one half of
smaller weight, and the other chords are chords of the halves with the same weights. Rulings of the halves avoiding their chords glue
(Proposition~\ref{prop:refine}) to a ruling of $ABCDA$ avoiding all chords.
\end{proof}

\formal{\lean{HPC.Moore.Space.Pocket}, \lean{HPC.Moore.chordWeight}, \lean{HPC.Moore.outer_tile},
\lean{HPC.Moore.exists_gap}, \lean{HPC.Moore.chords_rulable}.}

\begin{proposition}[A curve within]\label{prop:inside-rulable}
If the closed curve $h$ lies in $I$, there is a double ruling none of whose subdivisions meets $h$.
\end{proposition}

\begin{proof}
Build an arc $x$ from a point $P$ of the open side $AB$ to a point $H$ of the open side $CD$ out of an arc $Pa$ meeting $h$ only in $a$
(Theorem~\ref{thm:fa39-40}), an arc $ab$ of $h$, and an arc $bH$ meeting $h$ only in $b$ and missing $Pa$ (drawn in $I\setminus Pa$,
which is connected by Theorem~\ref{thm:fa33-34}). The other arc of $h$ from $a$ to $b$ is a chord of one half of the cut along $x$, and
the other half meets $h$ nowhere. Apply Proposition~\ref{prop:chords-rulable} to the first half and glue.
\end{proof}

\formal{\lean{HPC.Moore.inside_rulable}.}

\begin{proof}[Proof of Theorem~\ref{thm:theoremC}]
If $h$ meets $J$, the part of $h$ within $ABCDA$ lies on finitely many pairwise disjoint chords: write $h$ as the union of two arcs
between two of the common points, cut them at the (finitely many) common points, and keep the pieces within $ABCDA$; two such pieces
have no common point, since a common point would be an end of both and $h$ would touch $ABCDA$ from within there. By
Proposition~\ref{prop:chords-rulable} some ruling has no subdivision meeting these chords, hence none meeting $h$ (subdivisions lie in
$I$). If $h$ misses $J$, then $h\subseteq I$ (Proposition~\ref{prop:inside-rulable}) or $h\subseteq E(J)$ (the empty ruling: its only
subdivision $I$ misses $h$).
\end{proof}

\formal{\lean{HPC.Moore.extract_chords}, \lean{HPC.Moore.curve_rulable}, \lean{HPC.Moore.m19_theoremC_corrected},
\lean{HPC.Moore.m19_theoremC_crosscuts}, \lean{HPC.Moore.m19_theoremC_within}.}

\section{Theorem D}

\begin{proposition}[A nice curve about a compact set]\label{prop:nice-curve}
Let $F$ be a compact subset of the interior $V$ of a closed curve. There is a closed curve $h$ with interior $W\supseteq F$, $h\cup
W\subseteq V$, meeting $J$ in finitely many points, none of them $A,B,C,D$, and crossing $J$ at each of them.
\end{proposition}

\begin{proof}[Proof (outline)]
Enlarge $F$ to a compact connected set $K\subseteq V$ (by arcs, Theorem~\ref{thm:fa15}). About each point $O$ of $K\cap J$ draw a
\emph{lens}: a closed curve made of an arc within $J$ and an arc without $J$ with the same two ends $u,v$ on $J$, lying in $V$, whose
interior contains an open arc of $J$ about $O$ (the inner and outer arcs by Theorem~\ref{thm:fa28}, the interior computed by
Theorem~\ref{thm:fa27}); such a lens crosses $J$ at $u$ and $v$. Cover $K\cap J$ by finitely many lenses with all ends distinct, and the
rest of $K$ by finitely many closed curves missing $J$ with interiors in $V$ (Theorem~\ref{thm:fa35-36}). The union of the interiors is
connected (each meets the connected set $K$, which they cover), so Theorem~\ref{thm:fa42} merges these curves into one curve $h$ lying on
their union, whose interior contains all of them. Near each of its points on $J$, $h$ lies on a single lens, so it crosses $J$ there.
\end{proof}

\formal{\lean{HPC.Moore.lens_data}, \lean{HPC.Moore.lens_family}, \lean{HPC.Moore.exists_nice_curve}.}

A curve that crosses $J$ at a point does not touch it from within there (an arc ending at $P$ contains no two arcs from $P$ meeting only
in $P$), so Theorem~\ref{thm:theoremC} applies to the curve of Proposition~\ref{prop:nice-curve}.

\begin{theorem}[Moore 1919, Theorem~D]\label{thm:theoremD}
If every point on or within the quadrilateral $ABCDA$ is within some curve of a set $\mathcal G$ of closed curves, then $ABCDA$ has a double
ruling each of whose subdivisions lies within some curve of $\mathcal G$.
\end{theorem}

\begin{proof}
By compactness of $J\cup I$, finitely many curves $g_1,\dots,g_k$ of $\mathcal G$ suffice; we induct on $k$, for all quadrilaterals at once,
with Good the property ``lies within some curve of $\mathcal G$''. If $k=0$ then $J\cup I=\emptyset$, which is impossible; if $k\ge1$, let
$F$ be the set of points of $J\cup I$ not within $g_2,\dots,g_k$. It is compact and lies within $g_1$. Let $h$, with interior $W$, be a
nice curve about $F$ inside $I(g_1)$ (Proposition~\ref{prop:nice-curve}), and take, by Theorem~\ref{thm:theoremC}, a ruling none of whose
subdivisions meets $h$. Each subdivision $U$ is then within $W\subseteq I(g_1)$, and so Good; or $U$ misses the open set $W$, and then $\cl U\subseteq(J\cup I)\setminus W$ misses $F$, so every point on or within the
quadrilateral $U$ is within one of $g_2,\dots,g_k$, and by induction $U$ (made into a quadrilateral in any way) is rulable for Good. By the refinement proposition
(Proposition~\ref{prop:refine}), $ABCDA$ is rulable for Good.
\end{proof}

\formal{\lean{HPC.Moore.rulable_of_cover}, \lean{HPC.Moore.m19_theoremD}.}

\begin{corollary}[Theorem~D for regions]\label{cor:theoremD-regions}
If every point on or within $ABCDA$ lies in some region of a set $\mathcal G$ of regions, then $ABCDA$ has a double ruling each of whose
subdivisions lies in some region of $\mathcal G$.
\end{corollary}

\begin{proof}
Each point of a region lies within a closed curve whose interior lies in the region (Theorem~\ref{thm:fa35-36}(b)); apply
Theorem~\ref{thm:theoremD} to these curves.
\end{proof}

\formal{\lean{HPC.Moore.m19_theoremD_regions_of_curves}, \lean{HPC.Moore.m19_theoremD_regions}.}

\begin{remark}
Moore's Theorem~B (an arc from $A$ to $B$ within $J$ meeting a second closed curve $L$ in finitely many points) is proved in the formal
development with Moore's argument (\lean{HPC.Moore.m19_theoremB}); the corrected route to Theorem~D above does not need it.
\end{remark}

\section*{Exercises}
\addcontentsline{toc}{section}{Exercises}

\begin{exercise}
Draw a double ruling of the unit square with $2$ and $3$ arcs that is not made of straight segments, and check
Proposition~\ref{prop:cells-count}.
\end{exercise}

\begin{exercise}
Find the double ruling required by Theorem~\ref{thm:theoremC} for $ABCDA$ the unit square and $h$ a circle crossing the side $AB$ twice.
\end{exercise}

\begin{exercise}
Show that the weights in the proof of Proposition~\ref{prop:chords-rulable} decrease: the piece $cb$ of $t$ has smaller weight in its half
than $t$ has in $ABCDA$.
\end{exercise}

%======================================================================
\chapter{Grids, the Schoenflies theorem and Moore's characterization of the plane}\label{ch:plane}
%======================================================================

We now finish the proof that a space satisfying Moore's Axioms~1--8 is homeomorphic to $\R^2$. The route is
the one recommended by Moore's own construction, organized around the Schoenflies theorem:
\begin{enumerate}
\item a compact set carrying two families of ``grid lines'' of vanishing mesh is a closed square
(Section~\ref{sec:grid}; pure topology);
\item nested double rulings (Theorem~D) give such grid lines in the closure of the interior of any quadrilateral, so
this closure is a square with the sides of the quadrilateral going to the sides of the square (the
\emph{Schoenflies theorem} in Moore's space, Section~\ref{sec:schoenflies});
\item the closed region between two nested closed curves is a collar (the \emph{annulus lemma},
Section~\ref{sec:annulus});
\item the space is exhausted by nested closed curves (Moore's Theorem~F), and the square of the first curve and the
collars between consecutive ones are pasted onto the squares $[-(n+1),n+1]^2$ (Section~\ref{sec:paste}).
\end{enumerate}
Moore's 1919 paper reaches the plane through his Theorems~E, F and~G (a ``net'' of two families of open curves);
here Theorems~E and~G are obtained as corollaries of the main theorem.

\section{A compact set with a grid is a square}\label{sec:grid}

Let $K$ be a compact subset of a metric space.

\begin{definition}\label{def:family-lines}
A \emph{family of lines} in $K$ consists of indices $i$ with positions $t_i\in[0,1]$, including an index $i_0$ with
$t_{i_0}=0$ and an index $i_1$ with $t_{i_1}=1$, and for each $i$ two closed sets $L_i$, $R_i$ (the two closed
``sides'' of the line) such that:
\begin{enumerate}[label=(\arabic*)]
\item $L_i\cup R_i=K$; $L_i=K$ if $t_i=1$ and $R_i=K$ if $t_i=0$;
\item if $t_i<t_j$ then $L_i\subseteq L_j$, $R_j\subseteq R_i$ and $L_i\cap R_j=\emptyset$;
\item the \emph{line} $\ell_i=L_i\cap R_i$ is connected, and lies in $\cl{K\setminus L_i}$ if $t_i<1$ and in $\cl{K\setminus R_i}$ if
$t_i>0$.
\end{enumerate}
A \emph{pre-grid} is a pair of families of lines $(\mathcal F_X,\mathcal F_Y)$ in $K$ such that each line of $\mathcal F_X$ meets
each line of $\mathcal F_Y$ in exactly one point. Given finitely many lines of each family, including the extreme ones, the
\emph{box} between consecutive lines $j,j'$ of $\mathcal F_X$ and $m,m'$ of $\mathcal F_Y$ is $K\setminus(L_j\cup R_{j'}\cup
L_m\cup R_{m'})$. A \emph{grid} is a pre-grid such that
\begin{enumerate}[label=(\arabic*),resume]
\item \emph{(mesh)} for every $\varepsilon>0$ and any lines $i\in\mathcal F_X$, $k\in\mathcal F_Y$ there are finite subfamilies
containing $i$, $k$ and the extreme lines all of whose boxes have closures of diameter $\le\varepsilon$;
\item \emph{(parametrized lines)} some line $k_\varphi$ of $\mathcal F_Y$ carries an injective map $\varphi\colon[0,1]\to\ell_{k_\varphi}$ with
$\varphi(s)\in L_i\iff s\le t_i$ and $\varphi(s)\in R_i\iff t_i\le s$ for all lines $i$ of $\mathcal F_X$; likewise some line of $\mathcal F_X$
carries $\psi$ compatible with the positions of $\mathcal F_Y$.
\end{enumerate}
\end{definition}

\formal{\lean{HPC.Sq.Family}, \lean{HPC.Sq.PreGrid}, \lean{HPC.Sq.box}, \lean{HPC.Sq.Grid}.}

The model: $K=[0,1]^2$, the lines of $\mathcal F_X$ the horizontal segments $y=t$ with $L=\set{y\le t}$, $R=\set{y\ge t}$, and
those of $\mathcal F_Y$ the vertical ones.

\begin{theorem}[A compact grid is a square]\label{thm:grid-square}
If the compact set $K$ carries a grid, there is a homeomorphism $f$ of $[0,1]^2$ onto $K$ sending the sides $x=0$, $x=1$,
$y=0$, $y=1$ onto the first and last lines of $\mathcal F_Y$ and of $\mathcal F_X$. Moreover $f$ can be chosen to restrict to
the parametrizations $\varphi$, $\psi$ on the corresponding lines.
\end{theorem}

\begin{proof}
For $z\in K$ define its \emph{position} in $\mathcal F_X$, $\mathrm{pos}_X(z)=\inf\set{t_i:z\in L_i}$, and similarly $\mathrm{pos}_Y$.
Then $t_i\le\mathrm{pos}_X(z)\le t_j$ whenever $z\in R_i\cap L_j$, and $\mathrm{pos}_X=t_i$ on the line $\ell_i$.

\emph{Positions of lines are dense in $[0,1]$.} The parametrized line gives $\mathrm{pos}_X(\varphi(s))=s$. If no line had its
position in an interval $(a,b)\subseteq[0,1]$, the points $\varphi(a)\ne\varphi(b)$ would lie, for every finite subfamily, in the
closure of one box (between the consecutive lines around $(a,b)$), against the mesh condition with $\varepsilon<d(\varphi(a),
\varphi(b))$.

\emph{Limit lemma.} If $\mathrm{pos}_X(z)=t_i$ then $z\in\ell_i$: by density there are lines at positions approaching $t_i$ from
both sides, and $z$ lies on their $R$-sides from below and $L$-sides from above; by (2), (3) and closedness, $z\in L_i\cap R_i$.

\emph{Continuity.} $\set{z:\mathrm{pos}_X(z)<c}=\bigcup_{t_i<c}(K\setminus R_i)$ and $\set{z:\mathrm{pos}_X(z)>c}=\bigcup_{t_i>c}
(K\setminus L_i)$, by density and the limit lemma; both are open in $K$.

\emph{Injectivity of $\Phi=(\mathrm{pos}_Y,\mathrm{pos}_X)$.} If $z,w$ have the same positions, then for every finite subfamily they lie in
the closure of the same box (or on its boundary lines), so $d(z,w)\le\varepsilon$ for every $\varepsilon$ by the mesh condition.

\emph{Surjectivity.} The crossing point of the lines $i\in\mathcal F_X$, $k\in\mathcal F_Y$ has image $(t_k,t_i)$. These points are dense in
$[0,1]^2$ and $\Phi(K)$ is compact, so $\Phi(K)=[0,1]^2$.

So $\Phi\colon K\to[0,1]^2$ is a continuous bijection from a compact space onto a Hausdorff space, a homeomorphism; let $f=\Phi^{-1}$. The
sides go where stated since $\mathrm{pos}_X(z)=0$ iff $z\in L_{i_0}$, and so on. On the parametrized line $\ell_{k_\varphi}$ we have
$\Phi(\varphi(s))=(t_{k_\varphi},s)$, which gives the last claim.
\end{proof}

\formal{\lean{HPC.Sq.Family.pos}, \lean{HPC.Sq.Grid.dense_t}, \lean{HPC.Sq.Grid.mem_L_of_pos_eq},
\lean{HPC.Sq.Grid.continuousOn_pos}, \lean{HPC.Sq.Grid.eq_of_pos_eq}, \lean{HPC.Sq.Grid.image_pos},
\lean{HPC.Sq.Grid.exists_square}, \lean{HPC.Sq.Grid.exists_square_param}.}

\begin{lemma}[Boxes are components]\label{lem:boxes}
Given finitely many lines of each family of a pre-grid (including the extreme ones), the open boxes cut out by consecutive lines
are pairwise disjoint, nonempty, and cover $K$ minus the lines. If $U\subseteq K$ is a set such that $U$ minus the lines has
exactly as many components as there are boxes, every box lies in a single component.
\end{lemma}

\begin{proof}
Disjointness and covering follow from (1)--(2); nonemptiness from the crossing of lines and (3). Each component of $U$ minus the
lines is a union of (connected pieces of) boxes; each box meets some component, and a counting argument shows that if some box met
two components there would be more components than boxes.
\end{proof}

\formal{\lean{HPC.Sq.PreGrid.box_nonempty}, \lean{HPC.Sq.PreGrid.box_sub_component}.}

\section{The Schoenflies theorem in Moore's space}\label{sec:schoenflies}

We return to a space $S$ satisfying Axioms~1--8, with a fixed metric (Proposition~\ref{prop:bridge}).

\begin{lemma}[Nested rulings]\label{lem:nested-rulings}
For any sequence $\mathcal G_1,\mathcal G_2,\dots$ of covers of $S$ by regions there are double rulings $\beta_1,\beta_2,\dots$ of a
quadrilateral $ABCDA$ such that every subdivision of $\beta_k$ lies in a region of $\mathcal G_k$, every arc of $\beta_k$ is an arc of
$\beta_{k+1}$, and every subdivision of $\beta_{k+1}$ lies in a subdivision of $\beta_k$.
\end{lemma}

\begin{proof}
$\beta_1$ is given by Theorem~D for regions (Corollary~\ref{cor:theoremD-regions}). Given $\beta_k$, every subdivision of $\beta_k$,
made into a quadrilateral in any way, is rulable with subdivisions in regions of $\mathcal G_{k+1}$ (Corollary~\ref{cor:theoremD-regions}
again); the refinement proposition (Proposition~\ref{prop:refine}) gives $\beta_{k+1}$.
\end{proof}

\formal{\lean{HPC.Moore.nested_rulings}, \lean{HPC.Moore.nested_rulings_fund}.}

\begin{theorem}[Schoenflies; Moore 1919, Theorem~E]\label{thm:schoenflies}
For every quadrilateral $ABCDA$ there is a homeomorphism $f$ of $[0,1]^2$ onto $J\cup I$ with $f(0,0)=A$, $f(0,1)=B$, $f(1,1)=C$,
$f(1,0)=D$, sending the sides $x=0$, $y=1$, $x=1$, $y=0$ of the square onto $AB$, $BC$, $CD$, $DA$. The homeomorphism can be chosen
to agree with any given parametrizations of the four sides.
\end{theorem}

\begin{proof}
Let $\mathcal G_k$ be the cover by regions of diameter $<1/k$, and $\beta_k$ the nested rulings of Lemma~\ref{lem:nested-rulings}. Let
$\mathcal A_X$ be the set of all arcs from $AB$ to $CD$ of all the $\beta_k$, and $\mathcal A_Y$ the set of all arcs from $BC$ to $DA$. Arcs of
$\mathcal A_X$ are pairwise disjoint and each meets each arc of $\mathcal A_Y$ in exactly one point (two arcs occur together in some $\beta_k$,
since the rulings are nested).

\emph{Lines.} Fix a parametrization $f_{AB}$ of $AB$ from $A$ to $B$. Record an arc $x\in\mathcal A_X$ from $X\in AB$ to $X'\in CD$ by its
\emph{line record}: the parameter $t$ of $X$, and the closures $L,R$ of the two halves $AXX'DA$ and $XBCX'X$ of the cut along $x$. Record
the side $DA$ as $(0,DA,J\cup I)$ and $BC$ as $(1,J\cup I,BC)$. These form a family of lines in $J\cup I$ (Definition~\ref{def:family-lines}):
the two sides of a cut are closed, cover $J\cup I$ and meet in the arc, which lies in the closure of each complement; along $AB$ the point
$f_{AB}(s)$ is on the $L$-side iff $s\le t$; and two disjoint arcs are nested in the order of their ends on $AB$. Together with the family
built from $\mathcal A_Y$ and $f_{BC}$ (with $AB$, $CD$ as extreme lines) this is a pre-grid, and the sides $AB$, $BC$ with $f_{AB},f_{BC}$
are parametrized lines.

\emph{Mesh.} Given a finite set of lines, all of them belong to some $\beta_k$ (together with the sides). The $(n+1)(m+1)$ boxes cut out by
the arcs of a ruling are nonempty and disjoint (Lemma~\ref{lem:boxes}) and there are exactly $(n+1)(m+1)$ subdivisions
(Proposition~\ref{prop:cells-count}); so each box lies in one subdivision, hence in a region of diameter $<1/k$. Adding lines only refines
boxes, so the mesh condition holds.

By Theorem~\ref{thm:grid-square} the compact set $J\cup I$ is a square with the stated sides. For prescribed parametrizations of all four
sides, compose with a homeomorphism of the square that fixes two adjacent sides pointwise and reparametrizes the other two
(\lean{HPC.Sq.exists_square_reparam}).
\end{proof}

\formal{\lean{HPC.Moore.lineFam}, \lean{HPC.Moore.lr_mono}, \lean{HPC.Moore.preGrid},
\lean{HPC.Moore.preGrid_box_sub_cell}, \lean{HPC.Moore.m19_schoenflies_proof}, \lean{HPC.Moore.m19_schoenflies_param},
\lean{HPC.Moore.m19_schoenflies_boundary_proof}, \lean{HPC.Moore.m19_schoenflies_boundary}.}

Theorem~E as Moore states it (two sets of arcs across $ABCDA$, one through each interior point in each direction, each arc of one set
meeting each arc of the other once) follows by taking the images of the horizontal and vertical segments under $f$
(\lean{HPC.Moore.m19_theoremE}).

\begin{corollary}\label{cor:schoenflies-curve}
For every closed curve $J$ there is a homeomorphism of $[0,1]^2$ onto $J\cup I(J)$ sending the boundary of the square onto $J$.
\end{corollary}

\begin{proof}
Choose four points on $J$ to make it a quadrilateral.
\end{proof}

\formal{\lean{HPC.Moore.exists_quad}, \lean{HPC.Moore.exists_square_of_interior}.}

\section{The annulus lemma}\label{sec:annulus}

\begin{theorem}[Annulus lemma]\label{thm:annulus}
Let $J,J'$ be closed curves with $J\cup I(J)\subseteq I(J')$. There is a homeomorphism $H$ of $J\times[0,1]$ onto the closed annulus
$(J'\cup I(J'))\setminus I(J)$ with $H(z,0)=z$ for $z\in J$ and $H(J\times\set1)=J'$.
\end{theorem}

\begin{proof}
\emph{Two cuts.} By the Schoenflies theorem for $J'$ (Corollary~\ref{cor:schoenflies-curve}) there are coordinates on $J'\cup I(J')$ in
which it is a square. On a horizontal line meeting the compact set $J\cup I(J)$, take its extreme points $P,Q$ on $J\cup I(J)$ and the
horizontal segments from them to the boundary: they give disjoint arcs $c_1=PP'$, $c_2=QQ'$ from $J$ to $J'$, each meeting $J\cup I(J)$
only in its first end and $J'$ only in its last end.

\emph{Two quadrilaterals.} Let $e$ be an arc from $P$ to $Q$ with open part in $I(J)$ (Theorem~\ref{thm:fa39-40}). Then $c_1\cup e\cup c_2$
is a cross-cut of $I(J')$, splitting it into the interiors $R_1,R_2$ of $c_1\cup e\cup c_2\cup b_i$, where $b_1,b_2$ are the arcs of $J'$
from $P'$ to $Q'$ (Corollary~\ref{cor:split-interior}). One arc $a_1$ of $J$ from $P$ to $Q$ lies in $\cl{R_1}$ and the other, $a_2$, in
$\cl{R_2}$; cutting $R_i$ along $a_i$ leaves the interior $A_i$ of the closed curve $K_i=c_1\cup b_i\cup c_2\cup a_i$. The two closed sets
$K_i\cup A_i$ fill the closed annulus and meet only along $c_1$ and $c_2$.

\emph{Gluing.} Each $K_i\cup A_i$ is a quadrilateral with sides $a_i$, $c_2$, $b_i$, $c_1$. Choose parametrizations of $c_1$, $c_2$ and of
the arcs $a_i$ (compatible at $P,Q$), and apply the Schoenflies theorem with prescribed sides (Theorem~\ref{thm:schoenflies}) to both
quadrilaterals so that the two maps agree on $c_1$ and $c_2$. Define $H(z,s)$, for $z\in a_i$, as the value at height $s$ above the
parameter of $z$ in the square of $K_i\cup A_i$. The two halves agree on $c_1,c_2$, so $H$ is a continuous bijection from the compact
$J\times[0,1]$ onto the annulus, hence a homeomorphism.
\end{proof}

\formal{\lean{HPC.Moore.exists_cuts}, \lean{HPC.Moore.annulus_decomp}, \lean{HPC.Moore.exists_collar},
\lean{HPC.Moore.m19_annulus}, \lean{HPC.Moore.m19_annulus_square}.}

Moore asserts that the annulus can be divided by a double ruling into nine cells, the central one bounded by the inner curve; this is not
needed (two quadrilaterals suffice) and is not proved in the formal development.

\section{Exhaustion and pasting}\label{sec:paste}

\begin{lemma}\label{lem:curve-about-compact}
$S$ is connected, and every compact set $M\subseteq S$ lies within some closed curve.
\end{lemma}

\begin{proof}
Connectedness is Theorem~\ref{thm:fa-elementary}(22). Cover $M$ by finitely many interiors of closed curves (Theorem~\ref{thm:fa35-36}(b))
and join each of them by an arc to a base point (Theorem~\ref{thm:fa15}); the union $M_1$ of these arcs and the closures of the interiors
is compact and connected. Cover $M_1$ by finitely many interiors of closed curves, each meeting $M_1$; their union is connected (each meets
the connected $M_1$, which they cover). By Theorem~\ref{thm:fa42} some closed curve has an interior containing all of them, hence $M$.
\end{proof}

\formal{\lean{HPC.Moore.preconnectedSpace}, \lean{HPC.Moore.exists_curve_about_compact}.}

\begin{theorem}[Moore 1919, Theorem~F]\label{thm:theoremF}
There are closed curves $J_0,J_1,\dots$ with $J_n\cup I(J_n)\subseteq I(J_{n+1})$ and $\bigcup_nI(J_n)=S$.
\end{theorem}

\begin{proof}
Let $K_1,K_2,\dots$ be the fundamental regions of Axiom~1 (every point lies in some $K_n$). Take $J_0$ about $\cl{K_0}$, and $J_{n+1}$ a
closed curve within which lies the compact set $J_n\cup I(J_n)\cup\cl{K_{n+1}}$ (Lemma~\ref{lem:curve-about-compact}).
\end{proof}

\formal{\lean{HPC.Moore.m19_theoremF_proof}, \lean{HPC.Moore.m19_theoremF}.}

\begin{proposition}[Pasting]\label{prop:paste}
Let $X$ be a Hausdorff space and $J_n,I_n\subseteq X$ ($n\ge0$) with $I_n$ open, $J_n\cap I_n=\emptyset$, $J_n\cup I_n\subseteq I_{n+1}$ and
$\bigcup I_n=X$. Suppose $J_0\cup I_0$ is the image of a homeomorphism of the square $[-1,1]^2$ sending its boundary onto $J_0$, and each
$J_n$ has a collar in $J_{n+1}\cup I_{n+1}$ ending on $J_{n+1}$ (as in Theorem~\ref{thm:annulus}). Then $X$ is homeomorphic to $\R^2$.
\end{proposition}

\begin{proof}
Inductively extend the embedding $\Psi_n$ of $[-(n+1),n+1]^2$ onto $D_n=J_n\cup I_n$ to $\Psi_{n+1}$ on $[-(n+2),n+2]^2$, mapping the square
annulus $n+1\le\nrm p\le n+2$ onto the collar by $\Psi_{n+1}(p)=H\big(\Psi_n((n+1)p/\nrm p),\,\nrm p-(n+1)\big)$ (radial extension). Then
$D_n\subseteq\interior D_{n+1}$, and the $\Psi_n$ agree where both are defined, so they define a bijection $\Psi\colon\R^2\to X$. It is
continuous (each point of $\R^2$ has a neighbourhood in some $[-(n+1),n+1]^2$) and open (each point of $X$ lies in the interior of some
$D_n$, on which $\Psi^{-1}=\Psi_{n+1}^{-1}$ is continuous).
\end{proof}

\formal{\lean{HPC.Sq.exhaust_homeomorph}, \lean{HPC.Sq.plane_of_collars}.}

\begin{theorem}[Moore's characterization of the plane]\label{thm:moore-plane}
A space satisfying Axioms~1--8 is homeomorphic to $\R^2$ (with the topology generated by its regions).
\end{theorem}

\begin{proof}
Take the curves of Theorem~\ref{thm:theoremF}. $J_0\cup I(J_0)$ is a square with boundary $J_0$ (Corollary~\ref{cor:schoenflies-curve},
rescaled to $[-1,1]^2$), and each annulus is a collar (Theorem~\ref{thm:annulus}). Proposition~\ref{prop:paste} applies, since $S$ is
Hausdorff.
\end{proof}

\formal{\lean{HPC.Moore.m19_plane_proof}, \lean{HPC.Moore.m19_plane}.}

\begin{corollary}[Moore 1919, Theorem~G and the final paragraph]\label{cor:theoremG}
\leavevmode
\begin{enumerate}[label=(\alph*)]
\item There are two sets $\mathcal G_1,\mathcal G_2$ of open curves (images of closed embeddings of $\R$) such that through each point there is
exactly one curve of each set, and each curve of $\mathcal G_1$ meets each curve of $\mathcal G_2$ in exactly one point.
\item There is a bijection $P\mapsto(x_P,y_P)$ of $S$ onto $\R^2$ such that $P$ is the sequential limit of $P_1,P_2,\dots$ (in Moore's sense)
if and only if $x_{P_n}\to x_P$ and $y_{P_n}\to y_P$.
\end{enumerate}
\end{corollary}

\begin{proof}
Transport the horizontal and vertical lines, and convergence, by the homeomorphism of Theorem~\ref{thm:moore-plane}; Moore's sequential limits are
the limits of the region topology (Proposition~\ref{prop:bridge}).
\end{proof}

\formal{\lean{HPC.Moore.m19_theoremG}, \lean{HPC.Moore.m19_openCurve}, \lean{HPC.Moore.m19_final}.}

\section*{Exercises}
\addcontentsline{toc}{section}{Exercises}

\begin{exercise}
Show that the mesh condition in Definition~\ref{def:family-lines} cannot be dropped: give a compact set with two families of ``lines'' satisfying
(1)--(3) and (5) that is not a square.
\end{exercise}

\begin{exercise}
In the plane, carry out the proof of Theorem~\ref{thm:annulus} for two concentric circles, with $c_1,c_2$ horizontal segments.
\end{exercise}

\begin{exercise}
Show that Proposition~\ref{prop:paste} fails without the hypothesis $D_n\subseteq\interior D_{n+1}$.
\end{exercise}

%======================================================================
\chapter{The space of elements is a plane}\label{ch:quotient}
%======================================================================

We return to a Moore decomposition $G$ of the plane (Definition~\ref{def:moore-decomp}) and finish
Moore's 1925 paper: Theorems~13--22. The plan of the paper is simple. Theorems~1--12 (Chapters
\ref{ch:usc}--\ref{ch:jordan}) show that arcs and simple closed curves of elements behave as in the plane,
with the Jordan curve theorem for curves of elements (Theorem~\ref{thm:m25-11}) as the key result.
Theorems~13--19 develop the cross-cut theory of curves of elements, which shows that the
interiors of simple closed curves of elements, the \emph{regions in the restricted sense}, form a basis of the
topology. Theorem~21 checks that with these regions $G$ satisfies Moore's Axioms~1--8, and Theorem~22 then
follows from Moore's characterization of the plane (Theorem~\ref{thm:moore-plane}).

Throughout, $J$ denotes a simple closed curve of elements; by Theorem~\ref{thm:m25-11},
$G\setminus J$ has exactly two components, the bounded \emph{interior} $I(J)$ and the unbounded
\emph{exterior} $E(J)$, and $J$ is the frontier of each. (Formally, $I(J)$ is the set of elements off
$J$ whose component of $G\setminus J$ has bounded point set.) An element $x\notin J$ is
\emph{within} $J$ if $x\in I(J)$ and \emph{without} $J$ if $x\in E(J)$. For an arc $pxq$ of elements,
$\oarc{pxq}=pxq\setminus\set{p,q}$.

\formal{\lean{HPC.Moore.sccInterior}, \lean{HPC.Moore.sccExterior}, \lean{HPC.Moore.sccInterior_spec},
\lean{HPC.Moore.closure_sccInterior}, \lean{HPC.Moore.closure_sccExterior}.}

\section{Arcs and curves of elements are topological arcs and curves}

The cross-cut theory of F.A.\ was developed in Chapter~\ref{ch:curves} for a space satisfying Moore's
axioms, but most of it uses the axioms only through a few consequences. Those consequences were isolated
in topological form: a Hausdorff, locally connected space $X$ with a family of open sets (``regions'')
such that
\begin{enumerate}[label=(\roman*)]
\item regions and the complements of their closures are connected, and region closures are compact,
while $X$ is not;
\item two points of a connected open set are joined by an arc in it (F.A.\ Theorem~15);
\item every simple closed curve is the frontier of exactly two components of its complement, one a region
(Axiom~8 together with F.A.\ Theorem~21).
\end{enumerate}
To apply this to $G$ we need Moore's arcs and curves of elements to be the topological ones.

\begin{proposition}\label{prop:elt-top}
\leavevmode
\begin{enumerate}[label=(\alph*)]
\item A set of elements is a simple continuous arc of elements from $A$ to $B$ if and only if it is the image of
an embedding of $[0,1]$ sending $0$, $1$ to $A$, $B$.
\item A set of elements is a simple closed curve of elements if and only if it is the union of two such arcs
with only their ends in common, if and only if it is homeomorphic to a circle.
\item Two arcs of elements from $p$ to $q$ with only $p,q$ in common form a simple closed curve of
elements.
\item $G$ is Hausdorff and arcwise connected; every connected open set of elements is arcwise connected.
\end{enumerate}
\end{proposition}

\begin{proof}
(a) An arc of elements is compact (Proposition~\ref{prop:arc-elt-top}), metrizable, connected, and every
element other than $A,B$ cuts it: by the cut-point characterization of the arc
(Theorem~\ref{thm:irr-arc} and Proposition~\ref{prop:cutarc}) it is a topological arc. Conversely a
topological arc is closed, connected, compact (so bounded: the point set of a compact set of elements is
compact, Corollary~\ref{cor:quot-tools}), and every interior point cuts it.

(b) A curve of elements is split by any two of its elements into two arcs of elements
(Theorem~\ref{thm:m25-5-10}(e)), which are topological arcs by (a). Conversely, two topological arcs with only
their ends in common are cut by any two points (Proposition~\ref{prop:scc-cut}). The circle characterization is
Theorem~\ref{thm:circle-char}.

(c) is (b). (d) Hausdorff: Theorem~\ref{thm:quot-metrizable}. Arcs: Moore's proof of F.A.\ Theorem~15 goes
through in $G$ with the Moore regions (Theorem~\ref{thm:fa-elementary} applied as in Chapter~\ref{ch:axioms});
alternatively, $G$ is a connected, locally connected, complete metric space, so connected open sets are
arcwise connected.
\end{proof}

\formal{\lean{HPC.Moore.isArcElt_iff_tarc}, \lean{HPC.Moore.isSCCElt_iff_tscc}, \lean{HPC.Moore.isSCCElt_union},
\lean{HPC.Moore.t2Space_quot}, \lean{HPC.Moore.arcConn_quot}, \lean{HPC.Jones.tscc_homeomorph_circle}.}

\begin{definition}\label{def:restricted}
A \emph{region in the restricted sense} is the interior $I(J)$ of some simple closed curve of elements $J$.
The \emph{space of elements with the restricted regions} is $G$ with these sets as Moore's regions.
\end{definition}

\formal{\lean{HPC.Moore.restrictedRegion}, \lean{HPC.Moore.spaceRestricted}.}

\begin{lemma}\label{lem:restricted-moore}
A region in the restricted sense is a Moore region (Definition~\ref{def:moore-region}). The restricted
regions satisfy (i)--(iii) above.
\end{lemma}

\begin{proof}
By Theorem~\ref{thm:m25-11}, $I(J)$ is open, connected and bounded with frontier $J$, which is connected. Its
closure $I(J)\cup J$ has bounded point set, so it is compact; the complement of the closure is $E(J)$, connected.
$G$ is not compact, since $\pts G=\R^2$ is unbounded. (ii) is Proposition~\ref{prop:elt-top}(d), and (iii) is
Theorem~\ref{thm:m25-11} together with Proposition~\ref{prop:elt-top}(b).
\end{proof}

\formal{\lean{HPC.Moore.mooreRegion_of_restricted}, \lean{HPC.Moore.th21Prop_restricted},
\lean{HPC.Moore.ax8Prop_restricted}, \lean{HPC.Moore.isCompact_closure_restricted},
\lean{HPC.Moore.not_isCompact_univ_quot}.}

Moore remarks (p.~423) that F.A.\ Theorem~21 ``may be easily proved with the help of Theorem~11''. This is
correct as it stands: only Theorem~\ref{thm:m25-11} is needed for (iii).

\section{Theorems 13 and 14}

\begin{lemma}[Janiszewski's theorem for elements]\label{lem:quot-jan}
Let $A,B$ be closed sets of elements with bounded point sets such that $A\cap B$ is connected (or empty). If
neither $A$ nor $B$ separates the elements $x$ and $y$, then $A\cup B$ does not separate them.
\end{lemma}

\begin{proof}
Components of $G\setminus A$ correspond to components of $\R^2\setminus\pts A$ through $H\mapsto\pts H$
(Proposition~\ref{prop:pts-conn}). The point sets $\pts A$, $\pts B$ are compact, $\pts A\cap\pts B=\pts{(A\cap B)}$
is connected (or empty), and Janiszewski's theorem in the plane (\vfour{Theorem~2.7}) applies to
points of $x$ and $y$.
\end{proof}

\formal{\lean{HPC.Moore.quot_janiszewski}.}

\begin{theorem}[Moore, Theorem~13]\label{thm:m25-13}
Let $R$ be a Moore region and $k$ either a single element or an arc of elements all of whose elements,
except possibly one end, belong to $R$. Then $R\setminus k$ is open and connected.
\end{theorem}

\begin{proof}
$R\setminus k$ is open since $k$ is closed. Let $x,y\in R\setminus k$ and $B=\fr R$. $B$ does not separate $x$
from $y$, since $R$ is connected and misses $B$. Nor does $k$: $G\setminus k$ is connected
(Theorem~\ref{thm:m25-9}, or Corollary~\ref{cor:pts-nonsep} for a single element). Now $k\cap B$ is empty or the one
end of $k$ outside $R$, so by Lemma~\ref{lem:quot-jan} $k\cup B$ does not separate $x$ from $y$. The component of
$G\setminus(k\cup B)$ containing $x$ and $y$ misses $B$ and meets $R$, so lies in $R$, and it misses $k$.
Finally $R\setminus k\ne\emptyset$ (a point not in Moore's text), since an arc of elements $k$ contains no nonempty open set of elements.
Indeed, if $U\subseteq k$ is open and nonempty, choose $y\in U$ other than the ends of $k$ and a Moore region $V\ni y$
inside $U$ (Theorem~\ref{thm:m25-4}). Then $V$ is a connected open subset of the arc $k$, so an interval of the arc containing $y$
in its interior, and $V\setminus\set y$ is not connected; but by the first part of this proof (with $V$ for $R$ and
$\set y$ for $k$) it is.
\end{proof}

\formal{\lean{HPC.Moore.not_subset_arc_of_cut}, \lean{HPC.Moore.m25_theorem13}.}

\begin{theorem}[Moore, Theorem~14]\label{thm:m25-14}
Every Moore region contains a simple closed curve of elements.
\end{theorem}

\begin{proof}
A Moore region $R$ is infinite, so it contains elements $p\ne q$, joined by an arc $pxq$ in $R$
(Proposition~\ref{prop:elt-top}(d)). Let $x$ be an element of $pxq$ other than $p,q$. By
Theorem~\ref{thm:m25-13}, $R\setminus\set x$ is a domain, so it contains an arc $pyq$ from $p$ to $q$. Moore:
``it is easy to see that the sum of the arcs $pxq$ and $pyq$ contains as a subset a simple closed curve''.
Indeed, follow $pxq$ from $x$ towards $p$ until the first element $a$ on $pyq$, and towards $q$ until the first
element $b$ on $pyq$ (both exist, since $p,q\in pyq$ and $pyq$ is closed). The sub-arc $axb$ of $pxq$ meets $pyq$ only
in $a$ and $b$, and with the sub-arc of $pyq$ from $a$ to $b$ it forms a simple closed curve of elements
(Proposition~\ref{prop:elt-top}(c)).
\end{proof}

\formal{\lean{HPC.Moore.exists_tscc_of_arcs}, \lean{HPC.Moore.m25_theorem14}.}

\section{Theorems 15--18: cross-cuts}

Moore gives no proofs of Theorems~15--18, only pointers to F.A.\ Theorems~24--27
(Theorems~\ref{thm:fa24}--\ref{thm:fa27}). Since those were proved from (i)--(iii) only, and
Lemma~\ref{lem:restricted-moore} supplies (i)--(iii) for $G$, each pointer is made exact: the F.A.\ proof, read in
$G$, with the axioms it uses replaced by Theorem~\ref{thm:m25-11}.

Fix arcs of elements $pxq$, $pyq$ with $pxq\cap pyq=\set{p,q}$, $J=pxq\cup pyq$, and a third arc $pzq$; put
$J_1=pxq\cup pzq$ and $J_2=pyq\cup pzq$.

\begin{theorem}[Moore, Theorems~15 and~16]\label{thm:m25-15-16}
Suppose $\oarc{pzq}\subseteq I(J)$. Then
\begin{enumerate}[label=(\arabic*)]
\item $I(J_1)\subseteq I(J)$;
\item $\oarc{pyq}\subseteq E(J_1)$;
\item $I(J_1)\cap I(J_2)=\emptyset$;
\item $I(J)=I(J_1)\cup I(J_2)\cup\oarc{pzq}$.
\end{enumerate}
\end{theorem}

\formal{\lean{HPC.Moore.m25_theorem15}, \lean{HPC.Moore.m25_theorem16}.}

\begin{theorem}[Moore, Theorem~17]\label{thm:m25-17}
If $pxq$, $pyq$, $pzq$ are three arcs of elements, no two with a common element other than $p$, $q$, then the
interiors of the three curves they form in pairs are not mutually disjoint.
\end{theorem}

\formal{\lean{HPC.Moore.m25_theorem17}.}

\begin{theorem}[Moore, Theorem~18]\label{thm:m25-18}
Suppose $\oarc{pzq}\subseteq E(J)$. Then (1) either $\oarc{pyq}\subseteq E(J_1)$ or $\oarc{pxq}\subseteq E(J_2)$; and
(2) if $\oarc{pyq}\subseteq E(J_1)$ then $\oarc{pxq}\subseteq I(J_2)$ and $I(J_2)=I(J)\cup I(J_1)\cup\oarc{pxq}$.
\end{theorem}

\formal{\lean{HPC.Moore.m25_theorem18}.}

Moore's statement of Theorem~18 uses the letters $x$, $y$ for elements of the arcs $pxq$, $pyq$ without
introducing them, and says ``without'' for ``in the exterior of''. We read the statements for all elements of the
open arcs. This reading loses nothing: $\oarc{pyq}$ is connected and misses $J_1$, so it lies wholly within or wholly
without $J_1$, and likewise for $\oarc{pxq}$ and $J_2$.

\section{Theorem 19 and the restricted topology}

\begin{theorem}[Moore, Theorem~19]\label{thm:m25-19}
If $R$ is a Moore region and $O\in R$, there is a simple closed curve of elements $J\subseteq R$ with
$O\in I(J)\subseteq R$.
\end{theorem}

\begin{proof}
Moore: ``with the use of Theorems 11, 12, 13, 17 and 18 this theorem may be proved by a method closely analogous to
that employed to prove Theorem 36 in F.A.'' We follow the proof of F.A.\ Theorem~36
(Theorem~\ref{thm:fa35-36}(b)), with its correction, in $G$, using Moore regions as the neighbourhood basis
(Theorem~\ref{thm:m25-4}).
\begin{enumerate}[label=(\arabic*)]
\item Choose an element $X_0\in R$, $X_0\ne O$, and a Moore region $Q\ni X_0$ in $R$ with $O\notin\cl Q$
(Theorem~\ref{thm:m25-3}). By Theorem~\ref{thm:m25-14}, $Q$ contains a closed curve $J$; its interior $I$
lies in $Q$ by Theorem~\ref{thm:m25-12} (applied to $D_1=Q$, $D_2=I$), so $O\in E(J)$.
\item An arc from $O$ to an element of $J$ inside $R$, cut at its first element on $J$, gives an arc $OA$
meeting $J$ only in $A$.
\item Since $R$ is bounded and $G$ is not, $\fr R\ne\emptyset$. By Theorem~\ref{thm:m25-13} (for $R\setminus OA$,
which is connected), there is an arc from an element of $J$ to a boundary element $C$ of $R$, missing $OA$ and,
except for $C$, in $R$; cut at its last element on $J\cup I$ (and its first element on $\fr R$) it gives an arc
$BC$ meeting $J\cup I$ only in $B$.
\item The elements $A,B$ split $J$ into arcs $AF_1B$ and $AH_1B$ (Theorem~\ref{thm:m25-5-10}(e)). By
Theorem~\ref{thm:m25-13} for $R$ minus the arc $OA\cup AF_1B\cup BC$ (an arc with one end, $C$, on $\fr R$),
there is an arc $FZH$ from $AF_1B$ to $AH_1B$, missing $OA\cup AF_1B\cup BC$ except for $F$, lying except for its ends in
$E(J)$ and in $R$.
\item $B$ is without the curve $ZHAFZ$: the interior of that curve lies in $R$ (Theorem~\ref{thm:m25-12}), while the arc
$BC$ reaches $C\notin R$ without meeting the curve. By Theorem~\ref{thm:m25-18}, $A$ is within $ZHBFZ$, and so is $O$,
since $OA$ does not meet that curve. Its interior lies in $R$ by Theorem~\ref{thm:m25-12}.
\end{enumerate}
\end{proof}

\formal{\lean{HPC.Moore.m25_theorem19}.}

Of Moore's citations, Theorems~11, 12, 13 and~18 are used; Theorem~17 is not (nor is F.A.\ Theorem~26 in the
proof of F.A.\ Theorem~36), while Theorems~10 and~14, which Moore does not cite, are.

\begin{theorem}[Moore, Theorem~20]\label{thm:m25-20}
An element $p$ is a limit element of a set $H$ of elements if and only if every region in the restricted sense
containing $p$ contains an element of $H$ other than $p$. More precisely: the restricted regions form a basis of the
topology of $G$, so the topology they generate is the quotient topology.
\end{theorem}

\begin{proof}
Restricted regions are open (Lemma~\ref{lem:restricted-moore}). Given an open $V\ni p$, take a Moore region
$R\ni p$ in $V$ (Theorem~\ref{thm:m25-4}) and then a curve $J$ with $p\in I(J)\subseteq R$ (Theorem~\ref{thm:m25-19}).
So the restricted regions form a basis, and limit elements (Theorem~\ref{thm:m25-4}) are the accumulation
points for this basis.
\end{proof}

\formal{\lean{HPC.Moore.isTopologicalBasis_restricted}, \lean{HPC.Moore.m25_theorem20},
\lean{HPC.Moore.spaceRestricted_topology}.}

\section{Theorems 21 and 22}

\begin{theorem}[Moore, Theorem~21]\label{thm:m25-21}
The space of elements with the restricted regions satisfies Moore's Axioms~1--8
(Definition~\ref{def:moore-axioms}).
\end{theorem}

\begin{proof}
By Theorem~\ref{thm:m25-20} and the bridge (Proposition~\ref{prop:bridge}), Moore's notions for this space (limit
points, closure, boundary, connectedness, arcs and closed curves) are the topological notions of $G$.

\emph{Axiom 1.} $G$ is a separable metric space (Theorem~\ref{thm:quot-metrizable}). Take a countable dense set
$\set{d_m}$ and, for each $m$ and $j$, a restricted region about $d_m$ inside the ball of radius $2^{-j}$ about
$d_m$ (Theorems~\ref{thm:m25-3} and~\ref{thm:m25-19}); enumerate these, each repeated infinitely often in a
suitable order. If $P\ne\bar P$ lie in a region $R$, the fundamental regions about $P$ of small radius have closures
in $R\setminus\set{\bar P}$.

\emph{Axioms 2, 4.} A restricted region is connected, and its closure compact (Lemma~\ref{lem:restricted-moore}).

\emph{Axiom 3.} The exterior of $I(J)$ is $G\setminus(I(J)\cup J)=E(J)$, connected (Theorem~\ref{thm:m25-11}).

\emph{Axiom 5.} Moore's limit elements are the same for both systems of regions (Theorems~\ref{thm:m25-4}
and~\ref{thm:m25-20}), so this is the corresponding statement for the Moore regions: the elements containing the
points $(n,0)$, $n$ large, form an infinite set without limit element.

\emph{Axiom 8.} A simple closed curve (topological, Proposition~\ref{prop:elt-top}(b)) is a curve of elements $J$, and
it is the boundary of the restricted region $I(J)$ (Theorem~\ref{thm:m25-11}).

\emph{Axioms 6, 7.} Moore derives these from F.A.\ and the fact that every restricted region is the interior of a
closed curve. We prove them directly. Let $R=I(J)$, $\bar R$ a restricted region and $P\in\bar R$ an element of $J$.
\begin{itemize}
\item \emph{Cross-cuts near $P$.} Let $V_0\subseteq\bar R$ be a Moore region about $P$ missing some element of $J$, and
$D=I(J)$ (or $D=E(J)$). There are elements $a,b$ of $J$ splitting $J$ into arcs $X$, $Y$ with $P\in\oarc X$ and
$X\subseteq V_0$, and an arc $c$ from $a$ to $b$ in $V_0$ with $\oarc c\subseteq D$. Indeed: the two arcs of $J$ from $P$
(Theorem~\ref{thm:m25-5-10}(e)), cut near $P$ (Theorem~\ref{thm:m25-5-10}), give $a$, $b$ on either side of $P$, which
can be taken accessible from $D$ inside a small region (an arc from an element of $D$ near $J$, cut at its first
element on $J$). The two access arcs are joined in $D\cap V_0$ by Lemma~\ref{lem:quot-jan}, applied to $J$ and to
$\fr V_0$ together with the rest of $J$; these meet in a connected set.
\item \emph{Axiom 6.} Take $c$ interior to $J$ and $K=I(X\cup c)$. By Theorem~\ref{thm:m25-15-16}, $K\subseteq R$, and by
Theorem~\ref{thm:m25-12}, $K\subseteq V_0\subseteq\bar R$. By Theorem~\ref{thm:m25-19} there is a restricted region $\bar K$
about $P$ inside $V_0\setminus Y$; then $J\cap\bar K\subseteq X$, which is part of the boundary of $K$.
\item \emph{Axiom 7.} Take $c$ exterior to $J$ and $L=I(X\cup c)$. By Theorem~\ref{thm:m25-18}, either $\oarc Y\subseteq
E(X\cup c)$, and then $L$ misses $J$ and lies in $E(J)$ (as $\oarc c$ does); or $\oarc Y\subseteq L\subseteq V_0$, which is
excluded since $V_0$ misses an element of $J$.
\end{itemize}
Along the way, Moore's weaker Axioms~$6'$ and~$7'$ (his route to 6 and 7 through F.A.) also follow, from
Theorems~\ref{thm:m25-11} and~\ref{thm:m25-13}.
\end{proof}

\formal{\lean{HPC.Moore.m25r_axiom1}, \lean{HPC.Moore.m25r_axiom2}, \lean{HPC.Moore.m25r_axiom3},
\lean{HPC.Moore.m25r_axiom4}, \lean{HPC.Moore.m25r_axiom5}, \lean{HPC.Moore.m25r_axiom6},
\lean{HPC.Moore.m25r_axiom7}, \lean{HPC.Moore.m25r_axiom8}, \lean{HPC.Moore.m25_axioms6'7'},
\lean{HPC.Moore.m25_theorem21}.}

\begin{theorem}[Moore, Theorem~22]\label{thm:m25-22}
If $G$ is a Moore decomposition of the plane, the space of elements $G$ is homeomorphic to $\R^2$. In Moore's form:
there is a one-to-one correspondence between the elements of $G$ and the points of the plane such that a point $P$ is
a limit point of a point set $M$ if and only if the corresponding element is a limit element of the set of
elements corresponding to $M$.
\end{theorem}

\begin{proof}
By Theorem~\ref{thm:m25-21} and Moore's characterization of the plane (Theorem~\ref{thm:moore-plane}), $G$ with the
topology generated by the restricted regions is homeomorphic to $\R^2$; by Theorem~\ref{thm:m25-20} this topology is the
quotient topology. Homeomorphisms preserve accumulation points, and limit elements are the accumulation points of
the quotient topology (Theorem~\ref{thm:m25-4}).
\end{proof}

\formal{\lean{HPC.Moore.m25_theorem22}, \lean{HPC.Moore.m25_theorem22_moore}.}

\section{Monotone decompositions of a plane continuum}\label{sec:monotone-quotient}

The main use of Theorem~\ref{thm:m25-22} in this series is to recognise a monotone image of a plane continuum as
a plane set.

\begin{proposition}\label{prop:monotone}
Let $M$ be a compact space and $q\colon M\to Y$ a continuous surjection onto a Hausdorff space with connected
fibres $q^{-1}(y)$. Then
\begin{enumerate}[label=(\alph*)]
\item the preimage of a connected set is connected;
\item if $M$ is an indecomposable continuum, so is $Y$ (provided $Y$ has more than one point): a decomposition
$Y=A\cup B$ into proper subcontinua lifts to the decomposition $M=q^{-1}(A)\cup q^{-1}(B)$.
\end{enumerate}
\end{proposition}

\begin{proof}
(a) $q$ is closed (compact to Hausdorff). If $q^{-1}(C)=U\cup V$ with $U,V$ disjoint, nonempty and relatively closed,
each fibre over $C$, being connected, lies in $U$ or in $V$, so $C=q(U)\cup q(V)$ is a partition of $C$ into disjoint,
nonempty, relatively closed sets. (b) follows from (a), as $q^{-1}(A)$ and $q^{-1}(B)$ are proper, closed and connected.
\end{proof}

\formal{\lean{HPC.Hagopian.isPreconnected_preimage_of_monotone},
\lean{HPC.Hagopian.isDecomposableSpace_of_monotone_image}.}

Now let $M\subseteq\R^2$ be a continuum and $q\colon M\to Y$ as above. The \emph{fibre decomposition} of the plane
has as elements the fibres of $q$ and the single points of $\R^2\setminus M$.

\begin{lemma}\label{lem:fibre-usc}
The fibre decomposition is upper semicontinuous in Moore's sense (Definition~\ref{def:moore-decomp}(3)).
\end{lemma}

\begin{proof}
For fibres of $q$: if fibres $F_n=q^{-1}(y_n)$ came within $1/n$ of $F=q^{-1}(y)$ but not within $e$, pass to a
subsequence with $y_n\to y'$ (compactness of $Y$) and with points $u_n\in F_n$ near $F$ converging to $u\in F$,
and points $w_n\in F_n$ far from $F$ converging to $w$. By continuity $q(u)=y'=q(w)$, so $y'=y$ and $w\in F$,
a contradiction. Points off $M$ are handled by the compactness of $M$: a point $z\notin M$ has a neighbourhood missing $M$,
whose elements are single points.
\end{proof}

\formal{\lean{HPC.Hagopian.fibreSetoid}, \lean{HPC.Hagopian.fibre_usc}.}

\begin{corollary}\label{cor:monotone-plane}
Let $M\subseteq\R^2$ be a continuum and $q\colon M\to Y$ a continuous surjection onto a compact Hausdorff space with
connected fibres, none of which separates the plane. Then the fibre decomposition is a Moore decomposition, the
space of elements is a plane, and $Y$ is homeomorphic to a compact subset of $\R^2$ (the image of $\pi(M)$).
In the application, $M$ is a homogeneous indecomposable plane continuum and every fibre is a proper subcontinuum of
$M$; such subcontinua do not separate the plane.
\end{corollary}

\begin{proof}
The elements are compact and connected and do not separate the plane by hypothesis (points never do), and the
decomposition is upper semicontinuous by Lemma~\ref{lem:fibre-usc}. By Theorem~\ref{thm:m25-22} there is a
homeomorphism $h\colon G\to\R^2$. The map $Y\to G$, $y\mapsto q^{-1}(y)$, is a continuous injection of a compact space
into a Hausdorff space (continuity: the fibre map is the quotient map restricted to $M$, factored through $q$, which is a
quotient map as a closed surjection), so it is an embedding; composing with $h$ embeds $Y$ onto the compact set
$h(\pi(M))$.
\end{proof}

\formal{\lean{HPC.Hagopian.fibreSetoid_mooreDecomp}, \lean{HPC.Hagopian.exists_plane_quotient_of_monotone_aux},
\lean{HPC.Hagopian.exists_plane_quotient_of_monotone}.}

The formal corollary is stated for the case of the application (homogeneous indecomposable $M$, fibres proper);
the non-separation of proper subcontinua is the input that lets Theorem~\ref{thm:m25-22} apply.

\section{Further plane tools}

Once the Schoenflies theorem is available in the plane itself, several classical facts used in later volumes
follow. We record them, with the plane-specific forms that are proved.

\begin{theorem}[Cross-cuts in the plane]\label{thm:plane-split}
Let $V\subseteq\R^2$ be a bounded connected open set whose frontier is the simple closed curve $C_1\cup C_2$, with
$C_1,C_2$ arcs from $a$ to $b$ meeting only in $a,b$. If $\gamma$ is an arc from $a$ to $b$ with $\oarc\gamma\subseteq V$, then
$V\setminus\gamma=R_1\cup R_2$, where $R_1,R_2$ are disjoint bounded connected open sets with frontiers
$\gamma\cup C_1$ and $\gamma\cup C_2$.
\end{theorem}

\begin{proof}
This is Theorems~\ref{thm:fa24} and~\ref{thm:fa25} for the restricted regions of the plane: the plane with the
interiors of its simple closed curves satisfies (i)--(iii) (Jordan curve theorem, Theorem~\ref{thm:jct}), and
$V=I(C_1\cup C_2)$.
\end{proof}

\formal{\lean{HPC.Hagopian.plane_split}.}

\begin{theorem}[Schoenflies theorem in the plane, one side prescribed]\label{thm:plane-schoenflies}
Let $J\subseteq\R^2$ be a simple closed curve, $D$ its bounded complementary domain, and $\varphi\colon[0,1]\to J$ an
embedding. There is an embedding $f\colon[0,1]^2\to\R^2$ with image $J\cup D$, $f(t,0)=\varphi(t)$ for all $t$, and $f$
maps the open square onto $D$.
\end{theorem}

\begin{proof}
Apply Theorem~\ref{thm:m25-21} to the decomposition of the plane into points (Section~\ref{sec:points}): the plane,
with the interiors of its simple closed curves as regions, satisfies Moore's Axioms~1--8. Cut the arc of $J$
complementary to $\varphi([0,1])$ twice to obtain a quadrilateral whose bottom side is $\varphi([0,1])$, apply
Theorem~\ref{thm:schoenflies} to it, and reparametrize the bottom side by $\varphi$.
\end{proof}

\formal{\lean{HPC.Hagopian.plane_schoenflies_aux}.}

\begin{theorem}[Two-sided collars]\label{thm:collar}
Every arc $\gamma$ in the plane has a two-sided collar: an embedding $F$ of $[0,1]\times[-1,1]$ onto a neighbourhood of
$\oarc\gamma$ (relative to $\gamma$'s complement near the ends) with $F(t,0)=\gamma(t)$.
\end{theorem}

\begin{proof}
Join the ends of $\gamma$ by an arc $\beta$ meeting $\gamma$ only in them (both ends are accessible from the connected
complement of $\gamma$). An arc $\beta_2$ from end to end through the bounded side $I$ of $\gamma\cup\beta$ cuts $I$ into the
bounded sides $R_1$ of $\gamma\cup\beta_2$ and $R_2$ of $\beta\cup\beta_2$ (Theorem~\ref{thm:plane-split}). One half-collar is
taken in $R_1\cup\gamma\cup\beta_2$ (Theorem~\ref{thm:plane-schoenflies}); the other in the unbounded side of $\gamma\cup\beta$,
by the Schoenflies theorem after an inversion about a point of $I$, restricted to a strip along $\gamma$.
\end{proof}

\formal{\lean{HPC.Hagopian.IsCollar}, \lean{HPC.Hagopian.exists_collar}.}

\section*{Exercises}
\addcontentsline{toc}{section}{Exercises}

\begin{exercise}
Show that in Theorem~\ref{thm:m25-13} the hypothesis that at most one end of the arc lies outside $R$ cannot be
dropped: give a Moore region $R$ and an arc with both ends on $\fr R$ and its other elements in $R$ such that $R\setminus k$
is not connected.
\end{exercise}

\begin{exercise}
Prove Theorem~\ref{thm:m25-17} for the decomposition of the plane into points, directly from the Jordan curve theorem.
\end{exercise}

\begin{exercise}
Let $M$ be the closed unit disc and $q\colon M\to[0,1]$, $q(z)=\abs z$. Describe the fibre decomposition. Which hypothesis of
Corollary~\ref{cor:monotone-plane} fails, and what is the space of elements?
\end{exercise}

\begin{exercise}
Let $G$ be the decomposition of $\R^2$ whose elements are the vertical segments $\set{x}\times[-1,1]$, $\abs x\le1$, and the
single points elsewhere. Verify the hypotheses of Theorem~\ref{thm:m25-22} and write down an explicit homeomorphism
of $G$ onto $\R^2$.
\end{exercise}


\appendix
%======================================================================
\chapter{Correspondence with the formal proof}\label{app:formal}
%======================================================================

The table lists the definitions and results of this volume that carry a ``Formal counterpart''
line, together with the corresponding declarations of the Lean development. All declaration
names start with \lean{HPC.}. The formal statements are sometimes slightly more general, or are
split differently into lemmas; the table points to the place where the content of each result is
proved. Rows are in the order of the text.

The main theorems of the volume are, in the Lean development (prefix \lean{HPC.} omitted):
\begin{center}
\begin{tabular}{@{}p{0.42\textwidth}p{0.54\textwidth}@{}}
\toprule
Characterization of the arc (Corollary~\ref{cor:arc-char}) & \lean{Moore.isIrreducibleConnected_iff_arc}\\
Characterization of the circle (Theorem~\ref{thm:circle-char}) & \lean{Moore.circle_char}\\
Quotient of a Moore decomposition is metrizable (Theorem~\ref{thm:quot-metrizable}) & \lean{Moore.quotient_metrizable_plane}\\
Mullikin's theorem, corrected (Theorem~\ref{thm:mullikin}) & \lean{Moore.mullikin_theorem5}\\
Moore 1925, Theorem~11 (Theorem~\ref{thm:m25-11}) & \lean{Moore.m25_theorem11}\\
Jordan curve theorem (Theorem~\ref{thm:jct}) & \lean{Jones.jordan_curve_theorem}\\
Moore 1919, Theorem~D (Theorem~\ref{thm:theoremD}) & \lean{Moore.m19_theoremD}\\
A compact grid is a square (Theorem~\ref{thm:grid-square}) & \lean{Sq.Grid.exists_square}\\
Schoenflies theorem in a Moore space (Theorem~\ref{thm:schoenflies}) & \lean{Moore.m19_schoenflies_boundary_proof}\\
Moore's characterization of the plane (Theorem~\ref{thm:moore-plane}) & \lean{Moore.m19_plane}\\
Moore 1925, Theorem~21 (Theorem~\ref{thm:m25-21}) & \lean{Moore.m25_theorem21}\\
Moore 1925, Theorem~22 (Theorem~\ref{thm:m25-22}) & \lean{Moore.m25_theorem22}, \lean{Moore.m25_theorem22_moore}\\
Monotone images (Corollary~\ref{cor:monotone-plane}) & \lean{Hagopian.exists_plane_quotient_of_monotone}\\
Schoenflies theorem in the plane (Theorem~\ref{thm:plane-schoenflies}) & \lean{Hagopian.plane_schoenflies_aux}\\
Two-sided collars (Theorem~\ref{thm:collar}) & \lean{Hagopian.exists_collar}\\
\bottomrule
\end{tabular}
\end{center}
All of them are proved in the Lean development without \lean{sorry} and without axioms beyond the
standard ones of Lean and Mathlib. The printed forms that are false (Section~\ref{sec:intro-diff}) are
refuted by the declarations whose names begin with \lean{not_}.

\begin{longtable}{@{}p{0.40\textwidth}p{0.56\textwidth}@{}}
\toprule
In these notes & In the Lean development\\
\midrule
\endhead
\bottomrule
\endfoot
\multicolumn{2}{@{}l}{\emph{Chapter~\ref{ch:arcs}}}\\
Definition~\ref{def:arc} & \lean{HPC.Moore.Gen.TArc}, \lean{HPC.Moore.Gen.TSCC}.\\
Sub-arcs (Lemma~\ref{lem:subarc}) & \lean{HPC.Moore.Gen.TArc.image_Icc}, \lean{HPC.Moore.Gen.TArc.exists_sub}, \lean{HPC.Moore.Gen.TArc.nonempty_diff}, \lean{HPC.Moore.Gen.TArc.isPreconnected_diff}, \lean{HPC.Moore.Gen.TArc.subset_closure_diff}.\\
Lemma~\ref{lem:conn-frontier} & \lean{HPC.Moore.Gen.conn_sub_of_frontier}, \lean{HPC.Moore.Gen.conn_sub_ext_of_frontier}.\\
First points, gluing, joining (Lemma~\ref{lem:arc-ops}) & \lean{HPC.Moore.Gen.TArc.first_point}, \lean{HPC.Moore.Gen.TArc.union}, \lean{HPC.Moore.Gen.TArc.join}.\\
Lemma~\ref{lem:insert-conn} & \lean{HPC.Moore.ac_insert_isPreconnected}.\\
Cuts (Lemma~\ref{lem:cut}) & \lean{HPC.Moore.ac_compl_isPreconnected}, \lean{HPC.Moore.ac_cut}.\\
Proposition~\ref{prop:arc-order} & \lean{HPC.Moore.acL}, \lean{HPC.Moore.acL_asymm}, \lean{HPC.Moore.acL_trans}, \lean{HPC.Moore.acL_total}, \lean{HPC.Moore.acL_ray}, \lean{HPC.Moore.ac_exists_between}.\\
Irreducible continua are arcs (Theorem~\ref{thm:irr-arc}) & \lean{HPC.Moore.ac_exists_strictMono}, \lean{HPC.Moore.ac_homeomorph}.\\
Arc characterization (Corollary~\ref{cor:arc-char}) & \lean{HPC.Moore.isIrreducibleConnected_iff_arc}.\\
Definition~\ref{def:cutarc} & \lean{HPC.Moore.IsCutArc}.\\
Proposition~\ref{prop:cutarc} & (a) \lean{HPC.Moore.IsCutArc.irreducible}; (b) \lean{HPC.Moore.IsCutArc.exists_embedding}, \lean{HPC.Moore.isCutArc_range}, \lean{HPC.Moore.isCutArc_image_Icc}; (c) \lean{HPC.Moore.eq_image_Icc_of_subset_range}, \lean{HPC.Moore.ends_of_image_Icc}; (d) \lean{HPC.Moore.exists_nonCut_diff}.\\
Moore 1920, Theorem~4 (Theorem~\ref{thm:circle-char}) & \lean{HPC.Moore.circle_noCut}, \lean{HPC.Moore.circle_half}, \lean{HPC.Moore.circle_char}.\\
Proposition~\ref{prop:scc-cut} & \lean{HPC.Moore.Gen.tscc_cut}; \lean{HPC.Jones.glue_circle}, \lean{HPC.Jones.tscc_homeomorph_circle}.\\
Limit arcs (Theorem~\ref{thm:limit-arc}) & \lean{HPC.Moore.Gen.tarc_of_seq}; with trimming of arcs that meet more than consecutively, \lean{HPC.Moore.Gen.exists_tarc_of_seq}.\\
Replacement of a null family of sub-arcs (Theorem~\ref{thm:arc-replace}) & \lean{HPC.Moore.Gen.arc_replace}.\\[4pt]
\multicolumn{2}{@{}l}{\emph{Chapter~\ref{ch:usc}}}\\
Definition~\ref{def:moore-decomp} & \lean{HPC.Moore.MooreDecomp}, \lean{HPC.Moore.elt}, \lean{HPC.Moore.ptsOf}.\\
Definition~\ref{def:tube} & \lean{HPC.Moore.tube}.\\
Proposition~\ref{prop:closed-map} & \lean{HPC.Moore.isClosedMap_mk_plane}, \lean{HPC.Moore.isOpen_tube}, \lean{HPC.Moore.mem_nhds_iff_tube}, \lean{HPC.Moore.usc_of_isClosedMap}.\\
Theorem~\ref{thm:quot-metrizable} & \lean{HPC.Moore.t1Space_quot}, \lean{HPC.Moore.regularSpace_quot}, \lean{HPC.Moore.secondCountable_quot}, \lean{HPC.Moore.quotient_metrizable_plane}.\\
Proposition~\ref{prop:pts-conn} & \lean{HPC.Moore.isPreconnected_ptsOf}, \lean{HPC.Moore.ptsOf_connectedComponentIn}, \lean{HPC.Moore.isOpen_connectedComponentIn_quot}, \lean{HPC.Moore.locallyConnectedSpace_quot}, \lean{HPC.Moore.frontier_eq_image}.\\
Corollary~\ref{cor:quot-tools} & \lean{HPC.Moore.quot_brouwer}, \lean{HPC.Moore.quot_janiszewski}, \lean{HPC.Moore.m25_unbounded_unique}, \lean{HPC.Moore.isConnected_compl_singleton_quot}.\\
Proposition~\ref{prop:moore-notions} & \lean{HPC.Moore.isLimitElt_iff}, \lean{HPC.Moore.m25_domain_iff}, \lean{HPC.Moore.m25_boundary_iff}.\\
Moore, Theorem~1 (Theorem~\ref{thm:m25-1}) & \lean{HPC.Moore.m25_theorem1}.\\
Definition~\ref{def:moore-region} & \lean{HPC.Moore.mooreRegion}, \lean{HPC.Moore.spaceOfElts}.\\
Moore, Theorem~3 (Theorem~\ref{thm:m25-3}) & \lean{HPC.Moore.exists_region_in_tube}, \lean{HPC.Moore.m25_theorem3}.\\
Moore, Theorem~4 (Theorem~\ref{thm:m25-4}) & \lean{HPC.Moore.isTopologicalBasis_mooreRegion}, \lean{HPC.Moore.m25_theorem4}, \lean{HPC.Moore.spaceOfElts_topology}, \lean{HPC.Moore.cl_eq_closure_quot}.\\
Moore, Theorem~2 (Theorem~\ref{thm:m25-2}) & \lean{HPC.Moore.m25_theorem2}.\\
Proposition~\ref{prop:p420} & \lean{HPC.Moore.m25_axioms1245}, \lean{HPC.Moore.m25_FA_1_15}.\\[4pt]
\multicolumn{2}{@{}l}{\emph{Chapter~\ref{ch:elements}}}\\
Definition~\ref{def:arc-elt} & \lean{HPC.Moore.IsArcElt}, \lean{HPC.Moore.IsSCCElt}.\\
Proposition~\ref{prop:arc-elt-top} & \lean{HPC.Moore.isCompact_of_isClosed_bounded}, \lean{HPC.Moore.isArcElt_iff_tarc}, \lean{HPC.Moore.isSCCElt_iff_tscc}.\\
Moore, Theorems 5--8 and 10 (Theorem~\ref{thm:m25-5-10}) & \lean{HPC.Moore.m25_theorem5}, \lean{HPC.Moore.m25_theorem6}, \lean{HPC.Moore.m25_theorem7}, \lean{HPC.Moore.m25_theorem8}, \lean{HPC.Moore.m25_theorem10}.\\
Moore, Theorem~9 (Theorem~\ref{thm:m25-9}) & \lean{HPC.Moore.m25_theorem9}.\\
Corollary~\ref{cor:pts-nonsep} & \lean{HPC.Moore.not_separatesPlane_ptsOf}.\\
Moore, Theorem~12 (Theorem~\ref{thm:m25-12}) & \lean{HPC.Moore.m25_theorem12}.\\
Text of Section~\ref{sec:points} & \lean{HPC.Moore.mooreDecomp_ptSetoid}, \lean{HPC.Moore.ptHomeo}.\\[4pt]
\multicolumn{2}{@{}l}{\emph{Chapter~\ref{ch:jordan}}}\\
At least two components (Proposition~\ref{prop:mullikin-two}) & \lean{HPC.Moore.not_isConnected_compl_union}.\\
At most two components; Eilenberg (Proposition~\ref{prop:mullikin-three}) & \lean{HPC.Moore.exists_jump}, \lean{HPC.Moore.no_log_pow_of_bounded_component}, \lean{HPC.Moore.not_two_bounded_components}, \lean{HPC.Moore.compl_union_two_components}.\\
Mullikin's theorem (Theorem~\ref{thm:mullikin}) & \lean{HPC.Moore.mullikin_theorem5}.\\
Lemma~\ref{lem:frontier-scc} & \lean{HPC.Moore.frontier_eq_of_scc}.\\
Moore, Theorem~11 (Theorem~\ref{thm:m25-11}) & \lean{HPC.Moore.m25_theorem11}, \lean{HPC.Moore.sccInterior_spec}.\\
Jordan curve theorem (Theorem~\ref{thm:jct}) & \lean{HPC.Jones.jordan_curve_theorem}.\\[4pt]
\multicolumn{2}{@{}l}{\emph{Chapter~\ref{ch:axioms}}}\\
Moore's derived notions (Definition~\ref{def:moore-notions}) & \lean{HPC.Moore.Space}, \lean{HPC.Moore.Space.IsLimitPt}, \lean{HPC.Moore.Space.bdry}, \lean{HPC.Moore.Space.cl}, \lean{HPC.Moore.Space.ext}, \lean{HPC.Moore.Space.IsConn}, \lean{HPC.Moore.Space.IsBoundedM}, \lean{HPC.Moore.Space.IsDomain}, \lean{HPC.Moore.Space.IsArc}, \lean{HPC.Moore.Space.IsSCC}, \lean{HPC.Moore.Space.IsSimpleChain}.\\
Moore's axioms, F.A.\ pp.~132--133 (Definition~\ref{def:moore-axioms}) & \lean{HPC.Moore.Space.Axiom1} \dots\ \lean{HPC.Moore.Space.Axiom8}, \leansub{HPC.Moore.Space.Axiom5}{1}, \leansub{HPC.Moore.axiom5}{1}\lean{_of_axiom5}.\\
The bridge (Proposition~\ref{prop:bridge}) & \lean{HPC.Moore.Space.topology}, \lean{HPC.Moore.fa_theorem6}, \lean{HPC.Moore.isTopologicalBasis_region}, \lean{HPC.Moore.isLimitPt_iff}, \lean{HPC.Moore.cl_eq_closure}, \lean{HPC.Moore.isClosedM_iff}, \lean{HPC.Moore.isConn_iff}, \lean{HPC.Moore.isDomain_iff}, \lean{HPC.Moore.heineBorel_iff}, \lean{HPC.Moore.t1Space}, \lean{HPC.Moore.regularSpace}, \lean{HPC.Moore.fa_theorem4}, \lean{HPC.Moore.secondCountable}, \lean{HPC.Moore.metrizable}, \lean{HPC.Moore.isCompact_closure_region}, \lean{HPC.Moore.isBoundedM_of_isCompact}.\\
F.A.\ Theorems~1--3, 5, 7--9, 11--14, 17--19, 22 (Theorem~\ref{thm:fa-elementary}) & \lean{HPC.Moore.fa_theorem1}, \lean{HPC.Moore.fa_theorem2}, \lean{HPC.Moore.fa_theorem3}, \lean{HPC.Moore.fa_theorem5}, \lean{HPC.Moore.fa_theorem7}, \lean{HPC.Moore.fa_theorem8}, \lean{HPC.Moore.fa_theorem9}, \lean{HPC.Moore.fa_theorem11}, \lean{HPC.Moore.fa_theorem12}, \lean{HPC.Moore.fa_theorem13}, \lean{HPC.Moore.fa_theorem14}, \lean{HPC.Moore.fa_theorem17}, \lean{HPC.Moore.fa_theorem18}, \lean{HPC.Moore.fa_theorem19}, \lean{HPC.Moore.isConn_univ}, \lean{HPC.Moore.fa_theorem22}.\\
F.A.\ Theorem~10 (Theorem~\ref{thm:fa10}) & \lean{HPC.Moore.fa_theorem10}, \lean{HPC.Moore.exists_simpleChain_sublist}.\\
F.A.\ Theorem~15 (Theorem~\ref{thm:fa15}) & \lean{HPC.Moore.exists_chains}, \lean{HPC.Moore.exists_irreducible}, \lean{HPC.Moore.fa_theorem15}.\\
F.A.\ Theorem~16 (Corollary~\ref{cor:fa16}) & \lean{HPC.Moore.fa_theorem16}.\\
Proposition~\ref{prop:moore-arcs} & \lean{HPC.Moore.isArc_iff}, \lean{HPC.Moore.isSCC_iff_tscc}.\\[4pt]
\multicolumn{2}{@{}l}{\emph{Chapter~\ref{ch:curves}}}\\
F.A.\ Theorems~20 and~21 (Theorem~\ref{thm:fa20-21}) & \lean{HPC.Moore.fa_theorem20}, \lean{HPC.Moore.fa_theorem21}.\\
The interior of a closed curve (Corollary~\ref{cor:interior}) & \lean{HPC.Moore.existsUnique_interior}, \lean{HPC.Moore.jordan}, \lean{HPC.Moore.Space.IsInteriorOf}.\\
Text of the section ``The interior of a closed curve'' & \lean{HPC.Moore.Gen.ArcConn}, \lean{HPC.Moore.Gen.Th21Prop}, \lean{HPC.Moore.Gen.Ax8Prop}, \lean{HPC.Moore.Gen.TInt}.\\
Text of the section ``The interior of a closed curve'' & \lean{HPC.Moore.Gen.exists_access}, \lean{HPC.Moore.Gen.exists_cross}.\\
F.A.\ Theorem~23 (Theorem~\ref{thm:fa23}) & \lean{HPC.Moore.fa_theorem23}.\\
F.A.\ Theorem~24 (Theorem~\ref{thm:fa24}) & \lean{HPC.Moore.fa_theorem24}, \lean{HPC.Moore.Gen.th24_sub}, \lean{HPC.Moore.Gen.th24_ext}, \lean{HPC.Moore.Gen.th24_disj}.\\
F.A.\ Theorem~25 (Theorem~\ref{thm:fa25}) & \lean{HPC.Moore.fa_theorem25}, \lean{HPC.Moore.Gen.th25}.\\
F.A.\ Theorem~26 (Theorem~\ref{thm:fa26}) & \lean{HPC.Moore.fa_theorem26}, \lean{HPC.Moore.Gen.th26}.\\
F.A.\ Theorem~27 (Theorem~\ref{thm:fa27}) & \lean{HPC.Moore.fa_theorem27}, \lean{HPC.Moore.Gen.th27}.\\
Splitting an interior (Corollary~\ref{cor:split-interior}) & \lean{HPC.Moore.split_interior}, \lean{HPC.Moore.split_interior_of}.\\
F.A.\ Theorem~28 (Theorem~\ref{thm:fa28}) & \lean{HPC.Moore.fa_theorem28}.\\
Text of the section ``Arcs near a curve: Theorems 28--30'' & \lean{HPC.Moore.fa_theorem29}, \lean{HPC.Moore.fa_theorem30}.\\
F.A.\ Theorem~31; the pocket construction (Theorem~\ref{thm:fa31}) & \lean{HPC.Moore.fa_theorem31}, \lean{HPC.Moore.th31_pocket}.\\
The two sides (Corollary~\ref{cor:sides}) & \lean{HPC.Moore.sides}.\\
Propagation and accessibility; F.A.\ Theorem~32 (Theorem~\ref{thm:fa32}) & \lean{HPC.Moore.side_propagate}, \lean{HPC.Moore.accessible}, \lean{HPC.Moore.fa_theorem32}.\\
F.A.\ Theorems~33 and~34 (Theorem~\ref{thm:fa33-34}) & \lean{HPC.Moore.reach_end}, \lean{HPC.Moore.joinZ}, \lean{HPC.Moore.fa_theorem33}, \lean{HPC.Moore.fa_theorem34_region}, \lean{HPC.Moore.fa_theorem34}.\\
F.A.\ Theorems~39 and~40; accessibility of a closed curve (Theorem~\ref{thm:fa39-40}) & \lean{HPC.Moore.access}, \lean{HPC.Moore.fa_theorem39}, \lean{HPC.Moore.fa_theorem40}.\\
F.A.\ Theorems~35 and~36 (Theorem~\ref{thm:fa35-36}) & \lean{HPC.Moore.fa_theorem35}, \lean{HPC.Moore.fa_theorem36}, \lean{HPC.Moore.fa_theorem36_interior_subset}.\\
Chords form a null family; Moore 1919, Theorem~A (Lemma~\ref{lem:chords-null}) & \lean{HPC.Moore.chords_null}, \lean{HPC.Moore.m19_theoremA}.\\
Text of the section ``Merging interiors: Theorems 37--43'' & \lean{HPC.Moore.exists_maximal_chord}, \lean{HPC.Moore.fa_theorem37}, \lean{HPC.Moore.fa_theorem38}.\\
F.A.\ Theorem~41 (Theorem~\ref{thm:fa41}) & \lean{HPC.Moore.fa_theorem41_caseI_nested}, \lean{HPC.Moore.fa_theorem41_caseI}, \lean{HPC.Moore.fa_theorem41_of_subset}, \lean{HPC.Moore.fa_theorem41_caseII}, \lean{HPC.Moore.Gen.arc_replace}, \lean{HPC.Moore.fa_theorem41}.\\
F.A.\ Theorem~42 (Theorem~\ref{thm:fa42}) & \lean{HPC.Moore.fa_theorem42_of_41}, \lean{HPC.Moore.fa_theorem42}.\\
F.A.\ Theorem~43, corrected (Theorem~\ref{thm:fa43}) & \lean{HPC.Moore.dual_component}, \lean{HPC.Moore.fa_theorem43}.\\[4pt]
\multicolumn{2}{@{}l}{\emph{Chapter~\ref{ch:rulings}}}\\
Definition~\ref{def:quad} & \lean{HPC.Moore.Space.Quad}.\\
Definition~\ref{def:double-ruling} & \lean{HPC.Moore.Space.DoubleRuling}, \lean{HPC.Moore.Space.DoubleRuling.carrier}, \lean{HPC.Moore.Space.DoubleRuling.cells}.\\
Text of the section ``Quadrilaterals and double rulings'' & \lean{HPC.Moore.Space.ArcRuling}, \lean{HPC.Moore.Space.ArcRuling.toDoubleRuling}, \lean{HPC.Moore.Space.DoubleRuling.toArcRuling}, \lean{HPC.Moore.Space.ArcRuling.param_mono}.\\
Lemma~\ref{lem:tiling} & \lean{HPC.Moore.Space.Tiling}, \lean{HPC.Moore.Space.Tiling.add_arc}, \lean{HPC.Moore.Space.Tiling.add_arc_finite}, \lean{HPC.Moore.Space.Tiling.components}.\\
Subdivision lemma (Proposition~\ref{prop:cells-count}) & \lean{HPC.Moore.Space.DoubleRuling.tiling}, \lean{HPC.Moore.doubleRuling_cells_ncard}.\\
Text of the section ``Cutting, gluing, extending and refining'' & \lean{HPC.Moore.Space.Quad.splitX}, \lean{HPC.Moore.Space.Quad.rot}, \lean{HPC.Moore.Space.Quad.Cut}.\\
Cut and glue (Lemma~\ref{lem:cut-glue}) & \lean{HPC.Moore.Space.Quad.Cut.restrictA}, \lean{HPC.Moore.Space.Quad.Cut.restrictB}, \lean{HPC.Moore.Space.Quad.Cut.glue}, \lean{HPC.Moore.Space.ArcRuling.extendY}, \lean{HPC.Moore.Space.ArcRuling.extendY_finite}.\\
Text of the section ``Cutting, gluing, extending and refining'' & \lean{HPC.Moore.Space.Rulable}, \lean{HPC.Moore.Space.DoubleRulable}.\\
Refinement (Proposition~\ref{prop:refine}) & \lean{HPC.Moore.Space.ArcRuling.refine}, \lean{HPC.Moore.doubleRuling_refine}, \lean{HPC.Moore.Space.Quad.Cut.rulable}.\\
Definition~\ref{def:chord-quad} & \lean{HPC.Moore.Space.IsChord}, \lean{HPC.Moore.Space.Crosses}, \lean{HPC.Moore.Space.IsCrossCut}.\\
Text of the section ``Theorem C'' & \lean{HPC.Moore.no_ruling_of_two_branches}, \lean{HPC.Moore.tangent_configuration}, \lean{HPC.Moore.not_m19_theoremC_printed}.\\
Chords (Proposition~\ref{prop:chords-rulable}) & \lean{HPC.Moore.Space.Pocket}, \lean{HPC.Moore.chordWeight}, \lean{HPC.Moore.outer_tile}, \lean{HPC.Moore.exists_gap}, \lean{HPC.Moore.chords_rulable}.\\
A curve within (Proposition~\ref{prop:inside-rulable}) & \lean{HPC.Moore.inside_rulable}.\\
Text of the section ``Theorem C'' & \lean{HPC.Moore.extract_chords}, \lean{HPC.Moore.curve_rulable}, \lean{HPC.Moore.m19_theoremC_corrected}, \lean{HPC.Moore.m19_theoremC_crosscuts}, \lean{HPC.Moore.m19_theoremC_within}.\\
A nice curve about a compact set (Proposition~\ref{prop:nice-curve}) & \lean{HPC.Moore.lens_data}, \lean{HPC.Moore.lens_family}, \lean{HPC.Moore.exists_nice_curve}.\\
Moore 1919, Theorem~D (Theorem~\ref{thm:theoremD}) & \lean{HPC.Moore.rulable_of_cover}, \lean{HPC.Moore.m19_theoremD}.\\
Theorem~D for regions (Corollary~\ref{cor:theoremD-regions}) & \lean{HPC.Moore.m19_theoremD_regions_of_curves}, \lean{HPC.Moore.m19_theoremD_regions}.\\[4pt]
\multicolumn{2}{@{}l}{\emph{Chapter~\ref{ch:plane}}}\\
Definition~\ref{def:family-lines} & \lean{HPC.Sq.Family}, \lean{HPC.Sq.PreGrid}, \lean{HPC.Sq.box}, \lean{HPC.Sq.Grid}.\\
A compact grid is a square (Theorem~\ref{thm:grid-square}) & \lean{HPC.Sq.Family.pos}, \lean{HPC.Sq.Grid.dense_t}, \lean{HPC.Sq.Grid.mem_L_of_pos_eq}, \lean{HPC.Sq.Grid.continuousOn_pos}, \lean{HPC.Sq.Grid.eq_of_pos_eq}, \lean{HPC.Sq.Grid.image_pos}, \lean{HPC.Sq.Grid.exists_square}, \lean{HPC.Sq.Grid.exists_square_param}.\\
Boxes are components (Lemma~\ref{lem:boxes}) & \lean{HPC.Sq.PreGrid.box_nonempty}, \lean{HPC.Sq.PreGrid.box_sub_component}.\\
Nested rulings (Lemma~\ref{lem:nested-rulings}) & \lean{HPC.Moore.nested_rulings}, \lean{HPC.Moore.nested_rulings_fund}.\\
Schoenflies; Moore 1919, Theorem~E (Theorem~\ref{thm:schoenflies}) & \lean{HPC.Moore.lineFam}, \lean{HPC.Moore.lr_mono}, \lean{HPC.Moore.preGrid}, \lean{HPC.Moore.preGrid_box_sub_cell}, \lean{HPC.Moore.m19_schoenflies_proof}, \lean{HPC.Moore.m19_schoenflies_param}, \lean{HPC.Moore.m19_schoenflies_boundary_proof}, \lean{HPC.Moore.m19_schoenflies_boundary}.\\
Corollary~\ref{cor:schoenflies-curve} & \lean{HPC.Moore.exists_quad}, \lean{HPC.Moore.exists_square_of_interior}.\\
Annulus lemma (Theorem~\ref{thm:annulus}) & \lean{HPC.Moore.exists_cuts}, \lean{HPC.Moore.annulus_decomp}, \lean{HPC.Moore.exists_collar}, \lean{HPC.Moore.m19_annulus}, \lean{HPC.Moore.m19_annulus_square}.\\
Lemma~\ref{lem:curve-about-compact} & \lean{HPC.Moore.preconnectedSpace}, \lean{HPC.Moore.exists_curve_about_compact}.\\
Moore 1919, Theorem~F (Theorem~\ref{thm:theoremF}) & \lean{HPC.Moore.m19_theoremF_proof}, \lean{HPC.Moore.m19_theoremF}.\\
Pasting (Proposition~\ref{prop:paste}) & \lean{HPC.Sq.exhaust_homeomorph}, \lean{HPC.Sq.plane_of_collars}.\\
Moore's characterization of the plane (Theorem~\ref{thm:moore-plane}) & \lean{HPC.Moore.m19_plane_proof}, \lean{HPC.Moore.m19_plane}.\\
Moore 1919, Theorem~G and the final paragraph (Corollary~\ref{cor:theoremG}) & \lean{HPC.Moore.m19_theoremG}, \lean{HPC.Moore.m19_openCurve}, \lean{HPC.Moore.m19_final}.\\[4pt]
\multicolumn{2}{@{}l}{\emph{Chapter~\ref{ch:quotient}}}\\
Interiors and exteriors (chapter introduction) & \lean{HPC.Moore.sccInterior}, \lean{HPC.Moore.sccExterior}, \lean{HPC.Moore.sccInterior_spec}, \lean{HPC.Moore.closure_sccInterior}, \lean{HPC.Moore.closure_sccExterior}.\\
Proposition~\ref{prop:elt-top} & \lean{HPC.Moore.isArcElt_iff_tarc}, \lean{HPC.Moore.isSCCElt_iff_tscc}, \lean{HPC.Moore.isSCCElt_union}, \lean{HPC.Moore.t2Space_quot}, \lean{HPC.Moore.arcConn_quot}, \lean{HPC.Jones.tscc_homeomorph_circle}.\\
Definition~\ref{def:restricted} & \lean{HPC.Moore.restrictedRegion}, \lean{HPC.Moore.spaceRestricted}.\\
Lemma~\ref{lem:restricted-moore} & \lean{HPC.Moore.mooreRegion_of_restricted}, \lean{HPC.Moore.th21Prop_restricted}, \lean{HPC.Moore.ax8Prop_restricted}, \lean{HPC.Moore.isCompact_closure_restricted}, \lean{HPC.Moore.not_isCompact_univ_quot}.\\
Janiszewski's theorem for elements (Lemma~\ref{lem:quot-jan}) & \lean{HPC.Moore.quot_janiszewski}.\\
Moore, Theorem~13 (Theorem~\ref{thm:m25-13}) & \lean{HPC.Moore.not_subset_arc_of_cut}, \lean{HPC.Moore.m25_theorem13}.\\
Moore, Theorem~14 (Theorem~\ref{thm:m25-14}) & \lean{HPC.Moore.exists_tscc_of_arcs}, \lean{HPC.Moore.m25_theorem14}.\\
Moore, Theorems~15 and~16 (Theorem~\ref{thm:m25-15-16}) & \lean{HPC.Moore.m25_theorem15}, \lean{HPC.Moore.m25_theorem16}.\\
Moore, Theorem~17 (Theorem~\ref{thm:m25-17}) & \lean{HPC.Moore.m25_theorem17}.\\
Moore, Theorem~18 (Theorem~\ref{thm:m25-18}) & \lean{HPC.Moore.m25_theorem18}.\\
Moore, Theorem~19 (Theorem~\ref{thm:m25-19}) & \lean{HPC.Moore.m25_theorem19}.\\
Moore, Theorem~20 (Theorem~\ref{thm:m25-20}) & \lean{HPC.Moore.isTopologicalBasis_restricted}, \lean{HPC.Moore.m25_theorem20}, \lean{HPC.Moore.spaceRestricted_topology}.\\
Moore, Theorem~21 (Theorem~\ref{thm:m25-21}) & \lean{HPC.Moore.m25r_axiom1}, \lean{HPC.Moore.m25r_axiom2}, \lean{HPC.Moore.m25r_axiom3}, \lean{HPC.Moore.m25r_axiom4}, \lean{HPC.Moore.m25r_axiom5}, \lean{HPC.Moore.m25r_axiom6}, \lean{HPC.Moore.m25r_axiom7}, \lean{HPC.Moore.m25r_axiom8}, \lean{HPC.Moore.m25_axioms6'7'}, \lean{HPC.Moore.m25_theorem21}.\\
Moore, Theorem~22 (Theorem~\ref{thm:m25-22}) & \lean{HPC.Moore.m25_theorem22}, \lean{HPC.Moore.m25_theorem22_moore}.\\
Proposition~\ref{prop:monotone} & \lean{HPC.Hagopian.isPreconnected_preimage_of_monotone}, \lean{HPC.Hagopian.isDecomposableSpace_of_monotone_image}.\\
Lemma~\ref{lem:fibre-usc} & \lean{HPC.Hagopian.fibreSetoid}, \lean{HPC.Hagopian.fibre_usc}.\\
Corollary~\ref{cor:monotone-plane} & \lean{HPC.Hagopian.fibreSetoid_mooreDecomp}, \lean{HPC.Hagopian.exists_plane_quotient_of_monotone_aux}, \lean{HPC.Hagopian.exists_plane_quotient_of_monotone}.\\
Cross-cuts in the plane (Theorem~\ref{thm:plane-split}) & \lean{HPC.Hagopian.plane_split}.\\
Schoenflies theorem in the plane, one side prescribed (Theorem~\ref{thm:plane-schoenflies}) & \lean{HPC.Hagopian.plane_schoenflies_aux}.\\
Two-sided collars (Theorem~\ref{thm:collar}) & \lean{HPC.Hagopian.IsCollar}, \lean{HPC.Hagopian.exists_collar}.\\

\end{longtable}


\end{document}
